% This file is intended to be \input{} inside subsubsec__window6-cube-fold6-additional-results.tex.

\paragraph{并行对象链：算子交换子包络与抽象括号的审计约束}
“李括号的构造与唯一闭合”在 window-\(6\) 处被明确对象化为由离散可审计数据诱导的两条并行对象链。
其一由超立方体坐标翻转边在折叠 \(F_6\) 下的推送计数定义生成元 \(\{L_i\}\)，并在 \(\mathfrak{gl}(V)\) 内以交换子 \([A,B]=AB-BA\) 生成最小 Lie 子代数闭包，从而得到一个由有限边计数完全决定的交换子 Lie 包络 \(\mathfrak{g}_6\)。
其二以 \(X_6=X^{\mathrm{cyc}}_6\sqcup X^{\mathrm{bdry}}_6\) 的 \(18\oplus 3\) 分解为输入，在 \(V=\RR\langle X_6\rangle\) 上构造抽象括号并要求其与局域性、分块与奇偶等离散证书相容；在该链条中，“唯一闭合”指向在固定标号与可审计约束下的括号存在/唯一性，以及由几何 \(\ZZ_2\) 与 orbit--Levi 骨架诱导的紧致子结构极小化与残余代数选出问题。
下文先对第一条链给出闭包与结构识别的定理结论，再将第二条链中与猜想 \ref{conj:window6-intrinsic-bracket-uniqueness} 对应的存在/唯一性推进为证书化表述。

\begin{definition}[坐标翻转推送算子与交换子 Lie 包络]\label{def:window6-coordinate-pushforward-lie-envelope}
令 \(V:=\RR\langle X_6\rangle\simeq \RR^{X_6}\) 为以 \(X_6\) 为基的实向量空间。
对每个坐标方向 \(i\in\{1,\dots,6\}\)，令 \(L_i\in \mathrm{End}(V)\) 为由 \(Q_6\) 的第 \(i\) 位翻转边推送到 \(X_6\) 所得的计数算子：其矩阵元定义为
\[
(L_i)_{uv}:=\#\bigl\{\omega\in\Omega_6:\ F_6(\omega)=u,\ F_6(\omega\oplus e_i)=v\bigr\}
\qquad (u,v\in X_6),
\]
其中 \(e_i\) 为第 \(i\) 位标准基向量。
在 \(\mathfrak{gl}(V)\) 中以交换子
\(
[A,B]:=AB-BA
\)
作为李括号，并定义由 \(\{L_i\}\) 生成的最小 Lie 子代数
\[
\mathfrak{g}_6:=\mathrm{Lie}\langle L_1,\dots,L_6\rangle\ \subset\ \mathfrak{gl}(V),
\]
称为 window-\(6\) 的交换子 Lie 包络。
\end{definition}

\begin{lemma}[Fold\texorpdfstring{$^{\mathrm{bin}}_6$}{bin6} 的 Zeckendorf 低位投影实现]\label{lem:window6-foldbin6-zeckendorf-lowbit-projection}
\leanverified{paper\_window6\_foldbin6\_zeckendorf\_lowbit\_projection}
对任意整数 \(n\in\{0,\dots,63\}\)，取其 Zeckendorf 展开
\[
n=\sum_{k=2}^{10} d_k\,F_k,
\qquad
d_k\in\{0,1\},
\qquad
d_kd_{k+1}=0.
\]
则 window-\(6\) 的 dyadic 折叠读出可等价表述为
\[
\mathrm{Fold}^{\mathrm{bin}}_6(n)=(d_2,d_3,d_4,d_5,d_6,d_7)\in X_6.
\]
并且其纤维偏移结构由尾部三元组 \((F_8,F_9,F_{10})=(21,34,55)\) 的可行组合完全参数化（命题 \ref{prop:terminal-foldbin6-tail-cube-section}）。
\end{lemma}
\begin{proof}
该结论是命题 \ref{prop:terminal-foldbin6-tail-cube-section} 的“低位截断”重述：在 \(m=6\) 锚点上，\(n\in[0,2^6)\) 的 Zeckendorf 数字在低位 \(\{F_2,\dots,F_7\}\) 上无歧义，且纤维中的全部额外自由度恰由尾部权重 \((F_8,F_9,F_{10})\) 的可行组合给出。
因此 \(\mathrm{Fold}^{\mathrm{bin}}_6\) 可用低位 Zeckendorf 数字实现，而纤维参数化与文稿给出的尾部立方截面描述一致。证毕。
\leanverified{paper\_window6\_foldbin6\_zeckendorf\_lowbit\_projection}
\end{proof}

\begin{lemma}[window-\(6\) 纤维的差分谱]\label{lem:window6-foldbin6-fiber-difference-spectrum}
设 \(w\in X_6\)，并记其二进制折叠纤维为
\(
\mathrm{Fold}^{\mathrm{bin}}_6{}^{-1}(w)\subset\{0,\dots,63\}
\)。
则同一纤维中任意两点之差必属于集合
\[
\Delta_6:=\{0,\pm 13,\pm 21,\pm 34,\pm 55\}.
\]
\end{lemma}
\begin{proof}
命题 \ref{prop:terminal-foldbin6-tail-cube-section} 表明：任意纤维在尾部三权重 \((21,34,55)\) 的口径下可写为一个偏移集合的平移，并且该偏移集合包含于 \(\{0,21,34,55\}\)。
因此，同一纤维中两点之差为两偏移之差，落在
\(
\{0,\pm 21,\pm 34,\pm 55,\pm(34-21)\}
=\{0,\pm 21,\pm 34,\pm 55,\pm 13\}
\)。
证毕。
\leanverified{paper\_window6\_foldbin6\_fiber\_difference\_spectrum}
\end{proof}

\begin{theorem}[最小交换子闭包的存在与唯一性]\label{thm:window6-lie-envelope-existence-uniqueness}
令 \(S:=\{L_1,\dots,L_6\}\subset\mathfrak{gl}(V)\)。
则存在唯一 Lie 子代数
\[
\mathfrak{g}_6
=\bigcap\Bigl\{\mathfrak{h}\subseteq \mathfrak{gl}(V):\ \mathfrak{h}\ \text{为 Lie 子代数且}\ S\subseteq \mathfrak{h}\Bigr\},
\]
并且该交恰等于 \(\mathrm{Lie}\langle L_1,\dots,L_6\rangle\)。
\end{theorem}
\begin{proof}
任意 Lie 子代数族的交仍为 Lie 子代数，且显然包含 \(S\)，故上述交良定并给出包含 \(S\) 的极小性。
唯一性由“按包含关系取最小者”直接推出。证毕。
\leanverified{paper\_window6\_lie\_envelope\_existence\_uniqueness}
\end{proof}

\begin{proposition}[有限维情形的交换子闭包迭代终止]\label{prop:finite-dim-lie-closure-termination}
设 \(V\) 有限维，\(S\subset\mathfrak{gl}(V)\) 为有限集合。
令
\[
\mathfrak{g}^{(0)}:=\Span(S),\qquad
\mathfrak{g}^{(n+1)}:=\mathfrak{g}^{(n)}+\Span\bigl\{[A,B]:\ A\in S,\ B\in\mathfrak{g}^{(n)}\bigr\}.
\]
则存在 \(N\) 使得 \(\mathfrak{g}^{(N)}=\mathfrak{g}^{(N+1)}=\cdots\)，并且稳定值 \(\mathfrak{g}^{(N)}\) 等于定理 \ref{thm:window6-lie-envelope-existence-uniqueness} 所定义的最小 Lie 子代数。
\end{proposition}
\begin{proof}
由于 \(\dim\mathfrak{gl}(V)<\infty\)，序列 \(\dim\mathfrak{g}^{(n)}\) 为非降整数列且有上界，故必稳定：存在 \(N\) 使得 \(\dim\mathfrak{g}^{(N)}=\dim\mathfrak{g}^{(N+1)}\)。
由 \(\mathfrak{g}^{(N)}\subseteq\mathfrak{g}^{(N+1)}\) 得 \(\mathfrak{g}^{(N)}=\mathfrak{g}^{(N+1)}\)，从而对所有 \(n\ge N\) 恒等。\par
由构造可知 \(\mathfrak{g}^{(N)}\) 包含 \(S\)，且对与 \(S\) 的交换子封闭；结合 Jacobi 恒等式与线性闭包可推出对全体交换子封闭，故为 Lie 子代数。
又由于 \(\mathfrak{g}^{(N)}\) 由 \(S\) 的迭代闭包生成，满足包含 \(S\) 的极小性，从而与定理 \ref{thm:window6-lie-envelope-existence-uniqueness} 的刻画一致。证毕。
\leanverified{paper\_finite\_dim\_lie\_closure\_termination}
\leanverified{paper\_window6\_finite\_dim\_lie\_closure\_termination}
\end{proof}

\begin{lemma}[对称性]\label{lem:window6-li-symmetric}
对每个 \(i\in\{1,\dots,6\}\)，矩阵 \(L_i\) 为对称矩阵：\(L_i^{\mathsf T}=L_i\)。
\end{lemma}
\begin{proof}
映射 \(\omega\mapsto \omega\oplus e_i\) 为 \(\Omega_6\) 上的对合。
若 \(\omega\) 对 \((L_i)_{uv}\) 贡献一次计数，则 \(\omega\oplus e_i\) 对 \((L_i)_{vu}\) 贡献一次计数，故 \((L_i)_{uv}=(L_i)_{vu}\)。证毕。
\leanverified{paper\_window6\_li\_symmetric}
\end{proof}

\begin{lemma}[零对角与迹零]\label{lem:window6-li-traceless}
对每个 \(i\in\{1,\dots,6\}\)，有 \((L_i)_{uu}=0\ (\forall u\in X_6)\)，从而 \(\tr(L_i)=0\)。
因此 \(\mathfrak{g}_6\subseteq \mathfrak{sl}(V)\)。
\end{lemma}
\begin{proof}
记 \(n=\mathrm{int}_6(\omega)\in\{0,\dots,63\}\)。翻转第 \(i\) 位等价于 \(n\mapsto n\oplus 2^{6-i}\)，其差值的绝对值属于 \(\{1,2,4,8,16,32\}\)。
若存在 \(u\in X_6\) 使得 \((L_i)_{uu}>0\)，则存在 \(\omega\) 使 \(F_6(\omega)=F_6(\omega\oplus e_i)=u\)，等价于存在同一纤维中的两点 \(n,n\oplus 2^{6-i}\in \mathrm{Fold}^{\mathrm{bin}}_6{}^{-1}(u)\)。
由引理 \ref{lem:window6-foldbin6-fiber-difference-spectrum}，同一纤维中两点之差只能落在 \(\Delta_6=\{0,\pm 13,\pm 21,\pm 34,\pm 55\}\)。
由于 \(\{1,2,4,8,16,32\}\cap\{13,21,34,55\}=\varnothing\)，矛盾。
故 \((L_i)_{uu}=0\) 对所有 \(u\) 成立。
迹为零随即成立。\par
又因 \(\tr([A,B])=0\) 且 \(\mathfrak{sl}(V)\) 在交换子下封闭，故由 \(L_i\) 生成的 Lie 子代数包含于 \(\mathfrak{sl}(V)\)。证毕。
\leanverified{paper\_window6\_li\_traceless}
\end{proof}

\begin{remark}[对称/反对称分解与交换子闭合]\label{rem:window6-sl-decomposition-so-sym0}
设 \(\dim V=21\) 且特征不为 \(2\)。
记
\[
\mathrm{Sym}_0(V):=\{S\in\mathfrak{sl}(V):\ S^{\mathsf T}=S\},
\qquad
\mathfrak{so}(V):=\{K\in\mathfrak{sl}(V):\ K^{\mathsf T}=-K\}.
\]
则有向量空间直和分解
\[
\mathfrak{sl}(V)=\mathfrak{so}(V)\oplus \mathrm{Sym}_0(V),
\qquad
\dim\mathfrak{so}(V)=210,\ \dim \mathrm{Sym}_0(V)=230.
\]
并且交换子满足
\(
[\mathrm{Sym}_0,\mathrm{Sym}_0]\subset \mathfrak{so}
\)、
\(
[\mathfrak{so},\mathrm{Sym}_0]\subset \mathrm{Sym}_0
\)、
\(
[\mathfrak{so},\mathfrak{so}]\subset \mathfrak{so}
\)。
\end{remark}

\begin{theorem}[交换子派生部分的正交闭合（window-\(6\)）]\label{thm:window6-lie-envelope-so21-from-commutators}
在定义 \ref{def:window6-coordinate-pushforward-lie-envelope} 的设定下，令
\[
\mathfrak{k}_6:=\mathrm{Lie}\Bigl\langle [L_i,L_j]:\ 1\le i<j\le 6\Bigr\rangle\ \subset\ \mathfrak{gl}(V).
\]
则 \(\mathfrak{k}_6=\mathfrak{so}(V)\cong \mathfrak{so}(21)\)，并且 \(\dim\mathfrak{k}_6=210\)。
\leanverified{paper\_window6\_lie\_envelope\_so21\_from\_commutators}
\end{theorem}
\begin{proof}
由引理 \ref{lem:window6-li-symmetric}，对任意 \(i<j\) 有
\(
[L_i,L_j]^{\mathsf T}=-[L_i,L_j]
\)，故 \([L_i,L_j]\in\mathfrak{so}(V)\)。
因此 \(\mathfrak{k}_6\subseteq\mathfrak{so}(V)\)，从而 \(\dim\mathfrak{k}_6\le \dim\mathfrak{so}(V)=21\cdot 20/2=210\)。\par
另一方面，脚本 \texttt{\detokenize{scripts/exp_window6_lie_envelope_orthogonal_part.py}} 以有限穷举构造 \(L_i\)，并在模素数口径下对约化子代数 \((\mathfrak{k}_6)_p\subset \mathfrak{gl}(V_{\FF_p})\) 的闭包维数给出满维证书（见 \eqref{eq:window6_lie_envelope_orthogonal_part_dim}），即 \(\dim_{\FF_p}(\mathfrak{k}_6)_p=210\)。
令 \(V_{\QQ}:=\QQ\langle X_6\rangle\)，并记
\(
\mathfrak{k}_{6,\QQ}:=\Lie_{\QQ}\langle [L_i,L_j]:\,1\le i<j\le 6\rangle\subset\mathfrak{gl}(V_{\QQ})
\)。
由引理 \ref{lem:modp-lie-dimension-upperbound} 的维数提升不等式，得到 \(\dim_{\QQ}\mathfrak{k}_{6,\QQ}\ge 210\)，从而 \(\dim_{\RR}\mathfrak{k}_{6}=\dim_{\QQ}\mathfrak{k}_{6,\QQ}\ge 210\)。
故 \(\dim\mathfrak{k}_6=210\)，并由包含关系 \(\mathfrak{k}_6\subseteq\mathfrak{so}(V)\) 得 \(\mathfrak{k}_6=\mathfrak{so}(V)\)。证毕。
\end{proof}

\begin{lemma}[模素数线性独立性证书的提升]\label{lem:modp-independence-lift}
\leanverified{paper\_modp\_independence\_lift}
设 \(A_1,\dots,A_r\in \Mat_n(\ZZ)\)。
若存在素数 \(p\) 使得其约化 \(\overline{A}_1,\dots,\overline{A}_r\in \Mat_n(\FF_p)\) 在 \(\FF_p\) 上线性无关，则 \(A_1,\dots,A_r\) 在 \(\QQ\) 上线性无关。
\end{lemma}
\begin{proof}
若存在非平凡 \(\QQ\)-线性关系 \(\sum_i q_iA_i=0\)，乘以分母最小公倍数得非平凡整数关系 \(\sum_i c_iA_i=0\) 且 \(\gcd(c_1,\dots,c_r)=1\)。
模 \(p\) 化得到 \(\sum_i \overline{c}_i\,\overline{A}_i=0\)，其中至少一个 \(\overline{c}_i\neq 0\)，与 \(\{\overline{A}_i\}\) 的线性无关矛盾。证毕。
\end{proof}

\begin{lemma}[模素数维数证书的特征 \(0\) 提升]\label{lem:modp-lie-dimension-upperbound}
设 \(A_1,\dots,A_k\in\mathfrak{sl}(n,\ZZ)\)，并定义
\[
\mathfrak{g}_{\QQ}:=\Lie_{\QQ}\langle A_1,\dots,A_k\rangle\ \subset\ \mathfrak{sl}(n,\QQ).
\]
对任意素数 \(p\)，令 \(\overline{A}_i\) 为 \(A_i\) 在 \(\FF_p\) 上的约化，并定义
\[
\mathfrak{g}_{p}:=\Lie_{\FF_p}\langle \overline{A}_1,\dots,\overline{A}_k\rangle\ \subset\ \mathfrak{sl}(n,\FF_p).
\]
则有不等式
\[
\dim_{\FF_p}\mathfrak{g}_{p}\ \le\ \dim_{\QQ}\mathfrak{g}_{\QQ}.
\]
\end{lemma}
\begin{proof}
取 \(\mathfrak{g}_{\QQ}\) 的一组 \(\QQ\)-基并通分，可得一个自由 \(\ZZ\)-模 \(\Lambda\subset \mathfrak{sl}(n,\ZZ)\) 满足 \(\Lambda\otimes_{\ZZ}\QQ=\mathfrak{g}_{\QQ}\)。
将 \(\Lambda\) 模 \(p\) 约化得到 \(\overline{\Lambda}\subset\mathfrak{sl}(n,\FF_p)\)，其 \(\FF_p\)-线性张成维数至多为 \(\rank_{\ZZ}\Lambda=\dim_{\QQ}\mathfrak{g}_{\QQ}\)。
另一方面，\(\mathfrak{g}_{p}\) 由 \(\{\overline{A}_i\}\) 的 Lie 词张成，且每一 Lie 词都来自整数矩阵 Lie 词的约化，故包含于 \(\Span_{\FF_p}(\overline{\Lambda})\)。
于是 \(\dim_{\FF_p}\mathfrak{g}_{p}\le \dim_{\QQ}\mathfrak{g}_{\QQ}\)。证毕。
\end{proof}
\leanverified{paper\_modp\_lie\_dimension\_upperbound}

\begin{theorem}[window-\(6\) 推送交换子包络的结构识别]\label{thm:window6-lie-envelope-sl21}
在定义 \ref{def:window6-coordinate-pushforward-lie-envelope} 的设定下，\(|X_6|=21\)。
并且
\[
\mathfrak{g}_6=\mathfrak{sl}(V)\cong \mathfrak{sl}_{21}(\RR),
\qquad
\dim \mathfrak{g}_6=21^2-1=440.
\]
\end{theorem}
\begin{proof}
由引理 \ref{lem:window6-li-traceless}，\(\mathfrak{g}_6\subseteq\mathfrak{sl}(V)\)，故 \(\dim\mathfrak{g}_6\le \dim\mathfrak{sl}(V)=21^2-1=440\)。\par
令 \(V_{\QQ}:=\QQ\langle X_6\rangle\)，并记 \(\mathfrak{g}_{6,\QQ}:=\Lie_{\QQ}\langle L_1,\dots,L_6\rangle\subseteq \mathfrak{sl}(V_{\QQ})\)。
则 \(\mathfrak{g}_6=\mathfrak{g}_{6,\QQ}\otimes_{\QQ}\RR\)，从而 \(\dim_{\RR}\mathfrak{g}_6=\dim_{\QQ}\mathfrak{g}_{6,\QQ}\)。\par
取任意奇素数 \(p\)。
将整数矩阵 \(L_i\in \Mat_{21}(\ZZ)\) 模 \(p\) 约化得到 \(\overline{L}_i\in \Mat_{21}(\FF_p)\)，并记
\(
(\mathfrak{g}_6)_p:=\Lie_{\FF_p}\langle \overline{L}_1,\dots,\overline{L}_6\rangle\subseteq \mathfrak{sl}(21,\FF_p)
\)。
脚本 \texttt{\detokenize{scripts/exp_window6_lie_envelope.py}} 对 \((\mathfrak{g}_6)_p\) 的闭包维数给出证书化输出（见 \eqref{eq:window6_lie_envelope_invariants}），其结果为
\(
\dim_{\FF_p}(\mathfrak{g}_6)_p=440
\)。
由引理 \ref{lem:modp-lie-dimension-upperbound}，得到 \(\dim_{\QQ}\mathfrak{g}_{6,\QQ}\ge 440\)。
结合上界 \(\dim\mathfrak{g}_6\le 440\)，可得 \(\dim\mathfrak{g}_6=440\)，从而 \(\mathfrak{g}_6=\mathfrak{sl}(V)\)。证毕。
\end{proof}
\leanverified{paper\_window6\_lie\_envelope\_sl21}

\begin{remark}[结构层闭合的含义]\label{rem:window6-lie-envelope-compact-substructure}
定理 \ref{thm:window6-lie-envelope-sl21} 表明：若括号取交换子且闭包对象取定义 \ref{def:window6-coordinate-pushforward-lie-envelope} 的 \(\mathfrak{g}_6\)，则闭合后的 Lie 代数为满的 \(\mathfrak{sl}_{21}\)，而非低秩的 \(21\) 维简单代数。
因此，任何以 \(B_3/C_3\) 等根型为目标的结构性陈述，必须被理解为对 \(\mathfrak{sl}_{21}\) 内某一\emph{受限子代数}或\emph{可观测紧致部分}的抽取结论，其唯一性应来自离散分级与可审计选择律，而非来自交换子包络本身的闭合定义。
\end{remark}

\begin{remark}[关于“内部伴随实现”断言的必要修正]\label{rem:window6-internal-adjoint-realization-obstruction}
若存在一个以 \(V=\RR\langle X_6\rangle\) 为底空间的 \(21\) 维 Lie 代数结构 \((V,[-,-]_6)\)，使得每个 \(L_i\) 都可写为某个元素的伴随算子
\(
L_i=\mathrm{ad}_{x_i}
\)
（\(x_i\in V\)），则 \(\mathfrak{g}_6=\mathrm{Lie}\langle L_1,\dots,L_6\rangle\) 作为 \(\mathfrak{gl}(V)\) 的 Lie 子代数将包含于 \(\mathrm{ad}(V)\) 的像。
从而
\(
\dim\mathfrak{g}_6\le \dim V=21
\)，与定理 \ref{thm:window6-lie-envelope-sl21} 的 \(\dim\mathfrak{g}_6=440\) 矛盾。
因此，任何“\(L_i\) 由类型层内禀括号伴随实现”的主张，必须被理解为在附加约束（例如可观测截断、投影闭包或不变子代数抽取）后的受限版本，而不能与定义 \ref{def:window6-coordinate-pushforward-lie-envelope} 的无约束交换子闭包同时成立。
\end{remark}

\begin{theorem}[拉回括号：由交换子包络诱导类型层李代数]\label{thm:pullback-bracket-from-operator-envelope}
设 \(\mathfrak{g}\subseteq\mathfrak{gl}(V)\) 为 Lie 子代数，并给定线性同构
\(
\Phi:W\overset{\cong}{\to}\mathfrak{g}
\)。
定义
\[
[x,y]_W:=\Phi^{-1}\bigl([\Phi(x),\Phi(y)]\bigr),\qquad x,y\in W,
\]
其中右端括号为 \(\mathfrak{gl}(V)\) 中的交换子。
则 \((W,[-,-]_W)\) 为 Lie 代数，且 \(\Phi\) 为 Lie 代数同构。
\end{theorem}
\begin{proof}
双线性与反对称性由交换子与 \(\Phi\) 的线性性直接推出。
对任意 \(x,y,z\in W\)，由交换子在 \(\mathfrak{gl}(V)\) 中满足 Jacobi 恒等式，有
\[
\Phi\bigl([x,[y,z]_W]_W+[y,[z,x]_W]_W+[z,[x,y]_W]_W\bigr)
=[\Phi(x),[\Phi(y),\Phi(z)]]+[\Phi(y),[\Phi(z),\Phi(x)]]+[\Phi(z),[\Phi(x),\Phi(y)]]=0.
\]
由 \(\Phi\) 为同构得括号满足 Jacobi 恒等式。证毕。
\leanverified{paper\_pullback\_bracket\_from\_operator\_envelope}
\end{proof}

\begin{theorem}[根字典拉回的抽象括号：\texorpdfstring{$B_3/C_3$}{B3/C3} 型李代数结构的显式存在性]\label{thm:window6-root-dictionary-pullback-bracket}
记 boundary 扇区 \(X_6^{\mathrm{bdry}}\subset X_6\) 如推论 \ref{cor:terminal-foldbin6-boundary-pure-f9-alias}，并令 \(X_6^{\mathrm{cyc}}:=X_6\setminus X_6^{\mathrm{bdry}}\)。
设给定根字典双射 \(\iota_B:X_6^{\mathrm{cyc}}\to \Phi(B_3)\)（表 \ref{tab:fold6_b3c3_root_dictionary_b3}）或 \(\iota_C:X_6^{\mathrm{cyc}}\to \Phi(C_3)\)（表 \ref{tab:fold6_b3c3_root_dictionary_c3}），并取任一双射 \(\kappa:X_6^{\mathrm{bdry}}\to\{H_1,H_2,H_3\}\) 到 Cartan 基。
令 \(\mathfrak{g}_{B_3}\)（或 \(\mathfrak{g}_{C_3}\)）为相应紧致实型，其 Chevalley 基为 \(\{H_j\}_{j=1}^3\cup\{E_\alpha\}_{\alpha\in\Phi}\)。
定义线性同构 \(\varphi:V\to \mathfrak{g}_{B_3}\)（或 \(\mathfrak{g}_{C_3}\)）为
\[
\varphi(e_w)=
\begin{cases}
E_{\iota(w)},& w\in X_6^{\mathrm{cyc}},\\
\kappa(w),& w\in X_6^{\mathrm{bdry}}.
\end{cases}
\]
并在 \(V\) 上定义括号 \([x,y]_6:=\varphi^{-1}([\varphi(x),\varphi(y)])\)。
则 \((V,[-,-]_6)\) 为 \(21\) 维实 Lie 代数，其复化为 \(B_3\)（或 \(C_3\)）型简单 Lie 代数。
\end{theorem}
\begin{proof}
反对称性与 Jacobi 恒等式由 \(\mathfrak{g}_{B_3}\)（或 \(\mathfrak{g}_{C_3}\)）的李括号继承，且 \(\varphi\) 为线性同构，故拉回保持这些恒等式；根型由所选 Chevalley 基的结构给出。证毕。
\end{proof}
\leanverified{paper\_window6\_root\_dictionary\_pullback\_bracket}

\begin{theorem}[猜想 \ref{conj:window6-intrinsic-bracket-uniqueness} 的证书化版本]\label{thm:window6-intrinsic-bracket-certificate-theorem}
设给定根字典 \(\iota\) 与 Cartan 锚点双射 \(\kappa\) 如定理 \ref{thm:window6-root-dictionary-pullback-bracket}，并令 \([-,-]_6\) 为其拉回括号。
考虑如下三类可审计证书条件：
\begin{enumerate}
  \item \textbf{标号刚性：}类型邻接图刚性 \(\mathrm{Aut}(\Gamma_6)=\{1\}\) 与加权邻接保持置换的平凡性（命题 \ref{prop:terminal-gamma6-rigidity} 及其推论）。
  \item \textbf{分级相容：}boundary dyadic uplift 诱导的 sheet parity 分级与括号的相容性（推论 \ref{cor:terminal-delta-parity-rule}）。
  \item \textbf{分块相容：}orbit--Levi 骨架与 edge--flux 分块骨架对结构常数支撑与闭包模式的限制，与括号闭合相容。
\end{enumerate}
若上述证书条件均成立（其核验可由有限数据的确定性计算给出），则 \([-,-]_6\) 满足猜想 \ref{conj:window6-intrinsic-bracket-uniqueness} 所述的局域性/分级/分块约束，并且在保持这些审计语义不变的同构关系下是唯一的。
\end{theorem}
\begin{proof}
定理 \ref{thm:window6-root-dictionary-pullback-bracket} 已给出 \([-,-]_6\) 的存在性与 Jacobi 恒等式。
其余约束均为对结构常数支撑与分块闭包的有限判定条件：在固定基 \(X_6\) 下可以逐项核验。\par
唯一性来自标号刚性：若另一括号 \([-,-]_6'\) 与同一邻接/分级/分块审计语义相容，则允许的基变换必须诱导 \(\Gamma_6\) 的自同构并保持加权邻接语义；由 \(\mathrm{Aut}(\Gamma_6)=\{1\}\) 排除非平凡重标号，剩余的有限规范自由度再被分级与分块相容性条件锁定，从而在同构意义下唯一。证毕。
\end{proof}
\leanverified{paper\_window6\_intrinsic\_bracket\_certificate\_theorem}

\begin{conjecture}[标准模型规范代数的涌现刻画：作为 \texorpdfstring{$(\mathfrak{sl}_{21},\ZZ_2)$}{(sl21,Z2)} 的唯一可观测极小紧致子代数]\label{conj:window6-sm-gauge-as-minimal-compact-observable}
在定理 \ref{thm:window6-lie-envelope-sl21} 的设定下，并联立 window-\(6\) 的几何 \(\ZZ_2\) 稳定子（推论 \ref{cor:terminal-foldbin6-geo-stabilizer-z2}）及其在 boundary 上诱导的 sheet parity 分级（推论 \ref{cor:terminal-delta-parity-rule}），
在 \(\mathfrak{sl}(V)\) 内考虑所有与该 \(\ZZ_2\) 分级相容、并与 orbit--Levi 骨架及 edge--flux 分块接口一致的紧致约化子代数族。
则在“秩保持且不引入额外同构自由度”的极小化口径下，预期存在唯一（至同构）极小候选，其根型为
\[
A_2\oplus A_1\oplus A_0,
\qquad\text{亦即}\qquad
\mathfrak{su}(3)\oplus\mathfrak{su}(2)\oplus\mathfrak{u}(1).
\]
\end{conjecture}

\begin{remark}[从抽象括号到残余代数的约束性读出]\label{rem:window6-bracket-to-sm-residual}
在定理 \ref{thm:window6-lie-envelope-sl21} 的设定下，\(\mathfrak{sl}_{21}\) 仅提供算子链的闭包环境。
若进一步要求 orbit--Levi 骨架与阈值轻/重分裂在同一紧致包络内兼容闭合，则会强制出现 SM 型 Levi 子结构（命题 \ref{prop:window6-common-refinement-sm-levi}）。
几何 \(\ZZ_2\) 与 sheet parity 约束在该路径中提供“有效 \(u(1)\) 方向选出”的离散输入，并与猜想 \ref{conj:window6-sm-gauge-as-minimal-compact-observable} 的极小化口径相协调。
\end{remark}

\begin{conjecture}[内禀括号的存在性与唯一性：固定基 \texorpdfstring{$X_6$}{X6} 上的审计约束]\label{conj:window6-intrinsic-bracket-uniqueness}
设 \(V=\RR\langle X_6\rangle\)。
猜想存在反对称双线性运算
\[
[-,-]_6:\ V\times V\to V
\]
使得 \((V,[-,-]_6)\) 成为 Lie 代数，并满足以下可审计约束：
\begin{enumerate}
  \item \textbf{局域性/函子性：}结构常数仅依赖于由 \(Q_6\) 一跳推送诱导的邻接关系与边多重度；因此允许的类型重标号由类型邻接图刚性 \(\mathrm{Aut}(\Gamma_6)=\{1\}\) 强制消除（命题 \ref{prop:terminal-gamma6-rigidity}；亦见推论 \ref{cor:window6-type-superselection-autgamma6}）。
  \item \textbf{sheet parity 相容：}括号与 boundary dyadic uplift 所诱导的 \(\delta=34\) 二层覆盖分级相容（推论 \ref{cor:terminal-delta-parity-rule}）。
  \item \textbf{根--Cartan 接口：}在文稿固定的根--Cartan 字典接口下，cyclic 扇区与 boundary 扇区在 \([-,-]_6\) 下分别闭合为根向量生成子与 Cartan 子空间，并与秩 \(3\) 的根型审计接口相容。
  \item \textbf{Levi 分块相容：}在 orbit--Levi 骨架与 edge--flux 分块骨架的共同约束下，括号诱导的闭包保持分块语义并与 Levi 型一致性相容。
\end{enumerate}
并且在保持上述审计语义不变的同构关系下，该括号在同构意义下唯一。
\end{conjecture}

\begin{proposition}[抽象括号的存在/唯一性之证书化归约]\label{prop:window6-bracket-certificate-reduction}
在猜想 \ref{conj:window6-intrinsic-bracket-uniqueness} 的约束集合固定后，括号 \([-,-]_6\) 的存在性与唯一性等价于一组有限多项式方程的可解性与解唯一性判别：该方程组由
反对称性、Jacobi 恒等式以及局域/分块/parity 相容性所给出的结构常数约束构成。
此外，在要求结构常数仅依赖一跳推送邻接与边多重度的口径下，类型邻接图刚性 \(\mathrm{Aut}(\Gamma_6)=\{1\}\)（命题 \ref{prop:terminal-gamma6-rigidity}）排除任何通过类型重标号吸收非唯一性的自由度，从而“同构意义下唯一”可被压缩为给定标号下的代数唯一性。
\end{proposition}
\leanverified{paper\_window6\_bracket\_certificate\_reduction}

\begin{remark}[可审计证书接口]\label{rem:window6-lie-envelope-audit-route}
定义 \ref{def:window6-coordinate-pushforward-lie-envelope} 中的 \(L_i\) 完全由有限边计数确定，因此 \(\mathfrak{g}_6\) 的生成与闭包可被归约为有限矩阵代数的确定性计算。
脚本 \texttt{\detokenize{scripts/exp_window6_lie_envelope.py}} 输出包含 \(L_i\) 的完整计数矩阵与 \(\mathfrak{g}_6\) 的不变量，并生成可直接 \(\backslash\)input 的公式片段 \eqref{eq:window6_lie_envelope_invariants}。
\end{remark}

\begin{equation*}
\begin{aligned}
\dim\mathfrak{g}_6&=440,\\
\dim[\mathfrak{g}_6,\mathfrak{g}_6]&=440,\\
\dim Z(\mathfrak{g}_6)&=0,\\
\mathrm{rank}\,\kappa_6&=440,\\
\dim\,\mathrm{Comm}(\mathrm{ad}(\mathfrak{g}_6))&=1.
\end{aligned}
\end{equation*}
\begin{equation*}
\dim\mathfrak{k}_6=210.
\end{equation*}

\begin{remark}[从审计对象到定理链条的缺口定位]\label{rem:window6-bracket-closure-gaps}
在“交换子包络”层面，定理 \ref{thm:window6-lie-envelope-existence-uniqueness}--\ref{thm:window6-lie-envelope-sl21} 已将括号构造与闭合唯一性定理化，并将闭包结构识别为 \(\mathfrak{sl}_{21}\)。
在“抽象括号唯一性”层面，可将仍待补齐的内容压缩为两类可核验的构造/证书：
\begin{itemize}
  \item \textbf{类型层括号与算子交换子的桥接证书：}给出一个显式可核验的线性同构 \(\Phi:\RR\langle X_6\rangle\to\mathfrak{g}_6\subseteq\mathfrak{gl}(V)\)，并核验其为括号同态
  \(
  \Phi([x,y]_6)=[\Phi(x),\Phi(y)]
  \)
 ；一旦给出该证书，则定理 \ref{thm:pullback-bracket-from-operator-envelope} 将“类型层括号”为交换子诱导对象的主张形式化为定理。
  \item \textbf{审计约束的逐条验算：}在固定根--Cartan 字典、sheet parity 分级与分块骨架后，逐条核验猜想 \ref{conj:window6-intrinsic-bracket-uniqueness} 的局域性、分级与分块相容性条件；该验算等价于对有限个结构常数满足的多项式系统进行证书化判别（命题 \ref{prop:window6-bracket-certificate-reduction}）。
\end{itemize}
\end{remark}

\begin{remark}[Lee--Yang/Toeplitz--PSD 混合证书对紧致性与唯一性的增强]\label{rem:window6-leyang-toeplitz-mixed-certificates}
将 Lee--Yang 机制并入 window-\(6\) 的“整数—证书”体系后，与抽象括号链条相关的两类关键前提可以被重写为可复核的混合证书条件，从而避免将其作为外加的连续假设。

\paragraph{由可逆详细平衡核导出的紧性原则}
在 window-\(6\) 的推送核 \(P_6\) 口径下，详细平衡给出 \(\mathsf{T}_6\) 在 \(\ell^2(X_6,\pi_6)\) 上的自伴性，并且其交换子 \(*\)-代数 \(\mathcal C_6\) 的酉群 \(\mathrm U(\mathcal C_6)\) 为紧 Lie 群（推论 \ref{cor:window6-p6-compactness-principle}）。
因此，若将“可观测连续规范对称”限定为保持可逆读出核 \(P_6\)（等价于与 \(\mathsf{T}_6\) 交换）的连续对称，则紧致性由有限维详细平衡证书内生推出，从而排除以非紧实型为起点的分支歧义。

\paragraph{unitary slice 的零点几何译码与 Lee--Yang 条件}
自伴谱对象的代数化核验可通过 Joukowsky--Gödel 输运与 Lee--Yang 条件的单位圆零点刻画来组织。
在文稿的统一输运模板中，定理 \ref{thm:group-jg-elliptic-two-scale-stokes-mean} 给出 \(J_r\) 下的多项式拉回分解；
当所考虑的完成化多项式满足 Lee--Yang 条件（零点落在单位圆）时，可进一步得到椭圆输运的严格自由能恒等式（定理 \ref{thm:group-jg-leyang-mahler-free-energy}），从而将“幺正切片”要求等价转写为对整数/有理系数多项式之零点几何的可审计断言。

\paragraph{Toeplitz--PSD 的可编译化证书链}
零点几何断言可被进一步压缩为半正定证书链：在 Carath\'eodory 正性与 Toeplitz 主子矩阵半正定的等价框架下，盘内零点禁入可由有限级 Toeplitz--PSD 失败在有限维上直接见证（定理 \ref{thm:app-horizon-toeplitz-lmi}），并可通过反射系数与最小失败索引进行单标量定位（定理 \ref{thm:xi-toeplitz-det-verblunsky}）。
因此，若把猜想 \ref{conj:window6-intrinsic-bracket-uniqueness} 的“唯一性”从纯 Jacobi 多项式系统扩展为“Jacobi + Toeplitz--PSD”混合可行性，则其排他性可被提升为一个可复核的半代数证书问题：一方面由局域/分级/分块约束给出有限多项式条件（命题 \ref{prop:window6-bracket-certificate-reduction}），另一方面由 Toeplitz--PSD 层级给出可在有限维上核验的半正定约束。

\paragraph{挠子零点几何与整数锁定的交叉审计}
除谱侧证书外，挠子账本提供一条实现无关的零点几何交叉核验接口：有限 torsion 群 \(F\) 的圆维谱配分多项式 \(Z_F\) 的零点满足 Lee--Yang 型环带界（定理 \ref{thm:xi-cdim-spectrum-zeros-annulus}），并在同阶均匀 torsion 情形给出严格的单圆等角分布律（推论 \ref{cor:xi-cdim-spectrum-homocyclic-circle-law}）。
在 window-\(6\) 的耦合响应账本中，模 \(23\) 的零响应线唯一性（定理 \ref{thm:terminal-window6-coupling-mod23-kernel-line}）与相应 \(23\)-Sylow 证书共同给出一条纯整数的分支锁定输入；
将该 torsion 数据映射到 \(Z_F\) 的零点几何可形成互补约束，从而进一步压缩 \(u(1)\) 归一化与离散选择通道中的歧义空间。
\end{remark}

\begin{theorem}[由可逆核到 unitary slice 与 Toeplitz--PSD 的证书链（window-\(6\)）]\label{thm:window6-p6-toeplitz-certificate-chain}
\leanverified{paper\_window6\_p6\_toeplitz\_certificate\_chain}
沿用定理 \ref{thm:terminal-foldbin6-pushforward-markov} 的推送核 \(P_6\)、平稳分布 \(\pi_6\) 与 Markov 算子 \(\mathsf{T}_6\)。
则存在一条完全由有限对象诱导的证书链条，将“紧致连续包络”与“维数槽位逐项对齐的唯一性”置于同一可复核框架内：
\begin{enumerate}
  \item \textbf{可逆详细平衡 \(\Rightarrow\) 自伴与紧性原则.}
  \(\mathsf{T}_6\) 在 \(\ell^2(X_6,\pi_6)\) 上自伴，并且保持 \(P_6\)（等价保持 \(\mathsf{T}_6\)）的连续对称必为紧 Lie 群（推论 \ref{cor:window6-p6-compactness-principle}）。

  \item \textbf{自伴 Markov \(\Rightarrow\) 谱区间.}
  \(\mathsf{T}_6\) 的谱满足
  \[
  \mathrm{Spec}(\mathsf{T}_6)\subseteq[-1,1],
  \qquad\text{从而}\qquad
  \mathrm{Spec}(2\mathsf{T}_6)\subseteq[-2,2].
  \]

  \item \textbf{谱区间 \(\Rightarrow\) 单位圆零点证书.}
  令
  \[
  \chi_6(\lambda):=\det\!\bigl(\lambda I-2\mathsf{T}_6\bigr)\in\QQ[\lambda],
  \qquad n:=|X_6|=21.
  \]
  定义完成化多项式
  \[
  \widehat{\chi}_6(z):=z^{n}\,\chi_6(z+z^{-1})\in\QQ[z].
  \]
  则 \(\widehat{\chi}_6\) 的全部零点均满足 \(|z|=1\)。
  因而 unitary slice 条件可被等价改写为一个纯有理系数多项式的单位圆全根断言；并且清分母后可输出整系数证书。

  \item \textbf{单位圆零点 \(\Rightarrow\) Jensen 缺陷与半正定证书.}
  由于 \(\widehat{\chi}_6\) 在开圆盘 \(|z|<1\) 内无零点，故其 Jensen 缺陷在任意半径 \(0<\varrho<1\) 上恒为零（定义 \ref{def:app-jensen-defect} 与命题 \ref{prop:app-jensen-formula-defect} 的标准推论）。
  在本文的圆盘化完成化口径下，盘内零点禁入亦可等价组织为 Carath\'eodory 正性与 Toeplitz--PSD 层级（定理 \ref{thm:app-horizon-toeplitz-lmi}），从而得到一族有限维半正定核验接口。

  \item \textbf{紧性原则 + 维数逐项对齐 \(\Rightarrow\) 唯一因子锁定.}
  若进一步在边界塔审计口径下得到 \(D=12\) 的 Zeckendorf 原子性并要求逐项对齐，则紧李代数同构型被强制为
  \(
  \mathfrak u(1)\oplus\mathfrak{su}(2)\oplus\mathfrak{su}(3)
  \)
  且在逐项对齐意义下唯一（命题 \ref{prop:sm-zeckendorf-lie-algebra-rigidity}）。
\end{enumerate}
\end{theorem}

\begin{proof}
第 (1) 点已在定理 \ref{thm:terminal-foldbin6-pushforward-markov} 与推论 \ref{cor:window6-p6-compactness-principle} 中给出。
\par
对第 (2) 点：作为 Markov 算子，\(\mathsf{T}_6\) 在 \(\ell^1(X_6,\pi_6)\) 与 \(\ell^\infty(X_6)\) 上均为收缩映射（范数不增且 \(\mathsf{T}_6\mathbf 1=\mathbf 1\)）。
由 Riesz--Thorin 插值定理得 \(\|\mathsf{T}_6\|_{\ell^2(\pi_6)\to \ell^2(\pi_6)}\le 1\)。
结合自伴性，谱半径等于算子范数，从而 \(\mathrm{Spec}(\mathsf{T}_6)\subseteq[-1,1]\)。
\par
对第 (3) 点：由第 (2) 点，\(\chi_6\) 的全部零点 \(\lambda_j\) 属于 \([-2,2]\)。
对任意此类 \(\lambda\)，方程 \(z+z^{-1}=\lambda\) 等价于二次方程 \(z^2-\lambda z+1=0\)，其判别式 \(\lambda^2-4\le 0\)，且两根满足乘积为 \(1\) 并互为复共轭，故必有 \(|z|=1\)。
由 \(\widehat{\chi}_6(z)=0\iff \chi_6(z+z^{-1})=0\)（注意 \(z\neq 0\)），得 \(\widehat{\chi}_6\) 的全部零点均落在单位圆上。
\par
第 (4) 点的 Jensen 缺陷结论为命题 \ref{prop:app-jensen-formula-defect} 对一般圆盘解析函数的直接推论；Toeplitz--PSD 层级在本文口径下由定理 \ref{thm:app-horizon-toeplitz-lmi} 给出一种可复核的半正定核验范式。
\par
第 (5) 点为命题 \ref{prop:sm-zeckendorf-lie-algebra-rigidity} 的直接应用。
证毕。
\end{proof}

\begin{remark}[链条中已作为正文结论出现的环节与标准推论环节]\label{rem:window6-p6-toeplitz-chain-what-is-internal}
定理 \ref{thm:window6-p6-toeplitz-certificate-chain} 的第 (1) 点与第 (5) 点分别对应文稿中的推送核可逆性/紧性原则接口（定理 \ref{thm:terminal-foldbin6-pushforward-markov}、推论 \ref{cor:window6-p6-compactness-principle}）与 Zeckendorf 原子性下的 SM 刚性（命题 \ref{prop:sm-zeckendorf-lie-algebra-rigidity}）。
第 (2)--(4) 点为在本文圆盘化与 Toeplitz--PSD 证书语言中常用的标准推论：由 Markov 收缩与自伴性得到谱区间界，再由 \(z+z^{-1}\) 完成化把区间谱改写为单位圆零点，最后用 Jensen 缺陷与 Toeplitz--PSD 层级给出等价或互补的有限维核验接口（定义 \ref{def:app-jensen-defect}、命题 \ref{prop:app-jensen-formula-defect}、定理 \ref{thm:app-horizon-toeplitz-lmi}）。
\end{remark}

% This file is intended to be \input{} from subsubsec__window6_lie_bracket_closure.tex.

\paragraph{Lee--Yang 与整数证书链的合流：紧性与唯一性的可核验化}
在 window-\(6\) 的审计范式中，Lee--Yang 机制并非作为外加类比引入，而是作为一条与“整数—可审计证书”同型的约束通道：它把若干原本依赖于结构选择的环节改写为可复核的不变量断言，从而将紧致性/幺正切片与唯一性收束从假设推进为推论。
其核心作用可归纳为两类提升：其一以可逆详细平衡核为起点内生推出紧性原则与 unitary slice 的多项式证书；其二把“括号唯一性/极小紧致子代数”的排他性从纯 Jacobi 多项式系统增强为带半正定约束的混合证书问题。

\begin{theorem}[由可逆核到 unitary slice 的多项式证书与 Toeplitz--PSD]\label{thm:window6-leyang-p6-unitary-toeplitz}
\leanverified{paper\_window6\_leyang\_p6\_unitary\_toeplitz}
在定理 \ref{thm:terminal-foldbin6-pushforward-markov} 与推论 \ref{cor:window6-p6-compactness-principle} 的设定下，\(\mathsf{T}_6\) 的自伴性与紧性原则可被推进为一条 unitary slice 与半正定证书链条：存在一个由 \(\mathsf{T}_6\) 决定的有理系数完成化多项式 \(\widehat{\chi}_6(z)\)，其全部零点位于单位圆 \(|z|=1\)；
并且在本文的圆盘化完成化接口下，该单位圆零点断言可等价组织为 Jensen 缺陷恒零与 Toeplitz--PSD 层级可行性（定义 \ref{def:app-jensen-defect}、命题 \ref{prop:app-jensen-formula-defect}、定理 \ref{thm:app-horizon-toeplitz-lmi}）。
\end{theorem}

\begin{proof}
该结论的完整链式证明已在定理 \ref{thm:window6-p6-toeplitz-certificate-chain} 中给出；
其中“可逆详细平衡 \(\Rightarrow\) 自伴 \(\Rightarrow\) 紧性原则”来自推论 \ref{cor:window6-p6-compactness-principle}，
而“自伴谱区间 \(\Rightarrow\) 单位圆零点完成化 \(\Rightarrow\) Jensen/Toeplitz--PSD”属于圆盘化完成化框架下的标准推论。
\end{proof}

\begin{remark}[紧致性作为推论条件]\label{rem:window6-compactness-not-assumption}
由于推论 \ref{cor:window6-p6-compactness-principle} 直接把“保持 \(P_6\)”的连续对称压缩为紧 Lie 群，故后续以紧李代数分类为口径的识别论证无需额外引入“选取紧实型/幺正切片”为外加假设；
相应的切片条件亦可由定理 \ref{thm:window6-leyang-p6-unitary-toeplitz} 的单位圆零点与 Toeplitz--PSD 证书链离线复核。
\end{remark}

\begin{remark}[Lee--Yang 判别式的整数锁定与额外约束的必要性]\label{rem:window6-leyang-discriminant-necessity}
在本文的输运—消元框架中，谱碰撞点与分支点由判别式零点刻画，并且在参数连续延拓的意义下给出标号延拓的最大半径控制（定理 \ref{thm:group-jg-discriminant-factorization} 及其推论）。
当谱对象来自整系数（或有理系数清分母后的整系数）多项式时，判别式不仅提供“是否发生碰撞”的等价判据，同时把允许的碰撞位置压缩为一组由整数因子与坏素数集合所控制的算术指纹。\par
这一点与 window-\(6\) 的两条对象链形成互补：在推送交换子包络口径下，\(\mathfrak g_6=\mathfrak{sl}(V)\) 的最大闭合由有限矩阵代数的证书化审计给出（定理 \ref{thm:window6-lie-envelope-sl21}），因而任何低秩紧致规范代数都不可能从“仅含 \(\{L_i\}\) 的交换子闭包”中直接出现（命题 \ref{prop:window6-lowrank-needs-observable-compression}）。
相反，若希望在可观测层面稳定地产生低秩紧致结构，则必须引入额外的压缩/兼容约束以选出受限子代数；而 Lee--Yang 判别式机制提供了一条实现无关的算术证据，说明“可观测选择律”不能仅由最大闭合的存在性替代：碰撞/分支的出现被整数判别式严格锁定，从而将允许的非阿贝尔增强通道压缩为离散且可复核的算术条件集合。
\end{remark}

\begin{proposition}[Zeckendorf 逐项对齐下的 SM 同构型锁定]\label{prop:window6-sm-uniqueness-from-zeckendorf-under-compactness}
\leanverified{paper\_window6\_sm\_uniqueness\_from\_zeckendorf\_under\_compactness}
在 Zeckendorf 原子性与逐项对齐的审计口径下，若可观测紧致部分的维数审计为 \(D=12\)，则其紧李代数同构型必为
\[
\mathfrak u(1)\oplus \mathfrak{su}(2)\oplus \mathfrak{su}(3),
\]
并且在逐项对齐意义下唯一。
\end{proposition}

\begin{proof}
该断言即命题 \ref{prop:sm-zeckendorf-lie-algebra-rigidity} 的直接重述。
\end{proof}

\begin{remark}[中心三轴与 Pati--Salam 连通实现的既有闭合接口]\label{rem:window6-center-ps-chain-existing}
window-\(6\) 的 boundary 读出在 dyadic uplift 口径下诱导三轴中心寄存器 \((\ZZ/2\ZZ)^3\)。
文稿已在群论层面给出四条与“全局群形态”相关的可引用闭合接口：
\begin{itemize}
  \item \(\Lie(G)\cong \mathfrak u(1)\oplus\mathfrak{su}(2)\oplus\mathfrak{su}(3)\) 的紧连通包络在中心 \(2\)-挠层面无法容纳 \((\ZZ/2\ZZ)^3\) 的单射嵌入（定理 \ref{thm:xi-window6-sm-envelope-center-2torsion-obstruction}）。
  \item 更强地，任意连通紧单 Lie 群都不可能满足 boundary 中心兼容性；对连通紧半单包络，中心 \(2\)-挠秩至少必须满足 \(2t+u\ge 3\) 的必要条件（定理 \ref{thm:window6-boundary-center-compatibility-excludes-compact-simple}）。
  \item 在 Pati--Salam 覆盖 \(\widetilde G_{\mathrm{PS}}=SU(4)\times SU(2)_L\times SU(2)_R\) 上，若要求 canonical involution 子群在商映射下仍保真下降，则中心商 \(H\) 必平凡，从而连通实现唯一（定理 \ref{thm:xi-window6-ps-faithful-canonical-descent-unique}）。
  \item 任意满足 \(\Lie(G_{\mathrm{SM}})\cong \mathfrak u(1)\oplus\mathfrak{su}(2)\oplus\mathfrak{su}(3)\) 的紧连通群作为包络时，都不可避免地压缩 Pati--Salam 的中心 \(2\)-挠通道（定理 \ref{thm:xi-window6-ps-to-sm-center2-kernel}）。
\end{itemize}
这些结论与 \(\delta=34\) 的 sheet parity 锚点性（例如命题 \ref{prop:xi-window6-sheet-parity-anchor-m6}）共同给出一条实现无关的全局排他性：中心通道与分级通道的自由度均被离散证书锁死，不能通过中心商或重标号吸收。
\end{remark}

\begin{remark}[torsion 的 Lee--Yang 型零点几何指纹]\label{rem:window6-torsion-zeros-geometry}
有限交换群的圆维谱配分多项式 \(Z_F\) 把不变量因子端点转写为零点模长的环带界（定理 \ref{thm:xi-cdim-spectrum-zeros-annulus}），并在同阶均匀 torsion 情形给出严格的单圆等角分布律（推论 \ref{cor:xi-cdim-spectrum-homocyclic-circle-law}）。
因此，中心/挠结构既可在群论层面以 \(2\)-挠秩与中心商核检验，也可在多项式层面以零点半径与等角分布提供独立的几何核验接口，从而形成互不依赖的双证书链。
\end{remark}

\begin{theorem}[括号唯一性的混合证书化归约：Jacobi 与 Toeplitz--PSD 的交汇]\label{thm:window6-bracket-uniqueness-mixed-certificate-reduction}
\leanverified{paper\_window6\_bracket\_uniqueness\_mixed\_certificate\_reduction}
在猜想 \ref{conj:window6-intrinsic-bracket-uniqueness} 的设定下，固定由一跳推送邻接与边多重度所决定的局域支撑模板，并联立 \(\delta=34\) 的 parity 分级封闭与 Levi 分块相容性约束。
则任一候选括号的“存在/唯一性”可被组织为一个混合证书可行性问题：
\begin{enumerate}
  \item \textbf{代数侧：}反对称性与 Jacobi 恒等式在有限未知结构常数上给出有限多项式约束，并由局域支撑模板与分级/分块约束进一步稀疏化（命题 \ref{prop:window6-bracket-certificate-reduction}）。
  \item \textbf{半正定侧：}若从候选括号函子化地产生一族完成化多项式或等价的圆盘解析对象，并要求其满足 unitary slice 的“盘内零点禁入”，则该条件可在本文口径下等价改写为 Toeplitz--PSD 层级的半正定可行性（定理 \ref{thm:app-horizon-toeplitz-lmi}）。
  \item \textbf{共终端核验：}由于 Toeplitz 主子矩阵按包含关系单调（大阶矩阵的主子矩阵仍为主子矩阵），故
  \[
  (\forall N)\ \mathbf T_N\succeq 0
  \quad\Longleftrightarrow\quad
  (\exists N_k\nearrow\infty)\ (\forall k)\ \mathbf T_{N_k}\succeq 0,
  \]
  即“全层级”可在任意共终端子序列上核验。
\end{enumerate}
此外，类型邻接图刚性 \(\mathrm{Aut}(\Gamma_6)=\{1\}\) 排除通过类型重标号吸收非唯一性的自由度（命题 \ref{prop:terminal-gamma6-rigidity}），从而“同构意义下唯一”可被压缩为固定标号下的混合系统唯一可解性。
\end{theorem}

\begin{proof}
第 (1) 点为命题 \ref{prop:window6-bracket-certificate-reduction} 的直接应用。
第 (2) 点为本文圆盘化完成化与 Toeplitz--PSD 证书语言的既有接口：盘内零点禁入可由 Carath\'eodory 正性与 Toeplitz--PSD 层级制度化（定理 \ref{thm:app-horizon-toeplitz-lmi}）。
第 (3) 点为半正定在主子矩阵下的遗传性：若 \(\mathbf T_M\succeq 0\) 且 \(M\ge N\)，则 \(\mathbf T_N\) 作为 \(\mathbf T_M\) 的主子矩阵亦半正定；因此对任意 \(N\) 取 \(k\) 使 \(N_k\ge N\) 即得右蕴含左。
最后，标号刚性将同构自由度消去为有限规范自由度，故唯一性判别等价于固定标号下的可行集是否为单点。证毕。
\end{proof}

% This file is intended to be \input{} from subsubsec__window6_lie_bracket_closure.tex.

\paragraph{根--Cartan 字典与边界塔维数接口的合流：从离散骨架到唯一闭合}
在既定的 window-\(6\) 锚定口径下，“李括号的构造”与“唯一闭合”可被组织为一组可验证的定理化结论，其关键并非引入额外连续结构，而在于把连续李代数所需的输入锚定到离散可审计数据上，并将剩余自由度压缩到有限证书问题。

\begin{remark}[离散 Cartan/根骨架：\texorpdfstring{$18\oplus 3$}{18⊕3} 刚性分裂与根字典输入]\label{rem:window6-rootcartan-18plus3-skeleton}
window-\(6\) 的稳定类型集合满足刚性分解 \(X_6=X_6^{\mathrm{cyc}}\sqcup X_6^{\mathrm{bdry}}\) 且 \(|X_6^{\mathrm{cyc}}|=18\)、\(|X_6^{\mathrm{bdry}}|=3\)。
在根--Cartan 字典口径下，将 boundary 三元组视为 Cartan 锚点、将 cyclic 的 \(18\) 个类型视为根向量标签，则根字典给出一种完全显式的秩 \(3\) 根型输入，并在 \(V=\RR\langle X_6\rangle\) 上诱导紧致 \(B_3/C_3\) 型李括号结构的存在性（定理 \ref{thm:window6-root-dictionary-pullback-bracket}）。
\end{remark}

\begin{remark}[可见规范代数的维数槽位锁定：\texorpdfstring{$(1,3,8)$}{(1,3,8)} \(\Rightarrow\) \texorpdfstring{$\mathfrak u(1)\oplus\mathfrak{su}(2)\oplus\mathfrak{su}(3)$}{u(1)+su(2)+su(3)}]\label{rem:window6-boundary-tower-sm-uniqueness}
边界塔在 \((m=4,6,8)\) 的维数接口给出 \((|X_4^{\mathrm{bdry}}|,|X_6^{\mathrm{bdry}}|,|X_8^{\mathrm{bdry}}|)=(1,3,8)\)。
在 Zeckendorf 原子性与逐项对齐的审计口径下，若可观测连续包络取紧约化，则维数槽位的同构型被唯一锁定为 \(\mathfrak u(1)\oplus\mathfrak{su}(2)\oplus\mathfrak{su}(3)\)（命题 \ref{prop:sm-zeckendorf-lie-algebra-rigidity}）。
\end{remark}

\begin{remark}[高分辨率边界 uplift 的根型分块与对角 \texorpdfstring{$\mathfrak{su}(3)$}{su(3)} 的强迫]\label{rem:window6-m10-a2pm-diagonal-su3}
在 \(m=10\) 的边界 uplift 字典下，根方向的刚性四块分裂把边界层 \(21\) 模态组织为
\[
21\ =\ 6\ (\mathrm{short})\ \oplus\ 6\ (A_2(+))\ \oplus\ 6\ (A_2(-))\ \oplus\ 3\ (\mathrm{Cartan})
\]
且两套 \(A_2(\pm)\) 的判别仅依赖符号结构（定理 \ref{thm:boundary-m10-b3c3-rigid-four-block-split}；表 \ref{tab:boundary_uplift_m10_dictionary}）。
在等权对称被提升为协议不变量并要求其在连续包络上可见时，交换对合的固定点机制排除两份独立的 \(\mathfrak{su}(3)\) 因子在可见层面并存，从而强迫仅保留对角 \(\mathfrak{su}(3)\)（命题 \ref{prop:a2pm-isotropy-diagonal-su3}）。
\end{remark}

上述离散骨架还可进一步凝聚为四条完全结构性的 window-\(6\) 定理：它们分别刻画 cyclic 扇区中 Hamming 权阈值对根长轨道的决定性、\(21=18+3\) 分解对伴随权多重集的精确实现、奇偶荷格中 boundary parity 的零权本质，以及 \(C_6\) 轨道分解与 \(B_3\) 长短根分区之间的唯一交会方式。
\begin{theorem}[window-\(6\) 循环扇区中 Hamming 权阈值对根长 Weyl 轨道的刻画]\label{thm:window6-cyclic-weight-threshold-root-length}
沿用 \S\ref{subsec:bdry-tower-zeck-gut} 的记号，并记
\[
X_6=X_6^{\mathrm{cyc}}\sqcup X_6^{\mathrm{bdry}},
\qquad
|X_6|=21,\quad |X_6^{\mathrm{cyc}}|=18,\quad |X_6^{\mathrm{bdry}}|=3.
\]
cyclic 扇区的 Hamming 权直方图满足
\[
\mathrm{wt}\text{-hist}(X_6^{\mathrm{cyc}})=\{0:1,\ 1:6,\ 2:9,\ 3:2\},
\]
因而恰有六个词满足 \(\mathrm{wt}(w)=1\)，其余十二个词满足 \(\mathrm{wt}(w)\neq 1\)。
则在表 \ref{tab:fold6_b3c3_root_dictionary_b3} 的 \(B_3\) 规范下，\(\mathrm{wt}(w)=1\) 的六个词恰对应短根
\[
\{\pm e_1,\pm e_2,\pm e_3\},
\]
而 \(\mathrm{wt}(w)\neq 1\) 的十二个词恰对应长根
\[
\{\pm e_i\pm e_j:\ 1\le i<j\le 3\}.
\]
在表 \ref{tab:fold6_b3c3_root_dictionary_c3} 的 \(C_3\) 规范下，上述两块仅交换长短根名称：\(\mathrm{wt}(w)=1\) 的六个词恰对应长根
\[
\{\pm 2e_1,\pm 2e_2,\pm 2e_3\},
\]
而 \(\mathrm{wt}(w)\neq 1\) 的十二个词恰对应短根
\[
\{\pm e_i\pm e_j:\ 1\le i<j\le 3\}.
\]
因此 \(X_6^{\mathrm{cyc}}\) 在 \(W(B_3)=W(C_3)\) 的作用下恰分解为两条 Weyl 轨道，轨道大小分别为 \(6\) 与 \(12\)，并由阈值条件 \(\mathrm{wt}=1\) 与 \(\mathrm{wt}\neq 1\) 精确切出。
\end{theorem}
\begin{proof}
表 \ref{tab:fold6_b3c3_root_dictionary_b3} 中 \(\mathrm{wt}=1\) 的六行恰为
\[
\texttt{100000},\ \texttt{010000},\ \texttt{001000},\ \texttt{000100},\ \texttt{000010},\ \texttt{000001},
\]
其像分别为
\[
(1,0,0),\ (0,1,0),\ (0,0,1),\ (0,0,-1),\ (0,-1,0),\ (-1,0,0),
\]
即 \(\{\pm e_1,\pm e_2,\pm e_3\}\)。
另一方面，唯一的 \(\mathrm{wt}=0\) 词 \(\texttt{000000}\) 被送到 \((1,1,0)\)，已经落在
\(\{\pm e_i\pm e_j\}\) 中；其余 \(\mathrm{wt}=2,3\) 的十一行正好给出该集合中的其余十一条向量。
故在 \(B_3\) 规范下，\(\mathrm{wt}=1\) 与 \(\mathrm{wt}\neq 1\) 分别对应短根与长根。

对 \(C_3\) 规范，表 \ref{tab:fold6_b3c3_root_dictionary_c3} 仅将轴向六词统一改写为
\(\{\pm 2e_1,\pm 2e_2,\pm 2e_3\}\)，而其余十二条向量仍为 \(\{\pm e_i\pm e_j\}\)，故结论中的长短根名称恰好交换。

最后，\(W(B_3)=W(C_3)\) 作为 signed permutation 群分别在两类根长集合上传递，因此其拉回到 \(X_6^{\mathrm{cyc}}\) 上正给出大小为 \(6\) 与 \(12\) 的两条 Weyl 轨道。证毕。
\end{proof}

\begin{theorem}[window-\(6\) 的 \texorpdfstring{$21=18+3$}{21=18+3} 分解精确实现 \texorpdfstring{$B_3/C_3$}{B3/C3} 伴随权多重集]\label{thm:window6-21-adjoint-weight-multiset}
在表 \ref{tab:fold6_b3c3_root_dictionary_b3} 与表 \ref{tab:fold6_b3c3_root_dictionary_c3} 的两种规范下，集合 \(X_6\) 都可被规范识别为相应秩 \(3\) 半单 Lie 代数的伴随权多重集：
\[
X_6\ \cong\ \Phi(B_3)\sqcup \{0,0,0\},
\qquad
X_6\ \cong\ \Phi(C_3)\sqcup \{0,0,0\}.
\]
其中 cyclic 扇区 \(X_6^{\mathrm{cyc}}\) 给出全部 \(18\) 条根，而 boundary 三词
\[
X_6^{\mathrm{bdry}}=\{\texttt{100001},\texttt{100101},\texttt{101001}\}
\]
给出三重零权。换言之，window-\(6\) 的 \(21\) 个稳定类型不是“根与外点”的松散并置，而是一个完全显式的秩 \(3\) 伴随权模型。
\end{theorem}
\begin{proof}
两张根字典表分别列出 \(X_6^{\mathrm{cyc}}\) 的 \(18\) 个词及其像，并且像恰为 \(\Phi(B_3)\) 或 \(\Phi(C_3)\) 的全部 \(18\) 条根。
因此 cyclic 扇区在两种规范下都与相应根系完全一致。
另一方面，boundary 扇区恰由三条词
\(\texttt{100001},\texttt{100101},\texttt{101001}\)
组成，且与 cyclic 扇区不交。
于是把这三条 boundary 词规范地视为零权，即得
\[
X_6\cong \Phi(B_3)\sqcup\{0,0,0\}
\qquad\text{或}\qquad
X_6\cong \Phi(C_3)\sqcup\{0,0,0\}.
\]
对秩 \(3\) 的 \(B_3\) 与 \(C_3\)，伴随表示的权多重集正由全部根与三重零权组成，故结论成立。证毕。
\end{proof}

\paragraph{伴随权模型的二阶与四阶不变量}
为把上述 \(18+3\) 权重模型进一步提升为可计算张量不变量，定义两种权映射
\[
\Lambda_B:X_6\longrightarrow \ZZ^3,
\qquad
\Lambda_C:X_6\longrightarrow \ZZ^3
\]
如下：在 cyclic 扇区 \(X_6^{\mathrm{cyc}}\) 上，分别取表 \ref{tab:fold6_b3c3_root_dictionary_b3} 与表 \ref{tab:fold6_b3c3_root_dictionary_c3} 中给出的根向量；在 boundary 三元组 \(X_6^{\mathrm{bdry}}\) 上，统一规定
\[
\Lambda_B(w)=0,
\qquad
\Lambda_C(w)=0.
\]
由定理 \ref{thm:window6-21-adjoint-weight-multiset}，\(\Lambda_B(X_6)\) 与 \(\Lambda_C(X_6)\) 作为多重集分别恰为
\[
\Phi(B_3)\sqcup \{0,0,0\},
\qquad
\Phi(C_3)\sqcup \{0,0,0\}.
\]

\begin{theorem}[window-\(6\) 伴随权模型的二次各向同性与精确 frame 常数]\label{thm:window6-b3c3-adjoint-second-moment-isotropy}
沿用上述 \(\Lambda_B,\Lambda_C\) 的记号。则有精确二阶矩恒等式
\[
\sum_{w\in X_6}\Lambda_B(w)\Lambda_B(w)^{\mathsf T}=10\,I_3,
\qquad
\sum_{w\in X_6}\Lambda_C(w)\Lambda_C(w)^{\mathsf T}=16\,I_3.
\]
因此，\(X_6\) 在两种规范下都给出 \(\RR^3\) 上的各向同性紧框架；其归一化 frame 常数分别为
\[
\frac{10}{3}
\qquad\text{与}\qquad
\frac{16}{3}.
\]
等价地，对任意 \(u\in \RR^3\)，都有
\[
\sum_{w\in X_6}\langle \Lambda_B(w),u\rangle^2=10\,|u|^2,
\qquad
\sum_{w\in X_6}\langle \Lambda_C(w),u\rangle^2=16\,|u|^2.
\]
\end{theorem}
\begin{proof}
boundary 三元组被送到零向量，故不贡献二阶矩；只需对 \(18\) 个根求和。

在 \(B_3\) 规范下，短根部分为
\[
\{\pm e_1,\pm e_2,\pm e_3\},
\]
故
\[
\sum_{\alpha\in \{\pm e_1,\pm e_2,\pm e_3\}}
\alpha\alpha^{\mathsf T}
=2I_3.
\]
长根部分为
\[
\{\pm e_i\pm e_j:\ 1\le i<j\le 3\}.
\]
对固定 \(i<j\)，四个向量
\[
\pm e_i\pm e_j
\]
满足
\[
\sum_{\epsilon_i,\epsilon_j=\pm 1}
(\epsilon_i e_i+\epsilon_j e_j)(\epsilon_i e_i+\epsilon_j e_j)^{\mathsf T}
=4(E_{ii}+E_{jj}),
\]
交叉项由符号对称性抵消。对三对 \((i,j)\) 求和后，每个对角元恰出现两次，故长根总贡献为
\[
8I_3.
\]
于是
\[
\sum_{w\in X_6}\Lambda_B(w)\Lambda_B(w)^{\mathsf T}=2I_3+8I_3=10I_3.
\]

在 \(C_3\) 规范下，长根部分为
\[
\{\pm 2e_1,\pm 2e_2,\pm 2e_3\},
\]
其贡献为
\[
\sum_{i=1}^{3}\bigl((2e_i)(2e_i)^{\mathsf T}+(-2e_i)(-2e_i)^{\mathsf T}\bigr)=8I_3.
\]
其余十二条短根仍为 \(\{\pm e_i\pm e_j\}\)，贡献同上仍为 \(8I_3\)。故
\[
\sum_{w\in X_6}\Lambda_C(w)\Lambda_C(w)^{\mathsf T}=8I_3+8I_3=16I_3.
\]
对任意 \(u\in\RR^3\) 取二次型，即得最后一组等式。证毕。
\end{proof}

\begin{theorem}[window-\(6\) 伴随权模型的指数特征函数与四次 Weyl 不变量]\label{thm:window6-b3c3-adjoint-character-and-quartic-invariant}
沿用 \(\Lambda_B,\Lambda_C\) 的记号。定义指数和
\[
\Theta_B(t):=\sum_{w\in X_6}e^{\langle \Lambda_B(w),t\rangle},
\qquad
\Theta_C(t):=\sum_{w\in X_6}e^{\langle \Lambda_C(w),t\rangle},
\qquad
t=(t_1,t_2,t_3)\in \RR^3.
\]
则有完全显式公式
\[
\Theta_B(t)
=
3+2\sum_{i=1}^{3}\cosh t_i
+4\sum_{1\le i<j\le 3}\cosh t_i\,\cosh t_j,
\]
\[
\Theta_C(t)
=
3+2\sum_{i=1}^{3}\cosh(2t_i)
+4\sum_{1\le i<j\le 3}\cosh t_i\,\cosh t_j.
\]
因此二者分别是 \(B_3\)、\(C_3\) 伴随权多重集在分裂 Cartan 上的精确指数特征函数。

进一步，对任意 \(u=(u_1,u_2,u_3)\in\RR^3\)，令 \(t=su\)。则有四次矩闭式
\[
\sum_{w\in X_6}\langle \Lambda_B(w),u\rangle^4
=
10\sum_{i=1}^{3}u_i^4+24\sum_{1\le i<j\le 3}u_i^2u_j^2,
\]
\[
\sum_{w\in X_6}\langle \Lambda_C(w),u\rangle^4
=
40\sum_{i=1}^{3}u_i^4+24\sum_{1\le i<j\le 3}u_i^2u_j^2.
\]
\end{theorem}
\begin{proof}
由定理 \ref{thm:window6-21-adjoint-weight-multiset}，
\[
\Theta_B(t)=3+\sum_{\alpha\in \Phi(B_3)}e^{\langle \alpha,t\rangle},
\qquad
\Theta_C(t)=3+\sum_{\beta\in \Phi(C_3)}e^{\langle \beta,t\rangle}.
\]
把 \(B_3\) 根系拆成短根 \(\{\pm e_i\}\) 与长根 \(\{\pm e_i\pm e_j\}\)，则
\[
\sum_{\alpha\in \{\pm e_1,\pm e_2,\pm e_3\}}e^{\langle \alpha,t\rangle}
=2\sum_{i=1}^{3}\cosh t_i,
\]
而对每个 \(i<j\)，有
\[
e^{t_i+t_j}+e^{t_i-t_j}+e^{-t_i+t_j}+e^{-t_i-t_j}
=4\cosh t_i\,\cosh t_j.
\]
对三对 \((i,j)\) 求和，即得 \(\Theta_B\) 的闭式。

对 \(C_3\) 规范，仅把轴向六根改为 \(\{\pm 2e_i\}\)，故
\[
\sum_{\beta\in \{\pm 2e_1,\pm 2e_2,\pm 2e_3\}}e^{\langle \beta,t\rangle}
=2\sum_{i=1}^{3}\cosh(2t_i),
\]
而其余十二根仍给出同样的
\(
4\sum_{i<j}\cosh t_i\cosh t_j
\)，遂得 \(\Theta_C\) 的闭式。

现令 \(t=su\)。利用展开式
\[
\cosh(su_i)=1+\frac{s^2u_i^2}{2}+\frac{s^4u_i^4}{24}+O(s^6),
\]
\[
\cosh(2su_i)=1+2s^2u_i^2+\frac{2}{3}s^4u_i^4+O(s^6),
\]
代入上述两条闭式并比较 \(s^4\) 项，可得
\[
\Theta_B(su)
=
21+5s^2|u|^2
+s^4\!\left(
\frac{5}{12}\sum_{i=1}^{3}u_i^4
+\sum_{1\le i<j\le 3}u_i^2u_j^2
\right)+O(s^6),
\]
\[
\Theta_C(su)
=
21+8s^2|u|^2
+s^4\!\left(
\frac{5}{3}\sum_{i=1}^{3}u_i^4
+\sum_{1\le i<j\le 3}u_i^2u_j^2
\right)+O(s^6).
\]
另一方面，由定义
\[
\Theta_B(su)=\sum_{w\in X_6}\exp\!\bigl(s\langle \Lambda_B(w),u\rangle\bigr),
\]
\[
\Theta_C(su)=\sum_{w\in X_6}\exp\!\bigl(s\langle \Lambda_C(w),u\rangle\bigr),
\]
而权集关于原点成对对称，故奇次项消失，且 \(s^4\) 系数分别等于
\[
\frac{1}{24}\sum_{w\in X_6}\langle \Lambda_B(w),u\rangle^4,
\qquad
\frac{1}{24}\sum_{w\in X_6}\langle \Lambda_C(w),u\rangle^4.
\]
将两边比较，即得所述四次矩闭式。证毕。
\end{proof}

\paragraph{三次 Euclidean cubature 与唯一四次调和检测器}
在上述 \(B_3/C_3\) 伴随权模型中，三重 boundary 零权并不只提供静态权多重集分解；它还把 window-\(6\) 的 \(21\) 点配置提升为一对三次精确的 Euclidean cubature，并使四次层面的全部非径向误差收束到唯一一条 \(W(B_3)\)-不变调和方向。

\begin{theorem}[\texorpdfstring{$B_3$}{B3} 归一化下的三次精确 Euclidean cubature]\label{thm:window6-b3-degree3-euclidean-cubature}
\leanverified{paper\_window6\_b3\_degree3\_euclidean\_cubature}
沿用 \(\Lambda_B\) 的记号，并把 \(X_6\) 的像视为带重数的离散配置
\[
\mathcal X_B:=\Lambda_B(X_6)
=
\{0,0,0\}\sqcup \{\pm e_1,\pm e_2,\pm e_3\}\sqcup \{\pm e_i\pm e_j:\ 1\le i<j\le 3\}.
\]
记 \(\sigma_r\) 为半径 \(r\) 的球面概率测度，并定义径向参考测度
\[
\nu_B:=\frac{3}{21}\delta_0+\frac{6}{21}\sigma_1+\frac{12}{21}\sigma_{\sqrt2}.
\]
则对任意总次数不超过 \(3\) 的实多项式 \(P\in\RR[x_1,x_2,x_3]\)，成立严格恒等式
\[
\frac1{21}\sum_{x\in \mathcal X_B}P(x)=\int_{\RR^3}P\,\dd\nu_B.
\]
\end{theorem}
\begin{proof}
只需在总次数不超过 \(3\) 的单项式基上验证。

常数项显然成立。又由于 \(\mathcal X_B\) 关于原点中心对称，任意奇次单项式在离散平均下都为零；而 \(\nu_B\) 是径向测度，因此其奇次矩同样全部为零。故只需检查二次矩。

由定理 \ref{thm:window6-b3c3-adjoint-second-moment-isotropy}，
\[
\sum_{x\in\mathcal X_B}xx^{\mathsf T}=10I_3.
\]
因此
\[
\frac1{21}\sum_{x\in\mathcal X_B}x_ix_j=0
\qquad(i\neq j),
\]
并且
\[
\frac1{21}\sum_{x\in\mathcal X_B}x_i^2=\frac{10}{21}
\qquad(i=1,2,3).
\]
另一方面，由球面对称性，
\[
\int x_ix_j\,\dd\sigma_r=0\qquad(i\neq j),
\qquad
\int x_i^2\,\dd\sigma_r=\frac{r^2}{3}.
\]
于是
\[
\int x_i^2\,\dd\nu_B
=
\frac{6}{21}\cdot\frac13+\frac{12}{21}\cdot\frac23
=
\frac{10}{21},
\]
且交叉矩同样为零。

因此离散平均与 \(\nu_B\) 的全部 \(0,1,2,3\) 阶矩一致，所示 cubature 恒等式成立。证毕。
\end{proof}

\begin{theorem}[\texorpdfstring{$C_3$}{C3} 归一化下的三次精确 Euclidean cubature]\label{thm:window6-c3-degree3-euclidean-cubature}
\leanverified{paper\_window6\_c3\_degree3\_euclidean\_cubature}
沿用 \(\Lambda_C\) 的记号，并把 \(X_6\) 的像视为带重数的离散配置
\[
\mathcal X_C:=\Lambda_C(X_6)
=
\{0,0,0\}\sqcup \{\pm 2e_1,\pm 2e_2,\pm 2e_3\}\sqcup \{\pm e_i\pm e_j:\ 1\le i<j\le 3\}.
\]
定义径向参考测度
\[
\nu_C:=\frac{3}{21}\delta_0+\frac{6}{21}\sigma_{2}+\frac{12}{21}\sigma_{\sqrt2}.
\]
则对任意总次数不超过 \(3\) 的实多项式 \(P\in\RR[x_1,x_2,x_3]\)，成立
\[
\frac1{21}\sum_{x\in \mathcal X_C}P(x)=\int_{\RR^3}P\,\dd\nu_C.
\]
\end{theorem}
\begin{proof}
证明与定理 \ref{thm:window6-b3-degree3-euclidean-cubature} 同样，只需比较二次矩。

由定理 \ref{thm:window6-b3c3-adjoint-second-moment-isotropy}，
\[
\sum_{x\in\mathcal X_C}xx^{\mathsf T}=16I_3.
\]
故
\[
\frac1{21}\sum_{x\in\mathcal X_C}x_ix_j=0
\qquad(i\neq j),
\]
并且
\[
\frac1{21}\sum_{x\in\mathcal X_C}x_i^2=\frac{16}{21}
\qquad(i=1,2,3).
\]
另一方面，
\[
\int x_i^2\,\dd\nu_C
=
\frac{6}{21}\cdot\frac{4}{3}+\frac{12}{21}\cdot\frac{2}{3}
=
\frac{16}{21},
\]
且交叉矩同样因径向对称而为零。

于是离散平均与 \(\nu_C\) 的全部 \(0,1,2,3\) 阶矩一致，结论成立。证毕。
\end{proof}

\begin{theorem}[四次误差的秩一性与唯一的 \texorpdfstring{$W(B_3)$}{W(B3)}-不变调和检测器]\label{thm:window6-b3c3-quartic-rankone-harmonic-detector}
\leanverified{paper\_window6\_b3c3\_quartic\_rankone\_harmonic\_detector}
定义
\[
A:=x_1^4+x_2^4+x_3^4,
\qquad
B:=x_1^2x_2^2+x_1^2x_3^2+x_2^2x_3^2,
\]
\[
H_4:=A-3B,
\qquad
R_4:=(x_1^2+x_2^2+x_3^2)^2=A+2B.
\]
又定义四次误差泛函
\[
\mathcal D_B(P):=\frac1{21}\sum_{x\in\mathcal X_B}P(x)-\int_{\RR^3}P\,\dd\nu_B,
\]
\[
\mathcal D_C(P):=\frac1{21}\sum_{x\in\mathcal X_C}P(x)-\int_{\RR^3}P\,\dd\nu_C.
\]
则：
\begin{enumerate}
  \item \(H_4\) 是 \(W(B_3)\)-不变调和四次多项式；
  \item 在 \(W(B_3)\)-不变四次多项式空间
  \[
  \mathrm{span}\{R_4,H_4\}
  \]
  上，\(\mathcal D_B\) 与 \(\mathcal D_C\) 都是秩一泛函；
  \item 它们满足精确公式
  \[
  \mathcal D_B(R_4)=0,\qquad \mathcal D_B(H_4)=-\frac{2}{7},
  \]
  \[
  \mathcal D_C(R_4)=0,\qquad \mathcal D_C(H_4)=4.
  \]
\end{enumerate}
因此，三次精确性在四次层面只沿唯一一条 \(W(B_3)\)-不变调和方向 \(H_4\) 失效，而纯径向方向 \(R_4\) 仍被完全积分。
\end{theorem}
\begin{proof}
首先，
\[
\Delta A=12(x_1^2+x_2^2+x_3^2),
\qquad
\Delta B=4(x_1^2+x_2^2+x_3^2),
\]
故
\[
\Delta H_4=\Delta(A-3B)=0.
\]
又 \(A,B\) 都是 signed permutation 对称的，故 \(H_4\) 为 \(W(B_3)\)-不变调和四次多项式。

其次，任意 \(W(B_3)\)-不变四次多项式都可唯一写成 \(\alpha A+\beta B\)。等价地，该二维空间亦可由
\[
R_4=A+2B,
\qquad
H_4=A-3B
\]
作为一组基。

现计算 \(\mathcal D_B\)。在 \(\mathcal X_B\) 上，六个短根给出 \(A=1,B=0\)，十二个长根给出 \(A=2,B=1\)，三重原点给出 \(A=B=0\)。因此
\[
\frac1{21}\sum_{x\in\mathcal X_B}A(x)=\frac{6\cdot 1+12\cdot 2}{21}=\frac{10}{7},
\qquad
\frac1{21}\sum_{x\in\mathcal X_B}B(x)=\frac{12}{21}=\frac{4}{7}.
\]
另一方面，对半径为 \(r\) 的球面概率测度有
\[
\int x_i^4\,\dd\sigma_r=\frac{r^4}{5},
\qquad
\int x_i^2x_j^2\,\dd\sigma_r=\frac{r^4}{15}\quad(i\neq j).
\]
故
\[
\int_{\RR^3}A\,\dd\nu_B
=
\frac{6}{21}\cdot\frac{3}{5}
\,+\,
\frac{12}{21}\cdot\frac{12}{5}
=
\frac{54}{35},
\]
\[
\int_{\RR^3}B\,\dd\nu_B
=
\frac{6}{21}\cdot\frac{1}{5}
\,+\,
\frac{12}{21}\cdot\frac{4}{5}
=
\frac{18}{35}.
\]
从而
\[
\mathcal D_B(A)=\frac{10}{7}-\frac{54}{35}=-\frac{4}{35},
\qquad
\mathcal D_B(B)=\frac{4}{7}-\frac{18}{35}=\frac{2}{35}.
\]
于是
\[
\mathcal D_B(R_4)=\mathcal D_B(A+2B)=0,
\qquad
\mathcal D_B(H_4)=\mathcal D_B(A-3B)=-\frac{2}{7}.
\]

再计算 \(\mathcal D_C\)。在 \(\mathcal X_C\) 上，六个轴向长根给出 \(A=16,B=0\)，十二个短根仍给出 \(A=2,B=1\)，故
\[
\frac1{21}\sum_{x\in\mathcal X_C}A(x)=\frac{6\cdot 16+12\cdot 2}{21}=\frac{40}{7},
\qquad
\frac1{21}\sum_{x\in\mathcal X_C}B(x)=\frac{12}{21}=\frac{4}{7}.
\]
而
\[
\int_{\RR^3}A\,\dd\nu_C
=
\frac{6}{21}\cdot\frac{48}{5}
\,+\,
\frac{12}{21}\cdot\frac{12}{5}
=
\frac{144}{35},
\]
\[
\int_{\RR^3}B\,\dd\nu_C
=
\frac{6}{21}\cdot\frac{16}{5}
\,+\,
\frac{12}{21}\cdot\frac{4}{5}
=
\frac{48}{35}.
\]
因此
\[
\mathcal D_C(A)=\frac{40}{7}-\frac{144}{35}=\frac{8}{5},
\qquad
\mathcal D_C(B)=\frac{4}{7}-\frac{48}{35}=-\frac{4}{5}.
\]
遂有
\[
\mathcal D_C(R_4)=\mathcal D_C(A+2B)=0,
\qquad
\mathcal D_C(H_4)=\mathcal D_C(A-3B)=4.
\]

因为 \(\mathcal D_B,\mathcal D_C\) 都在二维空间 \(\mathrm{span}\{R_4,H_4\}\) 上消去 \(R_4\) 而不消去 \(H_4\)，故二者在该空间上的秩都恰为 \(1\)。证毕。
\end{proof}

\begin{corollary}[双归一化四次判别器的唯一性]\label{cor:window6-b3c3-unique-quartic-detector}
\leanverified{paper\_window6\_b3c3\_unique\_quartic\_detector}
在 \(W(B_3)\)-不变多项式代数的次数不超过 \(4\) 的部分中，同时满足
\[
\mathcal D_B(P)=0,
\qquad
\mathcal D_C(P)=0
\]
的多项式空间恰为
\[
\mathrm{span}\bigl\{1,\ x_1^2+x_2^2+x_3^2,\ (x_1^2+x_2^2+x_3^2)^2\bigr\}.
\]
等价地，任一 \(W(B_3)\)-不变四次多项式 \(P\) 都有唯一分解
\[
P=c_4R_4+c_HH_4,
\]
并满足
\[
\mathcal D_B(P)=-\frac{2}{7}c_H,
\qquad
\mathcal D_C(P)=4c_H.
\]
因此，在次数不超过 \(4\) 的 \(W(B_3)\)-不变世界中，唯一能够同时检测两种归一化偏离纯径向积分的方向正是 \(H_4\)。
\end{corollary}
\begin{proof}
由定理 \ref{thm:window6-b3-degree3-euclidean-cubature} 与定理 \ref{thm:window6-c3-degree3-euclidean-cubature}，所有次数不超过 \(3\) 的多项式都已落在 \(\ker \mathcal D_B\cap \ker \mathcal D_C\) 中。

在四次 \(W(B_3)\)-不变空间中，由定理 \ref{thm:window6-b3c3-quartic-rankone-harmonic-detector}，
\[
\mathcal D_B(R_4)=\mathcal D_C(R_4)=0,
\]
而
\[
\mathcal D_B(H_4)\neq 0,
\qquad
\mathcal D_C(H_4)\neq 0.
\]
故
\[
\ker \mathcal D_B\cap \ker \mathcal D_C
\]
在该二维空间上恰等于 \(\mathrm{span}\{R_4\}\)。

再把低次不变量
\[
1,\qquad x_1^2+x_2^2+x_3^2
\]
加入，即得完整交核空间
\[
\mathrm{span}\bigl\{1,\ x_1^2+x_2^2+x_3^2,\ R_4\bigr\}.
\]
最后，将任意四次不变量写成 \(P=c_4R_4+c_HH_4\)，直接代入上一条定理中的数值公式，即得
\[
\mathcal D_B(P)=-\frac{2}{7}c_H,
\qquad
\mathcal D_C(P)=4c_H.
\]
证毕。
\end{proof}

\paragraph{带 boundary 复制方向的五次球面 cubature 与六次首缺陷}
在三次 Euclidean cubature 与四次调和检测器之外，window-\(6\) 的 \(21\) 个标号可见点在单位球面上还携带一条更刚性的五次 cubature 结构。其关键差别在于：上一节处理的是 \(19\) 个几何不同的归一化支撑，而这里保留 boundary 三词作为沿正坐标轴的复制方向，因此必须把 \(21\) 个标号点整体视为一族带重数的球面原子。

为此，沿用表 \ref{tab:fold6_b3c3_root_dictionary_b3} 的 \(B_3\) 字典，并记与 \(\pm e_i\) 对应的 cyclic 短根词为
\[
s_1^+=\texttt{100000},\qquad s_2^+=\texttt{010000},\qquad s_3^+=\texttt{001000},
\]
\[
s_1^-=\texttt{000001},\qquad s_2^-=\texttt{000010},\qquad s_3^-=\texttt{000100}.
\]
又把 boundary 三词按
\[
b_1=\texttt{100001},\qquad b_2=\texttt{100101},\qquad b_3=\texttt{101001}
\]
排序。对每个 \(w\in X_6\) 定义单位可见向量
\[
u_w:=\frac{v_w}{|v_w|}\in S^2,
\]
其中 \(v_w=\Lambda_B(w)\) 若 \(w\in X_6^{\mathrm{cyc}}\)，而 \(v_{b_i}=e_i\)。
于是 \(\widetilde U_6:=\{u_w:\ w\in X_6\}\) 是一个带三重重复方向
\[
\{e_1,e_2,e_3\}
\]
的 \(21\) 点标号球面支撑。以下记 \(\sigma\) 为 \(S^2\) 上的球面概率测度，并称有限加权原子测度
\[
\mu=\sum_{w\in X_6}c_w\delta_{u_w}
\]
在 \(\widetilde U_6\) 上对次数不超过 \(t\) 的球面多项式给出精确 cubature，若对任意次数 \(\le t\) 的多项式 \(P\) 都有
\[
\sum_{w\in X_6}c_w P(u_w)=\mu(S^2)\int_{S^2}P(u)\,\dd\sigma(u).
\]

\begin{theorem}[window-\(6\) 标号球面支撑上的五次 cubature 族]\label{thm:window6-labeled-spherical-degree5-cubature-family}
设
\[
\mu=\sum_{w\in X_6}c_w\delta_{u_w}
\]
是支撑于 \(\widetilde U_6\) 的有限加权原子测度。则 \(\mu\) 对 \(S^2\) 的全部次数 \(\le 5\) 多项式给出精确球面 cubature，当且仅当存在参数
\[
\lambda,t_1,t_2,t_3\in \CC
\]
使得
\[
c_w=\lambda
\qquad
\text{对全部 \(12\) 个长根方向},
\]
\[
c_{s_i^-}=\frac{\lambda}{2},
\qquad
c_{s_i^+}=\frac{\lambda}{2}-t_i,
\qquad
c_{b_i}=t_i
\qquad (i=1,2,3).
\]
特别地，总质量恒为
\[
\mu(S^2)=15\lambda.
\]
因此，window-\(6\) 的 \(21\) 点标号球面支撑上的五次 cubature 并非离散孤点，而是一个四维仿射族；其真正模只来自三条正轴复制方向上的质量转移。
\end{theorem}
\begin{proof}
对每个 \(i\) 记
\[
m_i^+:=c_{s_i^+}+c_{b_i},
\qquad
m_i^-:=c_{s_i^-}.
\]
由于 \(u_{s_i^+}=u_{b_i}=e_i\)，球面矩条件只依赖于 \(m_i^+\)，而不依赖于它在 \((s_i^+,b_i)\) 之间如何分拆。
因此可先在 \(18\) 个几何不同的方向上分类所有五次 cubature，再把 \(m_i^+\) 在 \((s_i^+,b_i)\) 间自由分配。

对每个坐标平面 \(ij\in\{12,13,23\}\)，记四个长根方向
\[
\frac{\pm e_i\pm e_j}{\sqrt2}
\]
上的权为
\[
\lambda_{ij}^{++},\qquad
\lambda_{ij}^{+-},\qquad
\lambda_{ij}^{-+},\qquad
\lambda_{ij}^{--}.
\]
由于球面概率测度的一切奇次矩都为零，取混合矩
\[
x_i x_j,\qquad x_i x_j^2,\qquad x_i^2 x_j
\]
便得到
\[
\lambda_{ij}^{++}-\lambda_{ij}^{+-}-\lambda_{ij}^{-+}+\lambda_{ij}^{--}=0,
\]
\[
\lambda_{ij}^{++}+\lambda_{ij}^{+-}-\lambda_{ij}^{-+}-\lambda_{ij}^{--}=0,
\]
\[
\lambda_{ij}^{++}-\lambda_{ij}^{+-}+\lambda_{ij}^{-+}-\lambda_{ij}^{--}=0.
\]
解此线性系统得
\[
\lambda_{ij}^{++}=\lambda_{ij}^{+-}=\lambda_{ij}^{-+}=\lambda_{ij}^{--}=:\lambda_{ij}.
\]
故每个坐标平面中的四个长根必须等权。

再看一阶矩。由于三组长根在每个坐标方向上的总贡献已成对相消，取线性函数 \(x_i\) 即得
\[
m_i^+=m_i^-
\qquad (i=1,2,3).
\]

最后使用四阶 isotropy。球面概率测度满足
\[
\int_{S^2}x_i^4\,\dd\sigma=3\int_{S^2}x_i^2x_j^2\,\dd\sigma
\qquad(i\ne j).
\]
而在当前离散测度下，
\[
M_{iijj}:=\sum_{w\in X_6}c_w u_{w,i}^2u_{w,j}^2=\lambda_{ij},
\]
\[
M_{iiii}:=\sum_{w\in X_6}c_w u_{w,i}^4
=m_i^++m_i^-+\lambda_{ij}+\lambda_{ik}
=2m_i^++\lambda_{ij}+\lambda_{ik}.
\]
因此
\[
2m_i^+ + \lambda_{ij}+\lambda_{ik}=3\lambda_{ij}=3\lambda_{ik}.
\]
由此推出
\[
\lambda_{12}=\lambda_{13}=\lambda_{23}=:\lambda,
\qquad
m_i^+=m_i^-=\frac{\lambda}{2}.
\]
于是全部长根权都等于 \(\lambda\)，全部负短根权都等于 \(\lambda/2\)，而每条正轴上的总质量都等于 \(\lambda/2\)。
把该总质量在 \((s_i^+,b_i)\) 间分拆为
\[
c_{s_i^+}=\frac{\lambda}{2}-t_i,
\qquad
c_{b_i}=t_i,
\]
便得到所述参数化。总质量随即为
\[
12\lambda+3\cdot\frac{\lambda}{2}+3\cdot\frac{\lambda}{2}=15\lambda.
\]

反过来，若权满足题设公式，则一切奇次矩都因长根的符号对称与轴向的平衡条件
\[
m_i^+=m_i^-=\frac{\lambda}{2}
\]
而消失。对二次矩，三个坐标轴方向总贡献为 \(\lambda I_3\)，三组长根方向总贡献为 \(4\lambda I_3\)，故
\[
\sum_{w\in X_6}c_w\langle x,u_w\rangle^2=5\lambda |x|^2.
\]
对四次矩，三个坐标轴方向总贡献为
\[
\lambda(x_1^4+x_2^4+x_3^4),
\]
而每个坐标平面的四个长根方向给出
\[
\lambda\bigl(x_i^4+6x_i^2x_j^2+x_j^4\bigr).
\]
对三对 \((i,j)\) 求和后得到
\[
\sum_{w\in X_6}c_w\langle x,u_w\rangle^4=3\lambda |x|^4.
\]
因此
\[
\sum_{w\in X_6}c_w\langle x,u_w\rangle^k
=
15\lambda\int_{S^2}\langle x,u\rangle^k\,\dd\sigma(u)
\qquad (k=0,1,2,3,4,5),
\]
其中用到了
\[
\int_{S^2}\langle x,u\rangle^2\,\dd\sigma(u)=\frac{|x|^2}{3},
\qquad
\int_{S^2}\langle x,u\rangle^4\,\dd\sigma(u)=\frac{|x|^4}{5},
\]
以及奇次矩为零。由极化恒等式，这等价于对全部次数 \(\le 5\) 多项式的精确 cubature。证毕。
\end{proof}
\leanverified{paper\_window6\_labeled\_spherical\_degree5\_cubature\_family}

\begin{proposition}[正权归一化五次 cubature 的立方体模空间]\label{prop:window6-labeled-spherical-degree5-cubature-cube}
在定理 \ref{thm:window6-labeled-spherical-degree5-cubature-family} 的四维仿射族中，正权且归一化
\[
\sum_{w\in X_6}c_w=1
\]
的部分恰为闭立方体
\[
\lambda=\frac{1}{15},
\qquad
0\le t_i\le \frac{1}{30}
\qquad (i=1,2,3).
\]
故 window-\(6\) 的正权五次球面 cubature 模空间同构于
\[
{\Bigl[0,\frac{1}{30}\Bigr]}^{3}.
\]
其 \(8\) 个顶点给出全部支撑极小的正权解；在这些顶点上，支撑恰有 \(18\) 个点，并且顶点集天然构成一个 \(\ZZ_2^3\)-torsor。

此外，在整个 \(21\) 点标号支撑上不存在非平凡等权五次球面设计。
\end{proposition}
\begin{proof}
由定理 \ref{thm:window6-labeled-spherical-degree5-cubature-family}，总质量恒为 \(15\lambda\)，故归一化立即给出
\[
\lambda=\frac{1}{15}.
\]
正权条件等价于
\[
\lambda>0,
\qquad
\frac{\lambda}{2}-t_i\ge 0,
\qquad
t_i\ge 0.
\]
代入 \(\lambda=1/15\) 即得
\[
0\le t_i\le \frac{1}{30}.
\]
故正权归一化模空间正是三维闭立方体。

在该立方体中，\(12\) 个长根权恒为 \(1/15\)，三条负短根权恒为 \(1/30\)，而每条正轴上的总质量 \(1/30\) 在 \((s_i^+,b_i)\) 间分配。若 \(t_i\) 处于开区间，则 \(s_i^+\) 与 \(b_i\) 两个标号点都带正权；若 \(t_i\in\{0,1/30\}\)，则恰有一个点的权消失。故支撑最小要求三个 \(t_i\) 同时落在端点，从而一共保留
\[
12+3+3=18
\]
个非零权点。三条正轴各有两种选择，因而顶点集与三维超立方体的顶点集自然同构，即构成一个 \(\ZZ_2^3\)-torsor。

最后，若全部 \(21\) 个标号点都带同一非零权 \(c\)，则定理 \ref{thm:window6-labeled-spherical-degree5-cubature-family} 强迫
\[
c=\lambda
\qquad\text{在全部长根方向上},
\]
同时又有
\[
c=\frac{\lambda}{2}
\qquad\text{在三条负短根上},
\]
矛盾。故非平凡等权五次球面设计不存在。证毕。
\end{proof}

\begin{theorem}[整个五次 cubature 立方体共享同一六次谐和缺陷]\label{thm:window6-labeled-spherical-degree6-harmonic-defect}
设 \(\mu\) 是定理 \ref{thm:window6-labeled-spherical-degree5-cubature-family} 的任意五次精确 cubature。定义
\[
A_6(x):=x_1^6+x_2^6+x_3^6,
\qquad
B_6(x):=\sum_{i\ne j}x_i^4x_j^2,
\qquad
C_6(x):=x_1^2x_2^2x_3^2,
\]
以及六次谐和多项式
\[
\mathscr H_6(x):=
2A_6(x)-15B_6(x)+180C_6(x).
\]
则对任意 \(x\in \RR^3\)，有精确恒等式
\[
\sum_{w\in X_6}c_w\langle x,u_w\rangle^6
=
\frac{15\lambda}{7}|x|^6-\frac{\lambda}{14}\mathscr H_6(x),
\]
其中 \(\lambda\) 为定理 \ref{thm:window6-labeled-spherical-degree5-cubature-family} 中的长根公共权。
在归一化情形 \(\lambda=1/15\) 下，公式化为
\[
\sum_{w\in X_6}c_w\langle x,u_w\rangle^6
=
\frac{1}{7}|x|^6-\frac{1}{210}\mathscr H_6(x).
\]
特别地，六次矩缺陷与参数 \((t_1,t_2,t_3)\) 无关；整个五次 cubature 立方体在六次上共享同一个首个不可消去 defect。
\end{theorem}
\begin{proof}
由定理 \ref{thm:window6-labeled-spherical-degree5-cubature-family}，三条正轴上的总质量始终满足
\[
c_{s_i^+}+c_{b_i}=\frac{\lambda}{2},
\]
故六次矩只依赖于 \(\lambda\)，与 \((t_1,t_2,t_3)\) 如何分配无关。

沿三个坐标轴的总贡献为
\[
\lambda A_6(x).
\]
另一方面，对固定的 \(i<j\)，四个长根方向给出
\[
\lambda\sum_{\epsilon_i,\epsilon_j=\pm 1}
\left\langle x,\frac{\epsilon_i e_i+\epsilon_j e_j}{\sqrt2}\right\rangle^{6}
=
\frac{\lambda}{8}
\sum_{\epsilon_i,\epsilon_j=\pm 1}
{\bigl(\epsilon_i x_i+\epsilon_j x_j\bigr)}^{6}
=
\frac{\lambda}{2}\bigl(x_i^6+15x_i^4x_j^2+15x_i^2x_j^4+x_j^6\bigr).
\]
对三对 \((i,j)\) 求和后得到
\[
\sum_{w\in X_6}c_w\langle x,u_w\rangle^6
=
2\lambda A_6(x)+\frac{15\lambda}{2}B_6(x).
\]

又由
\[
|x|^6=A_6(x)+3B_6(x)+6C_6(x)
\]
可知
\[
\frac{15\lambda}{7}|x|^6-\frac{\lambda}{14}\mathscr H_6(x)
=
\lambda\left(
\frac{15}{7}(A_6+3B_6+6C_6)-\frac{1}{14}(2A_6-15B_6+180C_6)
\right),
\]
整理即得
\[
\frac{15\lambda}{7}|x|^6-\frac{\lambda}{14}\mathscr H_6(x)
=
2\lambda A_6(x)+\frac{15\lambda}{2}B_6(x).
\]
这就证明了题设恒等式。

最后，直接计算拉普拉斯算子可得
\[
\Delta \mathscr H_6=0,
\]
故 \(\mathscr H_6\) 的确是六次谐和多项式。归一化情形只需代入 \(\lambda=1/15\)。证毕。
\end{proof}

\begin{corollary}[window-\(6\) 标号球面支撑的 cubature 强度恰为 \(5\)]\label{cor:window6-labeled-spherical-cubature-strength-five}
window-\(6\) 的 \(21\) 点标号球面支撑 \(\widetilde U_6\) 上存在非平凡的五次精确球面 cubature；但不存在任何总质量非零、支撑于同一 \(21\) 点标号支撑的六次精确球面 cubature。
因此其最大精确阶恰为
\[
t_{\max}=5.
\]
\end{corollary}
\begin{proof}
存在性已由定理 \ref{thm:window6-labeled-spherical-degree5-cubature-family} 给出。
为证六次不可能，只需取球面六次单项式
\[
P(x):=x_1^2x_2^2x_3^2.
\]
对每个 \(w\in X_6\)，向量 \(u_w\) 至少有一个坐标为零：长根方向只占据两个坐标，而短根与 boundary 方向只占据一个坐标。故
\[
P(u_w)=0
\qquad(\forall\,w\in X_6),
\]
从而任意支撑于 \(\widetilde U_6\) 的原子测度 \(\mu\) 都满足
\[
\sum_{w\in X_6}c_w P(u_w)=0.
\]
另一方面，
\[
\int_{S^2}x_1^2x_2^2x_3^2\,\dd\sigma=\frac{1}{105}>0.
\]
因此若 \(\mu(S^2)\ne 0\)，便不可能有
\[
\sum_{w\in X_6}c_w P(u_w)=\mu(S^2)\int_{S^2}P(u)\,\dd\sigma(u).
\]
故六次精确 cubature 不存在。证毕。
\end{proof}

\begin{conjecture}[boundary 复制方向的低阶不可见性与首个统一谐和缺陷]\label{conj:window-even-boundary-duplicate-spherical-cubature}
设 \(m\) 为偶数，并假定该窗口存在与 window-\(6\) 同型的可见分裂
\[
X_m=X_m^{\mathrm{cyc}}\sqcup X_m^{\mathrm{bdry}},
\]
其中 \(X_m^{\mathrm{cyc}}\) 经某个审计字典嵌入约化根系
\[
\Phi_m\subset \mathfrak h_m^\ast,
\]
而 \(X_m^{\mathrm{bdry}}\) 恰给出若干与正短根方向重复的 boundary copies。
则应当存在以下普适现象：
\begin{enumerate}
  \item 低阶 \(SO(r_m)\)-moment exactness 只依赖于每条重复方向上的总质量，而不依赖于该总质量在 cyclic copy 与 boundary copy 之间的分配；
  \item 正权精确 cubature 的归一化模空间是一个 \(r_m\)-维立方体；
  \item 第一个超出精确阶的缺陷由一个与立方体参数无关的 \(W(\Phi_m)\)-谐和多项式统一控制。
\end{enumerate}
window-\(6\) 正是该原则的首个经审计样本：此时 \(r_6=3\)，正权模空间为
\[
{\Bigl[0,\frac{1}{30}\Bigr]}^{3},
\]
而首个缺陷由定理 \ref{thm:window6-labeled-spherical-degree6-harmonic-defect} 中的 \(\mathscr H_6\) 唯一控制。
\end{conjecture}

\paragraph{几何解释}
定理 \ref{thm:window6-labeled-spherical-degree5-cubature-family}--\ref{thm:window6-labeled-spherical-degree6-harmonic-defect} 表明：在一切不超过五阶的 \(SO(3)\) 角分辨可观测量中，boundary 正轴质量向 cyclic 正轴质量的转移完全不可见；真正开始分辨两者差异的，不是一阶偶极，也不是四阶各向异性，而是同一个六次谐和 defect \(\mathscr H_6\)。因此，window-\(6\) 的 boundary 扇区并不是低阶运动学中的显性扰动，而是高阶角动量通道中的首个可测破缺源。

\paragraph{根系支撑的精确消失理想、Hilbert twin 现象与四次插值饱和}
在 \(B_3/C_3\) 根字典、三次 Euclidean cubature 与唯一四次调和检测器都已固定之后，还可把 window-\(6\) 的 \(19\) 点归一化支撑进一步提升为一个完全显式的零维代数几何对象：其消失理想、标准单项式基、Hilbert 系列以及四次以内的全部插值关系都可被闭式写出。

\begin{theorem}[\texorpdfstring{$B_3$}{B3} 归一化支撑的精确消失理想与 Hilbert 系列]\label{thm:window6-b3-support-vanishing-ideal-hilbert}
沿用定理 \ref{thm:window6-21-adjoint-weight-multiset} 中的 \(B_3\) 规范，并记
\[
\mathscr X_B:=\{0\}\cup \Phi(B_3)\subset \mathbb A^3.
\]
则其消失理想恰为
\[
I_B=
\bigl\langle
x_1^3-x_1,\ x_2^3-x_2,\ x_3^3-x_3,\ x_1x_2x_3
\bigr\rangle
\subset \mathbb Q[x_1,x_2,x_3].
\]
并且在词典序 \(x_1>x_2>x_3\) 下，标准单项式恰为
\[
\mathcal B_B=
\bigl\{
x_1^a x_2^b x_3^c:\ 0\le a,b,c\le 2,\ \min(a,b,c)=0
\bigr\}.
\]
从而
\[
\dim_{\mathbb Q}\mathbb Q[x_1,x_2,x_3]/I_B=19,
\]
且其 Hilbert 系列为
\[
\boxed{\;
H_B(t)=1+3t+6t^2+6t^3+3t^4.\;}
\]
\leanpartial{paper\_window6\_b3\_support\_count}{仅形式化 19-点计数代数核 $|\{x \in \{-1,0,1\}^3 : x_1 x_2 x_3 = 0\}| = 27 - 8 = 19$；消失理想 $I_B$、Hilbert 系列 $H_B(t)$ 与 $B_3$ 根系结构留后轮}
\leanverified{cube27\_card}
\leanverified{supportNonzero\_card}
\leanverified{supportZero\_disjoint\_nonzero}
\leanverified{supportZero\_union\_nonzero}
\leanverified{supportZero\_card}
\end{theorem}
\begin{proof}
由方程
\[
x_i^3-x_i=0
\qquad (i=1,2,3)
\]
可知任一零点都满足
\[
x_i\in\{-1,0,1\}.
\]
再加上约束
\[
x_1x_2x_3=0,
\]
便得到
\[
V(I_B)
=
\bigl\{
(x_1,x_2,x_3)\in\{-1,0,1\}^3:\ x_1x_2x_3=0
\bigr\}.
\]
其基数为
\[
3^3-2^3=19.
\]
另一方面，\(\Phi(B_3)\) 在标准模型下恰为
\[
\{\pm e_1,\pm e_2,\pm e_3\}
\sqcup
\{\pm e_i\pm e_j:\ 1\le i<j\le 3\},
\]
亦即所有坐标属于 \(\{-1,0,1\}\)、且至少一个坐标为 \(0\) 的非零点。因此
\[
\mathscr X_B=\{0\}\cup \Phi(B_3)=V(I_B).
\]

现取词典序 \(x_1>x_2>x_3\)。四个生成元的首项分别为
\[
x_1^3,\qquad x_2^3,\qquad x_3^3,\qquad x_1x_2x_3.
\]
由于前三个首项两两互素，而其余与 \(x_1x_2x_3\) 的最小公倍式对应的 \(S\)-多项式都可直接由同一生成系约化到 \(0\)，故该生成系本身即为一个 Gröbner 基。于是
\[
\mathrm{in}(I_B)=\langle x_1^3,x_2^3,x_3^3,x_1x_2x_3\rangle.
\]
因此标准单项式正是所有指数都不超过 \(2\) 且不同时含三变量的单项式，即
\[
\mathcal B_B=
\bigl\{
x_1^a x_2^b x_3^c:\ 0\le a,b,c\le 2,\ \min(a,b,c)=0
\bigr\}.
\]
其个数为
\[
3^3-2^3=19.
\]
故
\[
\dim_{\mathbb Q}\mathbb Q[x_1,x_2,x_3]/I_B=19.
\]

又因 \(I_B\subseteq I(\mathscr X_B)\)，故存在满射
\[
\mathbb Q[x_1,x_2,x_3]/I_B\twoheadrightarrow \mathbb Q[x_1,x_2,x_3]/I(\mathscr X_B),
\]
右端维数等于 \(|\mathscr X_B|=19\)。左端维数亦为 \(19\)，故该满射只能是同构，从而
\[
I_B=I(\mathscr X_B).
\]

最后按总次数计数标准单项式即可得到
\[
H_B(t)=1+3t+6t^2+6t^3+3t^4.
\]
证毕。
\end{proof}
\leanverified{paper\_window6\_b3\_support\_vanishing\_ideal\_hilbert}

\begin{theorem}[\texorpdfstring{$C_3$}{C3} 归一化支撑的精确消失理想与 Hilbert 系列]\label{thm:window6-c3-support-vanishing-ideal-hilbert}
沿用定理 \ref{thm:window6-21-adjoint-weight-multiset} 中的 \(C_3\) 规范，并记
\[
\mathscr X_C:=\{0\}\cup \Phi(C_3)\subset \mathbb A^3,
\qquad
R:=x_1^2+x_2^2+x_3^2.
\]
则 \(\mathscr X_C\) 的消失理想可写为
\[
I_C=
\Bigl\langle
x_1(x_1^2-1)(x_1^2-4),\,
x_2(x_2^2-1)(x_2^2-4),\,
x_3(x_3^2-1)(x_3^2-4),\,
x_1x_2x_3,\,
R(R-2)(R-4)
\Bigr\rangle.
\]
并且在词典序 \(x_1>x_2>x_3\) 下，\(\mathbb Q[x_1,x_2,x_3]/I_C\) 的标准单项式恰为
\[
\mathcal B_C=
\{
1,\ x_1,\ x_2,\ x_3,\ x_1^2,\ x_1x_2,\ x_1x_3,\ x_2^2,\ x_2x_3,\ x_3^2,
\]
\[
x_1x_2^2,\ x_1x_3^2,\ x_2^3,\ x_2^2x_3,\ x_2x_3^2,\ x_3^3,\ x_2^4,\ x_2^2x_3^2,\ x_3^4
\}.
\]
从而
\[
\dim_{\mathbb Q}\mathbb Q[x_1,x_2,x_3]/I_C=19,
\]
且其 Hilbert 系列同样为
\[
\boxed{\;
H_C(t)=1+3t+6t^2+6t^3+3t^4.\;}
\]
\leanverified{paper\_window6\_c3\_support\_count}
\end{theorem}
\begin{proof}
前三类生成元逐坐标强制
\[
x_i\in\{0,\pm 1,\pm 2\},
\qquad i=1,2,3.
\]
再由
\[
x_1x_2x_3=0
\]
知至多有两个坐标非零；最后
\[
R(R-2)(R-4)=0
\]
强制平方范数满足
\[
R=x_1^2+x_2^2+x_3^2\in\{0,2,4\}.
\]
因此零点只能分成三类：
\[
(0,0,0),
\]
\[
(\pm 2,0,0),\ (0,\pm 2,0),\ (0,0,\pm 2),
\]
以及
\[
(\pm 1,\pm 1,0),\ (\pm 1,0,\pm 1),\ (0,\pm 1,\pm 1).
\]
它们的总数为
\[
1+6+12=19.
\]
而 \(\Phi(C_3)\) 的标准模型恰为
\[
\{\pm 2e_1,\pm 2e_2,\pm 2e_3\}
\sqcup
\{\pm e_i\pm e_j:\ 1\le i<j\le 3\},
\]
故
\[
V(I_C)=\mathscr X_C.
\]

对该生成系作 Buchberger 消元，在词典序 \(x_1>x_2>x_3\) 下得到的初始单项式理想由
\[
x_1^3,\ x_1^2x_2,\ x_1^2x_3,\ x_1x_2^3,\ x_1x_2x_3,\ x_1x_3^3,\ x_2^5,\ x_2^3x_3,\ x_2x_3^3,\ x_3^5
\]
生成。于是标准单项式恰为未被上述单项式整除的单项式，也就是题述的 \(19\) 个单项式 \(\mathcal B_C\)。因此
\[
\dim_{\mathbb Q}\mathbb Q[x_1,x_2,x_3]/I_C=19.
\]

同样地，由包含 \(I_C\subseteq I(\mathscr X_C)\) 诱导的满射
\[
\mathbb Q[x_1,x_2,x_3]/I_C\twoheadrightarrow \mathbb Q[x_1,x_2,x_3]/I(\mathscr X_C)
\]
以及两端同为 \(19\) 维，可知该满射为同构，从而
\[
I_C=I(\mathscr X_C).
\]
最后按总次数统计 \(\mathcal B_C\) 中单项式，即得
\[
H_C(t)=1+3t+6t^2+6t^3+3t^4.
\]
证毕。
\end{proof}
\leanverified{paper\_window6\_c3\_support\_vanishing\_ideal\_hilbert}

\begin{corollary}[\texorpdfstring{$B_3/C_3$}{B3/C3} 两支撑的 Hilbert twin 现象与四次唯一插值]\label{cor:window6-b3c3-hilbert-twins-degree4-interpolation}
沿用定理 \ref{thm:window6-b3-support-vanishing-ideal-hilbert} 与定理 \ref{thm:window6-c3-support-vanishing-ideal-hilbert} 的记号。则尽管 \(\mathscr X_B\) 与 \(\mathscr X_C\) 在几何上并不等距、其半径分层也不同，二者却满足完全相同的 Hilbert 系列
\[
H_B(t)=H_C(t)=1+3t+6t^2+6t^3+3t^4.
\]
因此它们在标准分次意义下构成一对 Hilbert twins。进一步：
\begin{enumerate}
  \item 对任意函数 \(f:\mathscr X_B\to \mathbb Q\)，存在且仅存在一个多项式
  \[
  P_B\in \mathrm{span}(\mathcal B_B)
  \]
  使得
  \[
  P_B|_{\mathscr X_B}=f;
  \]
  \item 对任意函数 \(g:\mathscr X_C\to \mathbb Q\)，存在且仅存在一个多项式
  \[
  P_C\in \mathrm{span}(\mathcal B_C)
  \]
  使得
  \[
  P_C|_{\mathscr X_C}=g.
  \]
\end{enumerate}
换言之，二者的插值复杂度严格相同，并且都在总次数 \(4\) 处闭合。
\leanverified{paper\_window6\_b3c3\_hilbert\_twins\_degree4\_interpolation}
\end{corollary}
\begin{proof}
第一条等式只是前两条定理的直接并列。

又由定理 \ref{thm:window6-b3-support-vanishing-ideal-hilbert}，
\(\mathcal B_B\) 是坐标环
\[
\mathbb Q[x_1,x_2,x_3]/I_B
\]
的一组标准单项式基，且其中每个单项式的总次数都不超过 \(4\)。因此
\[
\dim \mathrm{span}(\mathcal B_B)=19=|\mathscr X_B|.
\]
由于 \(I_B=I(\mathscr X_B)\)，评价映射
\[
\mathrm{ev}_B:\mathrm{span}(\mathcal B_B)\longrightarrow \mathrm{Fun}(\mathscr X_B,\mathbb Q)
\]
正是坐标环到函数代数的同构在这组基上的实现，故为同构，从而唯一插值成立。

对 \(\mathscr X_C\) 与 \(\mathcal B_C\) 完全同理。证毕。
\end{proof}

\begin{theorem}[\texorpdfstring{$B_3/C_3$}{B3/C3} 支撑的四次正交补空间维数为 \(16\)，且由消失理想完全控制]\label{thm:window6-b3c3-degree4-relation-space-saturation}
记
\[
\mathcal P_{\le 4}:=
\{P\in \mathbb Q[x_1,x_2,x_3]:\deg P\le 4\},
\qquad
\dim \mathcal P_{\le 4}=\binom{4+3}{3}=35.
\]
对 \(S\in\{B,C\}\)，定义评价映射
\[
\mathrm{ev}_S:\mathcal P_{\le 4}\to \mathrm{Fun}(\mathscr X_S,\mathbb Q).
\]
则有
\[
\mathrm{rank}(\mathrm{ev}_S)=19,
\qquad
\dim\ker(\mathrm{ev}_S)=16.
\]
并且
\[
\ker(\mathrm{ev}_B)=I_B\cap \mathcal P_{\le 4},
\qquad
\ker(\mathrm{ev}_C)=I_C\cap \mathcal P_{\le 4}.
\]
特别地，\(\mathscr X_B\) 与 \(\mathscr X_C\) 在四次以内的全部代数关系都已饱和：不存在任何额外的低次数关系超出各自消失理想所给出的那些。
\leanverified{paper\_window6\_b3c3\_degree4\_relation\_space\_saturation}
\end{theorem}
\begin{proof}
由推论 \ref{cor:window6-b3c3-hilbert-twins-degree4-interpolation}，对 \(S=B,C\) 都可在 \(\mathcal P_{\le 4}\) 中找到一个 \(19\) 维子空间，其评价映射到 \(\mathrm{Fun}(\mathscr X_S,\mathbb Q)\) 为同构。因此
\[
\mathrm{rank}(\mathrm{ev}_S)=19.
\]
再由秩--零度定理，
\[
\dim\ker(\mathrm{ev}_S)=35-19=16.
\]

另一方面，\(\ker(\mathrm{ev}_S)\) 按定义就是所有在 \(\mathscr X_S\) 上消失的次数不超过 \(4\) 的多项式，故恰为
\[
I(\mathscr X_S)\cap \mathcal P_{\le 4}.
\]
由定理 \ref{thm:window6-b3-support-vanishing-ideal-hilbert} 与定理 \ref{thm:window6-c3-support-vanishing-ideal-hilbert}，
\[
I(\mathscr X_B)=I_B,
\qquad
I(\mathscr X_C)=I_C,
\]
于是
\[
\ker(\mathrm{ev}_B)=I_B\cap \mathcal P_{\le 4},
\qquad
\ker(\mathrm{ev}_C)=I_C\cap \mathcal P_{\le 4}.
\]
这正说明四次以内的一切关系都已经由各自的完整消失理想给出，而不存在额外独立的低次数约束。证毕。
\end{proof}

\noindent
由此，window-\(6\) 的 \(B_3/C_3\) 根系字典不再只是一个离散标号装置，而是同时诱导出两套 \(19\) 点零维支撑的完整代数几何模型：一方面，二者分别具有显式消失理想与标准单项式基；另一方面，尽管其几何半径结构不同，却共享同一 Hilbert 系列、同一四次插值阈值以及同一 \(16\) 维低次关系空间规模。于是，\(B_3\) 与 \(C_3\) 两种归一化不只在根长交换意义下对偶，而且在有限支撑的分次代数复杂度上也形成一对严格的 Hilbert twins。

\endinput

\paragraph{三张 Levi 平面、三轴选择律与唯一 quartic 残差}
在 \(B_3/C_3\) 可见权字典、三次 cubature、四次调和检测器与支撑理想都已固定之后，window-\(6\) 的可见部门还可进一步压缩为一条更接近有效场论的局域结构链：全部非零支撑并不充满三维 Cartan 空间，而是精确贴在三张 rank-\(2\) Levi 平面上；任何真正的三轴单点耦合都逐点湮灭；二阶之后的首个非标量局域残差只沿唯一一条八面体型四次调和方向存活；而 \(C_3\) 规范相对于 \(B_3\) 规范的全部增厚则可在全偶矩层级上写成闭式。

\begin{theorem}[window-\(6\) 可见权的三 Levi 平面分解]\label{thm:window6-b3c3-visible-support-three-levi-planes}
沿用定理 \ref{thm:window6-21-adjoint-weight-multiset} 中的权映射
\[
\Lambda_B,\Lambda_C:X_6\to \ZZ^3
\]
及 boundary 三词映到零权的约定。
对 \(1\le i<j\le 3\)，记
\[
H_{ij}:=\operatorname{span}\{e_i,e_j\}\subset \RR^3.
\]
则有精确并分解
\[
\Lambda_B(X_6)\setminus\{0\}
=
\Phi^{B_2}_{12}\cup \Phi^{B_2}_{13}\cup \Phi^{B_2}_{23},
\qquad
\Phi^{B_2}_{ij}:=\{\pm e_i,\pm e_j,\pm e_i\pm e_j\}\subset H_{ij},
\]
以及
\[
\Lambda_C(X_6)\setminus\{0\}
=
\Phi^{C_2}_{12}\cup \Phi^{C_2}_{13}\cup \Phi^{C_2}_{23},
\qquad
\Phi^{C_2}_{ij}:=\{\pm 2e_i,\pm 2e_j,\pm e_i\pm e_j\}\subset H_{ij}.
\]
特别地，对任意可见权
\[
\lambda\in \Lambda_B(X_6)\cup \Lambda_C(X_6)
\]
都满足
\[
\lambda_1\lambda_2\lambda_3=0.
\]
\end{theorem}
\begin{proof}
由定理 \ref{thm:window6-21-adjoint-weight-multiset}，\(B_3\) 规范下的 \(18\) 条非零权恰为
\[
\{\pm e_1,\pm e_2,\pm e_3\}
\sqcup
\{\pm e_i\pm e_j:\ 1\le i<j\le 3\},
\]
而 \(C_3\) 规范下的 \(18\) 条非零权恰为
\[
\{\pm 2e_1,\pm 2e_2,\pm 2e_3\}
\sqcup
\{\pm e_i\pm e_j:\ 1\le i<j\le 3\}.
\]
按零坐标的位置分组，在 \(B_3\) 规范下每个 \(H_{ij}\) 内恰出现
\[
\{\pm e_i,\pm e_j,\pm e_i\pm e_j\},
\]
即一个 \(B_2\) 根系；在 \(C_3\) 规范下每个 \(H_{ij}\) 内恰出现
\[
\{\pm 2e_i,\pm 2e_j,\pm e_i\pm e_j\},
\]
即一个 \(C_2\) 根系。
而 boundary 三词都被送到零向量，不影响非零支撑。
故所示并分解成立。

每个非零支撑点都落在某个坐标平面 \(H_{ij}\) 上，故至少有一个坐标为零，从而
\(
\lambda_1\lambda_2\lambda_3=0
\)
对所有可见权成立。证毕。
\end{proof}

\begin{proposition}[三轴选择律与指数特征函数的三阶湮灭]\label{prop:window6-b3c3-triaxial-selection-ideal-factorization}
定义可见矩泛函
\[
\mathcal M_B(P):=\sum_{w\in X_6}P(\Lambda_B(w)),
\qquad
\mathcal M_C(P):=\sum_{w\in X_6}P(\Lambda_C(w)),
\qquad
P\in \CC[x_1,x_2,x_3].
\]
则 \(\mathcal M_B\) 与 \(\mathcal M_C\) 都经由商代数
\[
\CC[x_1,x_2,x_3]/(x_1x_2x_3)
\]
因子化。
等价地，只要 \(a,b,c\ge 1\)，便有
\[
\sum_{w\in X_6}\Lambda_B(w)_1^a\Lambda_B(w)_2^b\Lambda_B(w)_3^c=0,
\qquad
\sum_{w\in X_6}\Lambda_C(w)_1^a\Lambda_C(w)_2^b\Lambda_C(w)_3^c=0.
\]
进一步，沿用定理 \ref{thm:window6-b3c3-adjoint-character-and-quartic-invariant} 的指数特征函数
\[
\Theta_B(t):=\sum_{w\in X_6}e^{\langle \Lambda_B(w),t\rangle},
\qquad
\Theta_C(t):=\sum_{w\in X_6}e^{\langle \Lambda_C(w),t\rangle},
\qquad
t=(t_1,t_2,t_3)\in \RR^3,
\]
则恒有
\[
\partial_{t_1}\partial_{t_2}\partial_{t_3}\Theta_B\equiv 0,
\qquad
\partial_{t_1}\partial_{t_2}\partial_{t_3}\Theta_C\equiv 0.
\]
\end{proposition}
\begin{proof}
由定理 \ref{thm:window6-b3c3-visible-support-three-levi-planes}，对每个可见权 \(\lambda\) 都有
\[
\lambda_1\lambda_2\lambda_3=0.
\]
因此任意落在理想 \((x_1x_2x_3)\) 中的多项式在全部支撑点上逐点为零，故 \(\mathcal M_B,\mathcal M_C\) 都通过
\(
\CC[x_1,x_2,x_3]/(x_1x_2x_3)
\)
因子化。
把该结论应用到单项式 \(x_1^a x_2^b x_3^c\) 即得幂矩消失式。

对指数特征函数，逐项求导可得
\[
\partial_{t_1}\partial_{t_2}\partial_{t_3}
e^{\langle \lambda,t\rangle}
=
\lambda_1\lambda_2\lambda_3\,e^{\langle \lambda,t\rangle}
=
0.
\]
对全部 \(\lambda=\Lambda_B(w)\) 或 \(\Lambda_C(w)\) 求和，即得所示恒等式。证毕。
\end{proof}

\begin{theorem}[quartic 缺陷的一维性]\label{thm:window6-b3c3-quartic-defect-onedim}
定义
\[
s_2(u):=u_1^2+u_2^2+u_3^2,
\qquad
s_4(u):=u_1^4+u_2^4+u_3^4,
\]
并记归一化八面体型调和四次式
\[
H_4^{\mathrm{oct}}(u):=s_4(u)-\frac35 s_2(u)^2.
\]
则对任意 \(u\in \RR^3\)，有精确恒等式
\[
\sum_{w\in X_6}\langle \Lambda_B(w),u\rangle^2=10\,s_2(u),
\qquad
\sum_{w\in X_6}\langle \Lambda_C(w),u\rangle^2=16\,s_2(u),
\]
以及
\[
\sum_{w\in X_6}\langle \Lambda_B(w),u\rangle^4
=
\frac{54}{5}s_2(u)^2-2\,H_4^{\mathrm{oct}}(u),
\]
\[
\sum_{w\in X_6}\langle \Lambda_C(w),u\rangle^4
=
\frac{144}{5}s_2(u)^2+28\,H_4^{\mathrm{oct}}(u).
\]
因此，window-\(6\) 可见部门在二阶上完全各向同性，而其四阶偏离只沿唯一的一维调和方向 \(\RR H_4^{\mathrm{oct}}\) 存活。
\end{theorem}
\begin{proof}
前两条二阶恒等式正是定理 \ref{thm:window6-b3c3-adjoint-second-moment-isotropy} 的内容。
又由定理 \ref{thm:window6-b3c3-adjoint-character-and-quartic-invariant}，
\[
\sum_{w\in X_6}\langle \Lambda_B(w),u\rangle^4
=
10s_4(u)+24\sum_{1\le i<j\le 3}u_i^2u_j^2,
\]
\[
\sum_{w\in X_6}\langle \Lambda_C(w),u\rangle^4
=
40s_4(u)+24\sum_{1\le i<j\le 3}u_i^2u_j^2.
\]
由
\[
s_2(u)^2=s_4(u)+2\sum_{1\le i<j\le 3}u_i^2u_j^2
\]
可化为
\[
\sum_{w\in X_6}\langle \Lambda_B(w),u\rangle^4
=
12s_2(u)^2-2s_4(u),
\]
\[
\sum_{w\in X_6}\langle \Lambda_C(w),u\rangle^4
=
12s_2(u)^2+28s_4(u).
\]
再用
\[
s_4(u)=H_4^{\mathrm{oct}}(u)+\frac35 s_2(u)^2
\]
代回，即得
\[
12s_2^2-2s_4
=
\frac{54}{5}s_2^2-2H_4^{\mathrm{oct}},
\]
\[
12s_2^2+28s_4
=
\frac{144}{5}s_2^2+28H_4^{\mathrm{oct}}.
\]
故四阶非径向偏离只沿 \(H_4^{\mathrm{oct}}\) 这一条调和方向出现。证毕。
\end{proof}

\begin{theorem}[quartic 极值方向的 axis--diagonal 反转]\label{thm:window6-b3c3-quartic-axis-diagonal-reversal}
定义
\[
M_{4,B}(u):=\sum_{w\in X_6}\langle \Lambda_B(w),u\rangle^4,
\qquad
M_{4,C}(u):=\sum_{w\in X_6}\langle \Lambda_C(w),u\rangle^4.
\]
若 \(|u|=1\)，则
\[
M_{4,B}(u)=12-2\sum_{i=1}^{3}u_i^4,
\qquad
M_{4,C}(u)=12+28\sum_{i=1}^{3}u_i^4.
\]
因此
\[
\min_{|u|=1}M_{4,B}(u)=10,
\qquad
\max_{|u|=1}M_{4,B}(u)=\frac{34}{3},
\]
\[
\min_{|u|=1}M_{4,C}(u)=\frac{64}{3},
\qquad
\max_{|u|=1}M_{4,C}(u)=40.
\]
并且极值方向精确为：
\begin{enumerate}
  \item \(B_3\) 规范下，最小值在坐标轴方向
  \[
  \{\pm e_1,\pm e_2,\pm e_3\}
  \]
  处取得，而最大值在体对角方向
  \[
  \frac{1}{\sqrt3}(\pm 1,\pm 1,\pm 1)
  \]
  处取得；
  \item \(C_3\) 规范下，上述极值方向完全反转。
\end{enumerate}
\end{theorem}
\begin{proof}
由定理 \ref{thm:window6-b3c3-quartic-defect-onedim}，在单位球面 \(|u|=1\) 上有
\[
H_4^{\mathrm{oct}}(u)=\sum_{i=1}^{3}u_i^4-\frac35.
\]
代回即可得到
\[
M_{4,B}(u)=12-2\sum_{i=1}^{3}u_i^4,
\qquad
M_{4,C}(u)=12+28\sum_{i=1}^{3}u_i^4.
\]

令
\[
\Sigma_4(u):=\sum_{i=1}^{3}u_i^4.
\]
在约束 \(|u|=1\) 下，由 Hölder 不等式或 Lagrange 乘子法，
\[
\frac13\le \Sigma_4(u)\le 1,
\]
其中 \(\Sigma_4(u)=1\) 当且仅当 \(u\) 为坐标轴方向，而 \(\Sigma_4(u)=1/3\) 当且仅当
\[
|u_1|=|u_2|=|u_3|=\frac{1}{\sqrt3}.
\]
把这两个极值代回上述线性公式，即得全部极值与极值方向。证毕。
\end{proof}

\begin{theorem}[B/C 规范差的全偶矩层级与 Laplace 支配]\label{thm:window6-b3c3-even-moment-gap-laplace-domination}
对任意整数 \(m\ge 1\) 与任意 \(u\in \RR^3\)，记
\[
M_{2m}^{B}(u):=\sum_{w\in X_6}\langle \Lambda_B(w),u\rangle^{2m},
\qquad
M_{2m}^{C}(u):=\sum_{w\in X_6}\langle \Lambda_C(w),u\rangle^{2m}.
\]
则有精确恒等式
\[
M_{2m}^{C}(u)-M_{2m}^{B}(u)
=
2\bigl(2^{2m}-1\bigr)\sum_{i=1}^{3}u_i^{2m}.
\]
等价地，指数特征函数满足
\[
\Theta_C(t)-\Theta_B(t)
=
2\sum_{i=1}^{3}\bigl(\cosh(2t_i)-\cosh(t_i)\bigr)\ge 0
\qquad
(t\in \RR^3),
\]
且当且仅当 \(t=0\) 时取等。
因此，对每个固定方向 \(u\)，标量随机变量 \(\langle \Lambda_C,u\rangle\) 在 Laplace 序上严格支配 \(\langle \Lambda_B,u\rangle\)。
\end{theorem}
\begin{proof}
由定理 \ref{thm:window6-21-adjoint-weight-multiset}，\(\Lambda_B\) 与 \(\Lambda_C\) 的差别只发生在六个轴向词上：\(B_3\) 规范给出
\[
\{\pm e_1,\pm e_2,\pm e_3\},
\]
而 \(C_3\) 规范把它们改写为
\[
\{\pm 2e_1,\pm 2e_2,\pm 2e_3\}.
\]
其余十二条双坐标根和三重零权完全相同。
因此
\[
M_{2m}^{C}(u)-M_{2m}^{B}(u)
=
\sum_{i=1}^{3}\Bigl((2u_i)^{2m}+(-2u_i)^{2m}-u_i^{2m}-(-u_i)^{2m}\Bigr),
\]
即得
\[
M_{2m}^{C}(u)-M_{2m}^{B}(u)
=
2\bigl(2^{2m}-1\bigr)\sum_{i=1}^{3}u_i^{2m}.
\]

对该恒等式按 \(m\) 求指数母函数，或直接利用定理 \ref{thm:window6-b3c3-adjoint-character-and-quartic-invariant} 中 \(\Theta_B,\Theta_C\) 的显式公式，便得到
\[
\Theta_C(t)-\Theta_B(t)
=
2\sum_{i=1}^{3}\bigl(\cosh(2t_i)-\cosh(t_i)\bigr).
\]
因为 \(\cosh(2x)>\cosh(x)\) 对每个 \(x\neq 0\) 都成立，故右端非负，且只有在 \(t_1=t_2=t_3=0\) 时为零。
于是 \(\Theta_C(t)\ge \Theta_B(t)\) 对所有 \(t\in \RR^3\) 成立，并在 \(t\neq 0\) 时严格大于，故 \(C_3\) 规范在每个方向上的 Laplace 变换都严格支配 \(B_3\) 规范。证毕。
\end{proof}

\begin{conjecture}[GFF 标度下可见部门的唯一 quartic 残差原理]\label{conj:window6-visible-gff-unique-quartic-residual}
设某个临界六顶点系综中，window-\(6\) 的可见线性观测在适当归一化后收敛到标量 Gaussian free field 的线性像。
则单点可见有效作用的局域展开应同时满足：
\begin{enumerate}
  \item 不存在真正的三轴 cubic 项；
  \item 二阶项必为标量刚度；
  \item 首个非高斯、非标量且可见的局域修正必与
  \[
  H_4^{\mathrm{oct}}(u)=u_1^4+u_2^4+u_3^4-\frac35|u|^4
  \]
  成比例；
  \item \(B_3\) 与 \(C_3\) 两种规范只改变这一 quartic 项的系数，不改变其形式。
\end{enumerate}
\end{conjecture}

\begin{remark}[唯一 quartic 残差原则的离散依据]\label{rem:window6-visible-gff-unique-quartic-residual-evidence}
命题 \ref{prop:window6-b3c3-triaxial-selection-ideal-factorization} 排除了全部真正的三轴单点耦合；定理 \ref{thm:window6-b3c3-quartic-defect-onedim} 表明二阶完全各向同性，而四阶缺陷空间仅为一维；定理 \ref{thm:window6-b3c3-even-moment-gap-laplace-domination} 又说明 \(B/C\) 规范差在所有偶矩层级上都只能通过
\[
\sum_{i=1}^{3}u_i^{2m}
\]
这一轴向纯幂族进入。
因此，若 window-\(6\) 的可见部门在连续极限中确有首个非高斯残差，则它既不应首先出现在 cubic 通道，也不应分裂成多个独立的张量方向；最先显现的可见局域异常只能是唯一的八面体型 quartic operator。
\end{remark}

\endinput


\begin{theorem}[阿贝尔化奇偶荷格的根--Cartan 三分裂与 boundary parity 的零权本质]\label{thm:window6-abelianized-parity-charge-root-cartan-splitting}
记
\[
\mathcal G_6^{\mathrm{bin}}:=\Aut(\Fold_6^{\mathrm{bin}}),
\]
并经由纤维符号映射 \(\mathrm{sgn}_6\) 的坐标识别写成
\[
\bigl(\mathcal G_6^{\mathrm{bin}}\bigr)^{\mathrm{ab}}\cong (\ZZ_2)^{X_6}.
\]
在 \(B_3\) 规范下令
\[
X_6^{\mathrm{short}}:=\{w\in X_6^{\mathrm{cyc}}:\ \mathrm{wt}(w)=1\},
\qquad
X_6^{\mathrm{long}}:=X_6^{\mathrm{cyc}}\setminus X_6^{\mathrm{short}}.
\]
并记相应的根集合为
\[
\Phi_{\mathrm{short}}:=\{\pm e_1,\pm e_2,\pm e_3\},
\qquad
\Phi_{\mathrm{long}}:=\{\pm e_i\pm e_j:\ 1\le i<j\le 3\}.
\]
则存在规范直和分解
\[
\bigl(\mathcal G_6^{\mathrm{bin}}\bigr)^{\mathrm{ab}}
\cong
(\ZZ_2)^{X_6^{\mathrm{short}}}
\oplus
(\ZZ_2)^{X_6^{\mathrm{long}}}
\oplus
(\ZZ_2)^{X_6^{\mathrm{bdry}}}
\cong
\FF_2[\Phi_{\mathrm{short}}]
\oplus
\FF_2[\Phi_{\mathrm{long}}]
\oplus
\mathfrak h_{\FF_2},
\]
其中
\[
\dim_{\FF_2}\FF_2[\Phi_{\mathrm{short}}]=6,\qquad
\dim_{\FF_2}\FF_2[\Phi_{\mathrm{long}}]=12,\qquad
\dim_{\FF_2}\mathfrak h_{\FF_2}=3.
\]
若改用 \(C_3\) 规范，则前两项仅交换长短根名称，而三项维数保持不变。

进一步，记
\[
Z_{\partial}:=\langle \tau_w:\ w\in X_6^{\mathrm{bdry}}\rangle\cong (\ZZ_2)^3.
\]
则 Abel 化映射把 \(Z_{\partial}\) 同构到第三项 \(\mathfrak h_{\FF_2}\) 上。
因此 boundary sheet parity 在该伴随权模型中对应的正是零权的 Cartan 部分，而不是根方向电荷。
\end{theorem}
\begin{proof}
由定理 \ref{thm:window6-cyclic-weight-threshold-root-length} 与定理 \ref{thm:window6-21-adjoint-weight-multiset}，集合 \(X_6\) 已被规范分解为
\[
X_6=X_6^{\mathrm{short}}\sqcup X_6^{\mathrm{long}}\sqcup X_6^{\mathrm{bdry}},
\qquad
|X_6^{\mathrm{short}}|=6,\quad |X_6^{\mathrm{long}}|=12,\quad |X_6^{\mathrm{bdry}}|=3.
\]
另一方面，定理 \ref{thm:window6-foldbin-gauge-abelianization-even-parity} 给出
\[
\bigl(\mathcal G_6^{\mathrm{bin}}\bigr)^{\mathrm{ab}}\cong (\ZZ_2)^{21},
\]
并且该 Abel 化通过 \(\mathrm{sgn}_6\) 被识别为以 \(X_6\) 为坐标集的奇偶荷格。
因此对上述三块分割逐块取坐标子格，便得到所示的规范直和分解。

再由推论 \ref{cor:fold-bin-boundary-center-z2}，boundary 三条二点纤维的纤维内换位 involution 生成一个规范中心子群
\(
Z_{\partial}\cong(\ZZ_2)^3
\)，其在 Abel 化中的像恰对应 boundary 三个坐标方向。
这正是第三项 \((\ZZ_2)^{X_6^{\mathrm{bdry}}}=\mathfrak h_{\FF_2}\)。
故 boundary parity 的本质是零权的 Cartan 部分，而非根空间部分。证毕。
\end{proof}

\begin{theorem}[window-\(6\) 的 \texorpdfstring{$C_6$}{C6} 轨道分解与 \texorpdfstring{$B_3$}{B3} 长短根分区的交会刚性]\label{thm:window6-c6-orbit-b3-root-intersection-rigidity}
沿用 \(B_3\) 规范下的分块
\[
X_6^{\mathrm{cyc}}=X_6^{\mathrm{short}}\sqcup X_6^{\mathrm{long}}.
\]
则表 \ref{tab:window6_c6_orbit_decomposition} 给出的五条 \(C_6\) 轨道中，恰有且仅有一条长度为 \(6\) 的轨道满足
\[
X_6^{\mathrm{short}}
=
\{\texttt{000001},\texttt{000010},\texttt{000100},\texttt{001000},\texttt{010000},\texttt{100000}\},
\]
其余四条轨道全部包含于 \(X_6^{\mathrm{long}}\)。
特别地，另一条长度为 \(6\) 的轨道
\[
\{\texttt{000101},\texttt{001010},\texttt{010001},\texttt{010100},\texttt{100010},\texttt{101000}\}
\]
整体落在长根扇区之内。

从而，同时满足“在每条 \(C_6\) 轨道上常值”与“Weyl-不变”的函数空间恰为
\[
\left\{
f:X_6^{\mathrm{cyc}}\to\CC:\ 
f\text{ 在每条 }C_6\text{ 轨道上常值且为 Weyl-不变}
\right\}
=
\operatorname{span}_{\CC}\{\mathbf 1_{X_6^{\mathrm{short}}},\mathbf 1_{X_6^{\mathrm{long}}}\},
\]
因而其维数恰为 \(2\)。
\end{theorem}

\begin{proof}
由表 \ref{tab:window6_c6_orbit_decomposition}，\(X_6^{\mathrm{cyc}}\) 的两条长度为 \(6\) 的轨道分别为
\[
\{\texttt{000001},\texttt{000010},\texttt{000100},\texttt{001000},\texttt{010000},\texttt{100000}\}
\]
与
\[
\{\texttt{000101},\texttt{001010},\texttt{010001},\texttt{010100},\texttt{100010},\texttt{101000}\}.
\]
其中第一条轨道恰由全部 \(\mathrm{wt}=1\) 的六个词组成，故由定理
\ref{thm:window6-cyclic-weight-threshold-root-length} 正好等于 \(X_6^{\mathrm{short}}\)。
第二条长度为 \(6\) 的轨道在表 \ref{tab:fold6_b3c3_root_dictionary_b3} 中逐项对应长根，因而整体落在 \(X_6^{\mathrm{long}}\)。

其余三条轨道分别由 \(\texttt{000000}\)、\(\{\texttt{010101},\texttt{101010}\}\) 与
\(\{\texttt{001001},\texttt{010010},\texttt{100100}\}\) 生成；
它们的 Hamming 权分别为 \(0\)、\(3\) 与 \(2\)，故仍由定理
\ref{thm:window6-cyclic-weight-threshold-root-length} 全部落在 \(X_6^{\mathrm{long}}\)。
因此五条 \(C_6\) 轨道与长短根二分块的交会类型被唯一确定为：
唯一一条长度为 \(6\) 的轨道等于短根扇区，其余四条轨道全为长根扇区的子块。

现在设 \(f:X_6^{\mathrm{cyc}}\to\CC\) 在每条 \(C_6\) 轨道上常值且 Weyl-不变。
由于 \(X_6^{\mathrm{short}}\) 本身就是一条单独的 \(C_6\) 轨道，\(f\) 在其上只能取一个常数。
另一方面，\(X_6^{\mathrm{long}}\) 虽由四条 \(C_6\) 轨道组成，但 Weyl-不变性强迫 \(f\) 在整块 \(X_6^{\mathrm{long}}\) 上取同一个常数。
故该函数空间恰由 \(\mathbf 1_{X_6^{\mathrm{short}}}\) 与 \(\mathbf 1_{X_6^{\mathrm{long}}}\) 张成，维数为 \(2\)。证毕。
\end{proof}

\paragraph{visible Cartan quotient、syzygy lattice 与二元码闭链}
在 \(B_3\) 可见根字典、boundary 三基、四次调和检测器与 visible Laplace 支配都已固定之后，window-\(6\) 的 \(21\) 个可见模态还可继续压缩为一条更刚性的整数线性链：
boundary 三词先给出一个规范的 rank-\(3\) Cartan quotient，其核形成一个秩 \(18\) 的 syzygy lattice；该 lattice 的 Gram 判别式与 visible matroid 的全部三阶体积预算同时锁定到 \(11^3\)；四阶层面恰有唯一的 boundary counterterm 能把 cyclic cloud 各向同性化，而六阶则出现不可消去的各向异性障碍；最后，这个模 \(2\) Cartan quotient 自然生成一对互补的 \([21,3,9]\) 与 \([21,18,2]\) 码。

\leanverified{paper\_window6\_visible\_cartan\_quotient\_syzygy\_splitting}
\begin{theorem}[visible Cartan quotient 的规范分裂与 syzygy 基]\label{thm:window6-visible-cartan-quotient-syzygy-splitting}
固定 boundary 排序
\[
b_1:=\texttt{100001},
\qquad
b_2:=\texttt{100101},
\qquad
b_3:=\texttt{101001},
\]
并记表 \ref{tab:fold6_b3c3_root_dictionary_b3} 的 \(B_3\) 根字典为
\[
\alpha:X_6^{\mathrm{cyc}}\longrightarrow \Phi(B_3)\subset \ZZ^3.
\]
定义可见向量
\[
v_w:=
\begin{cases}
\alpha(w),& w\in X_6^{\mathrm{cyc}},\\
e_i,& w=b_i\in X_6^{\mathrm{bdry}},
\end{cases}
\]
以及整数线性映射
\[
\Pi_6:\ZZ^{X_6}\longrightarrow \ZZ^3,
\qquad
\Pi_6(\delta_w):=v_w,
\]
其中 \(\{\delta_w\}_{w\in X_6}\) 为标准坐标基。
则存在规范分裂正合列
\[
0\longrightarrow \Lambda_6^{\mathrm{syz}}
\longrightarrow \ZZ^{X_6}
\xrightarrow{\ \Pi_6\ }
\ZZ^3
\longrightarrow 0,
\]
其分裂由
\[
s:\ZZ^3\to \ZZ^{X_6},
\qquad
s(e_i)=\delta_{b_i}
\]
给出。更精确地，
\[
\Lambda_6^{\mathrm{syz}}=\ker \Pi_6
\]
是秩 \(18\) 的自由格，并有显式基
\[
r_w:=\delta_w-\sum_{i=1}^{3}\alpha_i(w)\,\delta_{b_i},
\qquad
w\in X_6^{\mathrm{cyc}}.
\]
因此
\[
\ZZ^{X_6}\cong \ZZ^3\oplus \ZZ^{18},
\]
其中三维直和因子是 visible Cartan 层，十八维直和因子是 cyclic syzygy 层。
\end{theorem}
\begin{proof}
由于 \(v_{b_i}=e_i\)，映射 \(\Pi_6\) 显然满射，且
\[
\Pi_6\circ s=\mathrm{id}_{\ZZ^3},
\]
故该正合列分裂。

对每个 \(w\in X_6^{\mathrm{cyc}}\)，有
\[
\Pi_6(r_w)
=
\alpha(w)-\sum_{i=1}^{3}\alpha_i(w)e_i
=
0,
\]
故 \(r_w\in \ker \Pi_6\)。
这些向量线性无关，因为每个 \(r_w\) 在其专属坐标 \(\delta_w\) 上系数为 \(1\)，而其余 \(r_{w'}\) 在该坐标上系数为 \(0\)。

另一方面，
\[
\operatorname{rank}\ker \Pi_6
=
|X_6|-3
=
21-3
=
18
=
|X_6^{\mathrm{cyc}}|,
\]
故 \(\{r_w\}_{w\in X_6^{\mathrm{cyc}}}\) 即为 \(\ker \Pi_6\) 的一组基。

最后，把该构造模 \(2\) 约化，并结合定理 \ref{thm:window6-foldbin-gauge-abelianization-even-parity} 中
\[
\bigl(\mathcal G_6^{\mathrm{bin}}\bigr)^{\mathrm{ab}}\cong (\ZZ_2)^{21}
\]
及推论 \ref{cor:fold-bin-boundary-center-z2} 给出的 boundary 三维坐标子格，可得 anomaly lattice 上的同样分裂。证毕。
\end{proof}

\begin{theorem}[syzygy lattice 的 Gram 谱与判别式]\label{thm:window6-syzygy-gram-spectrum-discriminant}
\leanverified{paper\_window6\_syzygy\_gram\_spectrum\_discriminant}
令
\[
R_6:=\bigl[r_w\bigr]_{w\in X_6^{\mathrm{cyc}}}\in M_{21\times 18}(\ZZ)
\]
为定理 \ref{thm:window6-visible-cartan-quotient-syzygy-splitting} 的 syzygy 基矩阵，并记
\[
K_6:=R_6^{\mathsf T}R_6.
\]
则
\[
\operatorname{Spec}(K_6)=\{11^{(3)},1^{(15)}\},
\]
从而
\[
\det K_6=11^3,
\qquad
\operatorname{tr}K_6=48.
\]
特别地，
\[
\operatorname{disc}(\Lambda_6^{\mathrm{syz}})=11^3.
\]
因此 \(\Lambda_6^{\mathrm{syz}}\) 并非均匀刚性，而是精确裂解为 \(15\) 个单位刚度方向与 \(3\) 个刚度 \(11\) 的特异方向；后者的个数恰与 Cartan 秩相同。
\end{theorem}
\begin{proof}
记
\[
A_6^{\mathrm{cyc}}
:=
\bigl[\alpha(w)\bigr]_{w\in X_6^{\mathrm{cyc}}}
\in M_{3\times 18}(\ZZ).
\]
由
\[
r_w=\delta_w-\sum_{i=1}^{3}\alpha_i(w)\delta_{b_i}
\]
的定义直接可见
\[
K_6=I_{18}+\bigl(A_6^{\mathrm{cyc}}\bigr)^{\mathsf T}A_6^{\mathrm{cyc}}.
\]
另一方面，定理 \ref{thm:window6-b3c3-adjoint-second-moment-isotropy} 给出 cyclic 根云的二阶矩恒等式
\[
A_6^{\mathrm{cyc}}\bigl(A_6^{\mathrm{cyc}}\bigr)^{\mathsf T}
=
\sum_{w\in X_6^{\mathrm{cyc}}}\alpha(w)\alpha(w)^{\mathsf T}
=
10I_3.
\]
故 \(\bigl(A_6^{\mathrm{cyc}}\bigr)^{\mathsf T}A_6^{\mathrm{cyc}}\) 的非零特征值恰为 \(10\)（重数 \(3\)），其余 \(15\) 个特征值为 \(0\)。
对其加上 \(I_{18}\) 后便得到
\[
\operatorname{Spec}(K_6)=\{11^{(3)},1^{(15)}\}.
\]
于是
\[
\det K_6=11^3,
\qquad
\operatorname{tr}K_6=3\cdot 11+15\cdot 1=48.
\]
由于 \(\{r_w\}\) 是 \(\Lambda_6^{\mathrm{syz}}\) 的整基，\(\det K_6\) 正是该格的判别式。证毕。
\end{proof}

\begin{theorem}[quartic isotropization 的唯一 boundary counterterm]\label{thm:window6-visible-quartic-shell-isotropization}
\leanverified{paper\_window6\_visible\_quartic\_shell\_isotropization}
对 \(d\in\{2,3,4\}\)，定义 cyclic 壳层
\[
\Phi_{6,d}^{\mathrm{cyc}}
:=
\{\alpha(w):\ w\in X_6^{\mathrm{cyc}},\ d_6^{\mathrm{bin}}(w)=d\}.
\]
对参数 \((a_2,a_3,a_4,b)\in \CC^4\)，定义四次型
\[
\mathcal Q_{a_2,a_3,a_4,b}(x):=
a_2\sum_{\beta\in \Phi_{6,2}^{\mathrm{cyc}}}\langle x,\beta\rangle^4
+
a_3\sum_{\beta\in \Phi_{6,3}^{\mathrm{cyc}}}\langle x,\beta\rangle^4
+
a_4\sum_{\beta\in \Phi_{6,4}^{\mathrm{cyc}}}\langle x,\beta\rangle^4
+
b\sum_{i=1}^{3}x_i^4.
\]
则存在常数 \(c\) 使
\[
\mathcal Q_{a_2,a_3,a_4,b}(x)=c|x|^4
\]
当且仅当
\[
a_2=a_3=a_4=:a,
\qquad
b=2a,
\qquad
c=12a.
\]
从而得到唯一的 shell-constant 正权公式
\[
\sum_{w\in X_6^{\mathrm{cyc}}}\langle x,\alpha(w)\rangle^4
+
2\sum_{i=1}^{3}x_i^4
=
12|x|^4.
\]
因此 boundary 三重态在四阶通道中不是附加饰项，而是把 cyclic root cloud 恢复为精确 \(O(3)\)-isotropy 的唯一 counterterm。
\end{theorem}
\begin{proof}
由 window-\(6\) 的刚性直方图 \(\#\{d_6=2\}=8\)、\(\#\{d_6=3\}=4\)、\(\#\{d_6=4\}=9\) 可知，cyclic 壳层的基数分别为 \(5,4,9\)，其中 \(d_6=2\) 的另外 \(3\) 个状态正好是 boundary 三词。
结合表 \ref{tab:fold6_b3c3_root_dictionary_b3} 的根字典，直接展开可得
\[
\sum_{\beta\in\Phi_{6,2}^{\mathrm{cyc}}}\langle x,\beta\rangle^4
=
x_1^4+4x_2^4+4x_3^4+24x_2^2x_3^2,
\]
\[
\sum_{\beta\in\Phi_{6,3}^{\mathrm{cyc}}}\langle x,\beta\rangle^4
=
4x_1^4+4x_3^4+24x_1^2x_3^2,
\]
\[
\sum_{\beta\in\Phi_{6,4}^{\mathrm{cyc}}}\langle x,\beta\rangle^4
=
5x_1^4+6x_2^4+2x_3^4+24x_1^2x_2^2.
\]
若
\[
\mathcal Q_{a_2,a_3,a_4,b}(x)=c(x_1^2+x_2^2+x_3^2)^2,
\]
则比较交叉项
\[
x_1^2x_2^2,\qquad x_1^2x_3^2,\qquad x_2^2x_3^2
\]
的系数，得到
\[
24a_4=2c,\qquad 24a_3=2c,\qquad 24a_2=2c,
\]
故
\[
a_2=a_3=a_4=\frac{c}{12}.
\]
再比较 \(x_1^4\) 系数，得到
\[
\frac{c}{12}+4\cdot\frac{c}{12}+5\cdot\frac{c}{12}+b=c,
\]
即
\[
b=\frac{c}{6}=2a.
\]
反之，把 \(a_2=a_3=a_4=a\) 与 \(b=2a\) 代回，直接化简便得
\[
\mathcal Q_{a,a,a,2a}(x)=12a\,|x|^4.
\]
证毕。
\end{proof}

\begin{proposition}[六阶旋转不变的不可实现性]\label{prop:window6-visible-degree6-isotropy-obstruction}
设 \(\{c_w\}_{w\in X_6}\subset \CC\) 任意，并定义六次齐次多项式
\[
P_c(x):=\sum_{w\in X_6}c_w\langle x,v_w\rangle^6.
\]
则 \(P_c\) 中单项式 \(x_1^2x_2^2x_3^2\) 的系数恒等于零。因而不存在非零常数 \(\lambda\) 使
\[
P_c(x)=\lambda |x|^6.
\]
特别地，window-\(6\) 的 visible support 上精确旋转不变的偶矩只能维持到四阶；六阶处即出现不可消去的各向异性障碍。
\end{proposition}
\leanverified{paper\_window6\_visible\_degree6\_isotropy\_obstruction}
\begin{proof}
由 \(v_w\) 的定义，任意可见向量都至少有一个零坐标：cyclic 根向量只支撑在至多两条坐标轴上，而 boundary 三向量正是 \(e_1,e_2,e_3\)。
写
\[
v_w=(a_1,a_2,a_3),
\]
则
\[
\langle x,v_w\rangle^6=(a_1x_1+a_2x_2+a_3x_3)^6
\]
中，单项式 \(x_1^2x_2^2x_3^2\) 的系数为
\[
90\,a_1^2a_2^2a_3^2,
\]
而这在某个 \(a_i=0\) 时恒为零。
故对所有 \(w\) 的线性组合 \(P_c\)，该单项式系数始终为零。

另一方面，
\[
|x|^6=(x_1^2+x_2^2+x_3^2)^3
\]
中 \(x_1^2x_2^2x_3^2\) 的系数等于 \(6\)。
因此若 \(P_c(x)=\lambda |x|^6\)，必有 \(\lambda=0\)。证毕。
\end{proof}

\begin{theorem}[visible matroid 的 minor census 与 cubic volume budget]\label{thm:window6-visible-matroid-minor-census}
令
\[
A_6:=\bigl[v_w\bigr]_{w\in X_6}\in M_{3\times 21}(\ZZ).
\]
对每个三元子集 \(I\subset X_6\)（\(|I|=3\)），记 \(A_6[I]\) 为对应的 \(3\times 3\) 列子矩阵。则
\[
\det A_6[I]\in\{0,\pm 1,\pm 2\}
\qquad
(\forall I\in \tbinom{X_6}{3}).
\]
若定义
\[
N_r:=
\#\Bigl\{
I\in\tbinom{X_6}{3}:\ |\det A_6[I]|=r
\Bigr\},
\qquad
r\in\{0,1,2\},
\]
则有精确计数
\[
N_1=675,
\qquad
N_2=164,
\qquad
N_0=491.
\]
特别地，
\[
\sum_{I\in\binom{X_6}{3}}\det\!\bigl(A_6[I]\bigr)^2
=
675+4\cdot 164
=
1331
=
11^3.
\]
因此，window-\(6\) 的 visible matroid 并非一般的 rank-\(3\) 整矩阵结构，而是被一个 cubic budget \(11^3\) 精确控制。
\end{theorem}
\begin{proof}
每一列 \(v_w\) 都属于 \(\ZZ^3\)，且列向量的平方范数至多为 \(2\)。
由 Hadamard 不等式，
\[
|\det A_6[I]|^2\le 2\cdot 2\cdot 2=8,
\]
故 \(|\det A_6[I]|\le 2\)。又行列式是整数，于是只能取 \(0,\pm 1,\pm 2\)。

下面计算 \(N_1\)。把 \(A_6\) 模 \(2\) 约化为
\[
\bar A_6\in M_{3\times 21}(\FF_2).
\]
由表 \ref{tab:fold6_b3c3_root_dictionary_b3} 与 boundary 三基可知，\(\bar A_6\) 的列只落在六种非零类型上：
\[
100,\ 010,\ 001
\quad\text{各重数 }3;
\qquad
110,\ 101,\ 011
\quad\text{各重数 }4.
\]
三列子矩阵有奇行列式，当且仅当这三列在 \(\FF_2^3\) 中线性无关。按类型分三类计数：

第一类，三个轴型 \(100,010,001\) 构成唯一一组基，贡献
\[
3\cdot 3\cdot 3=27.
\]
第二类，两个不同轴型加一个双坐标型。先选两条轴，有 \(\binom 32=3\) 种；与其和相等的那一条双坐标型会造成线性相关，故可选双坐标型恰有 \(2\) 个，因此共有 \(6\) 组类型基，每组贡献
\[
3\cdot 3\cdot 4=36,
\]
总计 \(216\)。
第三类，一条轴型加两条不同双坐标型。任取轴型有 \(3\) 种，再从三条双坐标型中任选两条有 \(\binom 32=3\) 种；这 \(9\) 组类型全都线性无关，每组贡献
\[
3\cdot 4\cdot 4=48,
\]
总计 \(432\)。
故
\[
N_1=27+216+432=675.
\]

另一方面，由定理 \ref{thm:window6-b3c3-adjoint-second-moment-isotropy}，cyclic 部分给出 \(10I_3\)；再加上 boundary 三基的贡献
\[
e_1e_1^{\mathsf T}+e_2e_2^{\mathsf T}+e_3e_3^{\mathsf T}=I_3,
\]
得到
\[
A_6A_6^{\mathsf T}=11I_3.
\]
故由 Cauchy--Binet，
\[
\sum_{I\in\binom{X_6}{3}}\det\!\bigl(A_6[I]\bigr)^2
=
\det(A_6A_6^{\mathsf T})
=
\det(11I_3)
=
11^3
=
1331.
\]
因此
\[
N_1+4N_2=1331,
\]
代入 \(N_1=675\) 可得
\[
N_2=\frac{1331-675}{4}=164.
\]
最后
\[
N_0=\binom{21}{3}-N_1-N_2
=
1330-675-164
=
491.
\]
证毕。
\end{proof}

\leanverified{paper\_window6\_visible\_matroid\_minor\_census}
\leanverified{paper\_window6\_cartan\_kernel\_code\_parameters}
\begin{theorem}[Cartan kernel code 与其对偶码的精确参数]\label{thm:window6-cartan-kernel-code-parameters}
令 \(\bar A_6\in M_{3\times 21}(\FF_2)\) 为 \(A_6\) 的模 \(2\) 约化，并定义
\[
C_6^\perp:=\operatorname{rowspan}_{\FF_2}(\bar A_6)\subset \FF_2^{21},
\qquad
C_6:=\ker \bar A_6.
\]
则：
\begin{enumerate}
  \item \(C_6^\perp\) 是一个 \([21,3,9]\) 码，其 Hamming 重量枚举多项式恰为
  \[
  W_{C_6^\perp}(z)=1+z^9+3z^{11}+3z^{14}.
  \]
  \item \(C_6\) 是一个 \([21,18,2]\) 码，且其 Hamming 重量枚举由 MacWilliams 变换写成
  \[
  W_{C_6}(x,y)
  =
  \frac18\Bigl(
  (x+y)^{21}
  +(x+y)^{12}(x-y)^9
  +3(x+y)^{10}(x-y)^{11}
  +3(x+y)^7(x-y)^{14}
  \Bigr).
  \]
\end{enumerate}
因此，window-\(6\) 的 \(21\) 比特 anomaly lattice 在 mod-\(2\) Cartan 投影之下自然生成一对互为正交补的二元线性码：三维对偶码记录 visible Cartan 信息，而十八维 kernel code 记录全部 gauge-neutral 组合。
\end{theorem}
\begin{proof}
由 \(\bar A_6\) 中三条 boundary 列正是标准基 \(e_1,e_2,e_3\)，可知其行秩为 \(3\)，故
\[
\dim C_6^\perp=3,
\qquad
\dim C_6=21-3=18.
\]

再由定理 \ref{thm:window6-visible-matroid-minor-census} 证明中给出的六类二进制列型，可直接计算这 \(8\) 个行空间码字的重量：
零码字重量为 \(0\)；三条单独的行各有重量
\[
3+4+4=11;
\]
任意两行之和的重量为 \(14\)；三行总和的重量为 \(9\)。
于是
\[
W_{C_6^\perp}(z)=1+z^9+3z^{11}+3z^{14},
\]
其最小非零重量为 \(9\)，故 \(C_6^\perp\) 的参数为 \([21,3,9]\)。

对偶码 \(C_6\) 的重量枚举由 MacWilliams 恒等式给出：
\[
W_{C_6}(x,y)
=
\frac1{|C_6^\perp|}
W_{C_6^\perp}(x+y,x-y),
\]
而 \(|C_6^\perp|=2^3=8\)，代入上式即得所示闭式。

最后，为证明 \(d(C_6)=2\)，只需注意
\[
\bar A_6(\delta_{\texttt{100000}}+\delta_{\texttt{100001}})=0,
\]
因为 \(\texttt{100000}\) 与 \(b_1=\texttt{100001}\) 在 \(\FF_2^3\) 中都映为 \(100\)。
故 \(C_6\) 含有重量 \(2\) 的非零码字，从而其最小距离恰为 \(2\)。证毕。
\end{proof}

\noindent
由此，window-\(6\) 的可见 \(21\) 点模型不再只是一个根字典或纤维直方图，而被压缩为同一条整数线性闭链：
boundary 三基给出 visible Cartan quotient，
其核形成判别式 \(11^3\) 的秩 \(18\) syzygy lattice；
四阶处存在且仅存在一个 boundary counterterm 能把 cyclic cloud 恢复为精确 \(O(3)\)-isotropy，
而六阶处又出现不可消去的各向异性；
最后，这个 quotient 的模 \(2\) 影子恰生成一对 \([21,3,9]\) 与 \([21,18,2]\) 的正交补码。
于是，Cartan quotient、syzygy lattice、quartic counterterm 与 binary code 四层结构便在同一对象上闭合。




% This file is intended to be \input{} from subsubsec__window6_lie_bracket_closure.tex.

\paragraph{交换子包络、根字典与残余结构：从离散审计到可验证闭合链}
本段将 window-\(6\) 处“李括号的构造”与“唯一闭合”的两条对象链作同一口径下的衔接：其一是由坐标翻转推送算子生成的交换子 Lie 包络（算子链），其二是由根--Cartan 字典在 \(V=\RR\langle X_6\rangle\) 上诱导的抽象括号（类型链）。
两条链均以有限可审计数据为输入，但其闭合目标不同；因此相关猜想的定理化应当分别落在这两条对象链的相应命题上。

\begin{remark}[交换子包络的严格构造与唯一闭合]\label{rem:window6-commutator-envelope-unique-closure}
对每个坐标方向 \(i\in\{1,\dots,6\}\)，推送计数算子 \(L_i\) 由有限边计数确定（定义 \ref{def:window6-coordinate-pushforward-lie-envelope}），并以交换子 \([A,B]=AB-BA\) 作为李括号。
由最小交刻画，交换子包络
\(
\mathfrak g_6=\mathrm{Lie}\langle L_1,\dots,L_6\rangle
\)
作为包含生成元集合的最小 Lie 子代数存在且唯一（定理 \ref{thm:window6-lie-envelope-existence-uniqueness}），并可通过有限维闭包迭代在有限步内稳定（命题 \ref{prop:finite-dim-lie-closure-termination}）。
这些结论仅依赖有限维线性代数与交换子封闭性，不引入任何连续假设。
\end{remark}

\begin{remark}[关于根型预期的对象层区分]\label{rem:window6-conj-10170-object-correction}
在交换子口径下，\(\mathfrak g_6\) 的结构识别已由证书化审计闭合为满的 \(\mathfrak{sl}_{21}\)（定理 \ref{thm:window6-lie-envelope-sl21}），因此“\(\dim\mathfrak g_6=21\) 且复化根型为 \(B_3/C_3\)”这一形式的闭合断言与定义 \ref{def:window6-coordinate-pushforward-lie-envelope} 的交换子包络对象不相容。
由此，秩 \(3\) 的 \(21\) 维简单根型结构应当被理解为类型链上的抽象括号对象（定理 \ref{thm:window6-root-dictionary-pullback-bracket} 与猜想 \ref{conj:window6-intrinsic-bracket-uniqueness}），或作为 \(\mathfrak{sl}_{21}\) 内受限可观测紧致部分的抽取目标（猜想 \ref{conj:window6-sm-gauge-as-minimal-compact-observable}），而非 \(\mathfrak g_6\) 的交换子闭包结构本身。
\end{remark}

\begin{proposition}[低秩规范代数必须来自可观测压缩/兼容约束]\label{prop:window6-lowrank-needs-observable-compression}
在定义 \ref{def:window6-coordinate-pushforward-lie-envelope} 的设定下，
令 \(\mathfrak h\subseteq \mathfrak{sl}(V)\) 为 Lie 子代数并满足 \(L_1,\dots,L_6\in\mathfrak h\)。
则 \(\mathfrak h=\mathfrak{sl}(V)\)。
因此，若目标为低秩紧致规范代数（例如 \(\mathfrak{su}(3)\oplus\mathfrak{su}(2)\oplus\mathfrak u(1)\) 或 Levi 型 \(\mathfrak{su}(4)\oplus\mathfrak{su}(2)_L\oplus\mathfrak{su}(2)_R\)），则其不能作为仅由 \(\{L_i\}\) 在交换子括号下生成的闭包对象出现；
该类低秩对象必须在 \(\mathfrak g_6=\mathfrak{sl}(V)\) 的最大闭合之上，引入额外的“可观测压缩/兼容约束”以抽取为受限子代数。
\end{proposition}
\begin{proof}
由定理 \ref{thm:window6-lie-envelope-existence-uniqueness}，
\(\mathfrak g_6=\Lie\langle L_1,\dots,L_6\rangle\) 为包含全部生成元的最小 Lie 子代数。
因此 \(\mathfrak g_6\subseteq \mathfrak h\)。
又由定理 \ref{thm:window6-lie-envelope-sl21}，
\(\mathfrak g_6=\mathfrak{sl}(V)\)，故 \(\mathfrak h=\mathfrak{sl}(V)\)。
最后一断言为前述等式的直接推论：若不改变“包含全部 \(L_i\)”与“交换子闭包”这两项结构性输入，则闭包对象固定为 \(\mathfrak{sl}(V)\)，从而任何低秩目标都只能通过附加约束实现为受限子对象。证毕。
\end{proof}

\begin{remark}[可观测压缩/兼容约束的典型形态]\label{rem:window6-observable-compression-forms}
在本文 window-\(6\) 的审计口径下，“可观测压缩/兼容约束”可被组织为对 \(\mathfrak{sl}(V)\) 的一族实现无关的受限条件，其典型来源包括：
\begin{itemize}
  \item 与 orbit--Levi 骨架或 edge--flux 四块骨架相容的分块保持条件，从而把 \(\mathfrak{sl}(V)\) 限制到保持给定粗粒化接口的子代数族；
  \item 与 \(\delta=34\) 的 sheet parity 分级相容的分级保持条件，从而把允许的非阿贝尔混合限制为分级同态可见的部分；
  \item 由交换子（或条件期望）接口诱导的 commutant/截断约束，使得“可见对称”以可复核的交换子方程而非以生成闭包来定义。
\end{itemize}
这些约束的作用不是削弱算子链的闭合唯一性，而是将其作为背景最大闭合，使低秩紧致结构的出现转化为一组离散可审计的抽取问题。
\end{remark}

\begin{remark}[两类分块骨架与 SM 残余结构的兼容闭合]\label{rem:window6-two-skeletons-to-sm-residual}
window-\(6\) 同时提供两条彼此独立的离散分块接口：其一是 edge--flux 口径下的四块分解 \(21=9\oplus 6\oplus 3\oplus 3\)（终端命题 \ref{thm:terminal-window6-edge-flux-skeleton}），其二是 orbit--Levi 口径下的 \(21=15\oplus 3\oplus 3\) 维数骨架及相应 Levi 候选
\(
\mathfrak g^{\mathrm{Levi}}_6\cong \mathfrak{su}(4)\oplus\mathfrak{su}(2)_L\oplus\mathfrak{su}(2)_R
\)
（参见 \S\ref{subsec:bdry-tower-zeck-gut} 与命题 \ref{prop:window6-common-refinement-sm-levi}）。
当要求两条分块在同一连续包络/可观测语义下兼容闭合时，标准模型型 Levi 子结构
\(
\mathfrak{su}(3)\oplus\mathfrak{su}(2)_L\oplus\mathfrak u(1)
\)
被强制出现（命题 \ref{prop:window6-common-refinement-sm-levi}）。
进一步地，在 Levi 骨架之上，\(\mathfrak{su}(4)\supset\mathfrak{su}(3)\oplus\mathfrak u(1)\) 的嵌入刚性（推论 \ref{cor:window6-su4-embedding-rigidity}）与几何 \(\ZZ_2\)/sheet parity 审计共同给出一条有效阿贝尔方向 \(\mathfrak u(1)_{\mathrm{eff}}\subset\mathfrak{su}(2)_R\)，从而得到标准模型型残余代数（推论 \ref{cor:window6-levi-rigidity-sm-residual}）。
\end{remark}

\begin{remark}[抽象括号的存在性、唯一性与混合证书化路径]\label{rem:window6-bracket-theoremization-route}
根--Cartan 字典在 \(V=\RR\langle X_6\rangle\) 上给出一条完全显式的括号存在性构造：将 \(X_6^{\mathrm{bdry}}\) 锚定到 Cartan 基、将 \(X_6^{\mathrm{cyc}}\) 通过根字典锚定到根向量，从而拉回紧致 \(B_3/C_3\) 型括号（定理 \ref{thm:window6-root-dictionary-pullback-bracket}）。
在固定标号刚性、分级相容与分块相容的审计语义下，其唯一性可被组织为证书化断言（定理 \ref{thm:window6-intrinsic-bracket-certificate-theorem}）。
\par
在该框架中，括号的“存在/唯一”可进一步增强为混合证书问题：代数侧以反对称性与 Jacobi 恒等式、局域支撑模板及 parity/分块约束给出有限多项式系统（命题 \ref{prop:window6-bracket-certificate-reduction}），谱侧以 unitary slice 的 Toeplitz--PSD 证书链给出半正定约束（定理 \ref{thm:window6-p6-toeplitz-certificate-chain}、定理 \ref{thm:app-horizon-toeplitz-lmi}，并见定理 \ref{thm:window6-bracket-uniqueness-mixed-certificate-reduction}）。
由此，原先的“纯 Jacobi 多项式系统”可被替换为“多项式可行域 \(\cap\) 半正定可行域”的交汇判别，从而对非紧/非幺正分支给出可复核的排除机制。
\end{remark}

\begin{remark}[将猜想 \(10.170\)–\(10.172\) 升格为定理的缺口定位]\label{rem:window6-10-170-10-172-gaps}
交换子包络链的闭合与结构识别已在定理 \ref{thm:window6-lie-envelope-sl21} 中完成证书化。
剩余的定理化工作集中在两处：其一是把“可观测极小紧致部分”的存在与极小性准则给出为可审计判别（对应猜想 \ref{conj:window6-sm-gauge-as-minimal-compact-observable} 的实质内容）；其二是在类型链上把抽象括号的审计约束（局域/分级/分块）与混合证书约束（Toeplitz--PSD 等）合并为可输出证书的唯一可解性结论（对应猜想 \ref{conj:window6-intrinsic-bracket-uniqueness} 的定理化版本）。
两处缺口均可被组织为有限对象的核验任务，并可与整数模素数审计与 Lee--Yang/PSD 证书链并行交叉复核。
\end{remark}

\paragraph{仿射纤维层与满控制动力学层的二层刚性}
沿用定理 \ref{thm:terminal-foldbin6-fiber-affine-geometry}、推论 \ref{cor:fold-bin-m6-fiber-hist}、定义 \ref{def:window6-coordinate-pushforward-lie-envelope} 与定理 \ref{thm:window6-lie-envelope-sl21} 的记号。
记
\[
V:=\RR\langle X_6\rangle,\qquad
F_w:=\bigl(\Fold^{\mathrm{bin}}_6\bigr)^{-1}(w)\subset \mathbb F_2^6,
\qquad
d_6(w):=\card{F_w},
\qquad
r(w):=\dim \operatorname{Aff}(F_w).
\]
再记
\[
L:=\{0,e_1\}\subset \mathbb A^1(\mathbb F_2),\qquad
P^\circ:=\mathbb A^2(\mathbb F_2)\setminus\{0\},\qquad
P:=\mathbb A^2(\mathbb F_2),\qquad
\Sigma:=\{0,e_1,e_2,e_3\}\subset \mathbb A^3(\mathbb F_2).
\]

\begin{theorem}[window-\texorpdfstring{$6$}{6} 全纤维的仿射同构分类]\label{thm:window6-affine-fiber-orbit-classification}
全部 \(21\) 条纤维按 \(\mathbb F_2\)-仿射同构恰落入四类，且各类基数分别为
\[
\card{\{w\in X_6:\ F_w\cong L\}}=8,\qquad
\card{\{w\in X_6:\ F_w\cong P^\circ\}}=4,
\]
\[
\card{\{w\in X_6:\ F_w\cong P\}}=3,\qquad
\card{\{w\in X_6:\ F_w\cong \Sigma\}}=6.
\]
因此，window-\(6\) 的纤维几何不是连续族，而是一个严格有限的四轨道分类。
\end{theorem}
\leanverified{paper\_window6\_affine\_fiber\_orbit\_classification}
\begin{proof}
由定理 \ref{thm:terminal-foldbin6-fiber-affine-geometry} 及表 \ref{tab:foldbin6_fiber_affine_geometry}，全部纤维仅出现下列四种组合型：
\[
\bigl(d_6(w),r(w),\text{\(F_w\) 是否为仿射子空间}\bigr)
\in
\{(2,1,\mathrm{Y}),(3,2,\mathrm{N}),(4,2,\mathrm{Y}),(4,3,\mathrm{N})\}.
\]
若 \((d_6(w),r(w),\mathrm{Y})=(2,1,\mathrm{Y})\)，则 \(F_w\) 必为一条仿射直线，故同构于 \(L\)。
若 \((d_6(w),r(w),\mathrm{N})=(3,2,\mathrm{N})\)，则 \(F_w\) 位于某个仿射平面中且恰缺一个点，故同构于 \(P^\circ\)。
若 \((d_6(w),r(w),\mathrm{Y})=(4,2,\mathrm{Y})\)，则 \(F_w\) 本身即整个仿射平面，故同构于 \(P\)。

最后设 \((d_6(w),r(w),\mathrm{N})=(4,3,\mathrm{N})\)。
任选 \(p\in F_w\) 并把 \(F_w-p\) 平移到经过原点的集合，则
\[
0\in F_w-p\subset \mathbb F_2^6,\qquad \dim\operatorname{span}(F_w-p)=3.
\]
除原点外的三点必须线性无关；否则其线性张成维数至多为 \(2\)，与 \(r(w)=3\) 矛盾。
因此 \(F_w-p\) 与 \(\{0,e_1,e_2,e_3\}\) 线性同构，从而 \(F_w\cong \Sigma\)。
四类出现次数由表 \ref{tab:foldbin6_fiber_affine_geometry} 逐行计数即得 \(8,4,3,6\)。证毕。
\end{proof}

\begin{corollary}[退化壳层的仿射裂解与零温 escort 极限]\label{cor:window6-degenerate-shell-affine-splitting}
\leanverified{paper\_window6\_degenerate\_shell\_affine\_splitting}
对 \(j\in\{2,3,4\}\) 定义退化壳层
\[
S_{6,j}:=\{w\in X_6:\ d_6(w)=j\}.
\]
再定义
\[
S_{6,4}^{\mathrm{pl}}:=\{w\in X_6:\ F_w\cong P\},
\qquad
S_{6,4}^{\mathrm{simp}}:=\{w\in X_6:\ F_w\cong \Sigma\}.
\]
则有严格分裂
\[
S_{6,2}=\{w:F_w\cong L\},\qquad \card{S_{6,2}}=8,
\]
\[
S_{6,3}=\{w:F_w\cong P^\circ\},\qquad \card{S_{6,3}}=4,
\]
\[
S_{6,4}=S_{6,4}^{\mathrm{pl}}\sqcup S_{6,4}^{\mathrm{simp}},
\qquad
\card{S_{6,4}^{\mathrm{pl}}}=3,
\qquad
\card{S_{6,4}^{\mathrm{simp}}}=6.
\]
特别地，最大退化壳层在仿射几何上并不均匀，而是按比率 \(3:6=1:2\) 裂解为平面型与单纯形型两部分。

对任意 \(q\ge 0\) 定义 escort 分布
\[
\pi_{6,q}(w):=\frac{d_6(w)^q}{\sum_{x\in X_6} d_6(x)^q},
\qquad w\in X_6,
\]
并记 \(u_A\) 为有限集 \(A\) 上的均匀分布。
则当 \(q\to+\infty\) 时有
\[
\pi_{6,q}\longrightarrow \frac13\,u_{S_{6,4}^{\mathrm{pl}}}+\frac23\,u_{S_{6,4}^{\mathrm{simp}}}.
\]
\end{corollary}
\begin{proof}
前半部分由定理 \ref{thm:window6-affine-fiber-orbit-classification} 与推论 \ref{cor:fold-bin-m6-fiber-hist} 的直方图
\[
2:8,\qquad 3:4,\qquad 4:9
\]
直接对应得到。

对 escort 极限，只需将分子分母同除以 \(4^q\)。
当 \(q\to+\infty\) 时，所有满足 \(d_6(w)<4\) 的项都消失，仅 \(S_{6,4}\) 保留质量；而在 \(S_{6,4}\) 内各点权重相同，故极限分布即为 \(S_{6,4}\) 上的均匀分布。
再利用
\[
\card{S_{6,4}^{\mathrm{pl}}}=3,\qquad \card{S_{6,4}^{\mathrm{simp}}}=6,
\]
便得到
\[
u_{S_{6,4}}
=\frac{3}{9}\,u_{S_{6,4}^{\mathrm{pl}}}+\frac{6}{9}\,u_{S_{6,4}^{\mathrm{simp}}}
=\frac13\,u_{S_{6,4}^{\mathrm{pl}}}+\frac23\,u_{S_{6,4}^{\mathrm{simp}}}.
\]
证毕。
\end{proof}

\begin{theorem}[仿射完成体积与 \texorpdfstring{$4\cdot 23$}{4*23} 因子化]\label{thm:window6-affine-completion-volume-4x23}
定义 window-\(6\) 的最小仿射完成体积
\[
\mathfrak A_6:=\sum_{w\in X_6}\card{\operatorname{Aff}(F_w)}
=\sum_{w\in X_6}2^{r(w)}.
\]
则
\[
\mathfrak A_6
=8\cdot 2+4\cdot 4+3\cdot 4+6\cdot 8
=92
=4\cdot 23.
\]
并且其以微态测度 \(\pi_6(w)=d_6(w)/64\) 加权的仿射完成平均满足
\[
\sum_{w\in X_6}\pi_6(w)\,\card{\operatorname{Aff}(F_w)}
=\frac1{64}\sum_{w\in X_6}d_6(w)\,\card{\operatorname{Aff}(F_w)}
=5.
\]
\end{theorem}
\begin{proof}
由定理 \ref{thm:window6-affine-fiber-orbit-classification}，四类纤维的仿射包络大小分别为
\[
\card{\operatorname{Aff}(L)}=2,\qquad
\card{\operatorname{Aff}(P^\circ)}=\card{\operatorname{Aff}(P)}=4,\qquad
\card{\operatorname{Aff}(\Sigma)}=8.
\]
故
\[
\mathfrak A_6=8\cdot 2+4\cdot 4+3\cdot 4+6\cdot 8=92=4\cdot 23.
\]
再按纤维大小加权，
\[
\sum_{w\in X_6}d_6(w)\,\card{\operatorname{Aff}(F_w)}
=8\cdot 2\cdot 2+4\cdot 3\cdot 4+3\cdot 4\cdot 4+6\cdot 4\cdot 8
=320.
\]
除以总微态数 \(64\) 即得所示平均值 \(5\)。证毕。
\end{proof}
\leanverified{paper\_window6\_affine\_completion\_volume\_4x23}

\begin{theorem}[edge--flux 结合包络的满矩阵饱和]\label{thm:window6-edge-flux-full-matrix-saturation}
令
\[
\mathcal B_6:=\RR\langle I,L_1,\dots,L_6\rangle\subset \mathrm{End}(V)
\]
为由六个 pushforward edge operators 与单位元生成的结合代数。
则
\[
\mathcal B_6=\mathrm{End}(V)\cong M_{21}(\RR).
\]
换言之，window-\(6\) 的 edge--flux 动力学在结合层面直接饱和为全部 \(21\times 21\) 实矩阵。
\end{theorem}
\begin{proof}
任意包含 \(L_1,\dots,L_6\) 的结合代数都对交换子封闭，因此必包含它们生成的 Lie 包络。
由定理 \ref{thm:window6-lie-envelope-sl21}，
\[
\mathfrak g_6:=\Lie\langle L_1,\dots,L_6\rangle=\mathfrak{sl}(V)\subseteq \mathcal B_6.
\]
又因为 \(I\in \mathcal B_6\)，故
\[
\mathfrak{gl}(V)=\mathfrak{sl}(V)\oplus \RR I\subseteq \mathcal B_6.
\]
另一方面 \(\mathcal B_6\subseteq \mathrm{End}(V)=\mathfrak{gl}(V)\)。
两侧夹逼即得 \(\mathcal B_6=\mathrm{End}(V)\)。证毕。
\end{proof}
\leanverified{paper\_window6\_edge\_flux\_full\_matrix\_saturation}

\begin{corollary}[精确投影的非交换综合]\label{cor:window6-exact-projector-synthesis}
对任意子集 \(Y\subseteq X_6\)，记 \(P_Y\in \mathrm{End}(V)\) 为坐标投影
\[
P_Y(e_w)=
\begin{cases}
e_w,& w\in Y,\\
0,& w\notin Y.
\end{cases}
\]
则存在一个带常数项的六变量非交换实多项式 \(p_Y\in \RR\langle x_1,\dots,x_6\rangle\)，使得
\[
P_Y=p_Y(L_1,\dots,L_6).
\]
特别地，下列几何子空间都可由 edge--flux generators 精确综合：
\[
\RR\langle X_6^{\mathrm{bdry}}\rangle,\qquad
\RR\langle S_{6,2}\rangle,\ \RR\langle S_{6,3}\rangle,\ \RR\langle S_{6,4}\rangle,
\]
\[
\RR\langle \{w:F_w\cong L\}\rangle,\ \RR\langle \{w:F_w\cong P^\circ\}\rangle,\ 
\RR\langle \{w:F_w\cong P\}\rangle,\ \RR\langle \{w:F_w\cong \Sigma\}\rangle,
\]
以及每一条单独的稳定类型坐标直线 \(\RR e_w\)。
\end{corollary}
\begin{proof}
由定理 \ref{thm:window6-edge-flux-full-matrix-saturation}，
\[
\mathcal B_6=\mathrm{End}(V).
\]
而任意坐标投影 \(P_Y\) 都属于 \(\mathrm{End}(V)\)，故 \(P_Y\in \mathcal B_6\)。
按 \(\mathcal B_6\) 的定义，\(P_Y\) 必可写为 \(I,L_1,\dots,L_6\) 的有限实线性组合与有限字乘积之和，即某个非交换多项式 \(p_Y(L_1,\dots,L_6)\)。
后述各子空间只是不同比例尺下的 \(Y\) 所对应的特例。证毕。
\end{proof}
\leanverified{paper\_window6\_exact\_projector\_synthesis}

\begin{theorem}[joint commutant 的标量刚性与无不变子空间]\label{thm:window6-edge-flux-commutant-rigidity}
六个 pushforward edge operators 的 joint commutant 恰为标量：
\[
\operatorname{Comm}(L_1,\dots,L_6)
:=
\{T\in \mathrm{End}(V):\ TL_i=L_iT\ \forall i\}
=\RR I.
\]
特别地，\(V\) 上不存在任何同时对全部 \(L_i\) 不变的非零真子空间，因而也不存在非平凡的公共不变直和分解。
因此，boundary/cyclic、轻/重以及 \(L/P^\circ/P/\Sigma\) 等全部几何分类都不是完整 edge--flux 动力学的对称扇区；它们只是可被精确综合的投影阴影（projector-shadow）。
\end{theorem}
\begin{proof}
由定义，\(\operatorname{Comm}(L_1,\dots,L_6)\) 等于结合代数 \(\mathcal B_6\) 的交换子代数。
定理 \ref{thm:window6-edge-flux-full-matrix-saturation} 给出
\[
\mathcal B_6=\mathrm{End}(V),
\]
故
\[
\operatorname{Comm}(L_1,\dots,L_6)=\operatorname{Comm}(\mathrm{End}(V))=\RR I.
\]

若 \(W\subset V\) 是一个同时对所有 \(L_i\) 不变的非零子空间，则它也对 \(\mathcal B_6\) 不变，从而对整个 \(\mathrm{End}(V)\) 不变。
但 \(\mathrm{End}(V)\) 对标准模的作用显然不可约：任取 \(0\neq v\in W\) 与任意 \(u\in V\)，总可选取某个 \(T\in \mathrm{End}(V)\) 使 \(Tv=u\)，故 \(u\in W\)，于是 \(W=V\)。
因此不存在非零真公共不变子空间；非平凡公共不变直和分解更不可能存在。证毕。
\end{proof}

\leanverified{paper\_window6\_edge\_flux\_commutant\_rigidity}

\begin{proposition}[满 edge--flux 动力学下不存在非零不变双线性型]\label{prop:window6-edge-flux-no-invariant-bilinear-form}
不存在非零双线性型
\[
B:V\times V\to \RR
\]
使得对所有 \(X\in \mathfrak g_6=\mathfrak{sl}(V)\) 与 \(u,v\in V\) 恒有
\[
B(Xu,v)+B(u,Xv)=0.
\]
特别地，full edge--flux 动力学不保留任何非平凡内积、辛型或一般双线性配对。
\end{proposition}
\begin{proof}
取 \(V\) 的标准基并将 \(B\) 写成矩阵形式
\[
B(u,v)=u^{\mathsf T}Sv,\qquad S\in M_{21}(\RR).
\]
不变性等价于
\[
X^{\mathsf T}S+SX=0\qquad(\forall X\in \mathfrak{sl}(V)).
\]
对任意 \(i\neq j\) 取矩阵单位 \(X=E_{ij}\in \mathfrak{sl}(V)\)，则
\[
E_{ji}S+SE_{ij}=0.
\]
固定列指标 \(i\) 与行指标 \(a\)。
由于 \(\dim V=21\ge 3\)，可选 \(j\) 使 \(j\neq i\) 且 \(j\neq a\)。
取上式的 \((a,j)\)-元便得
\[
0=(E_{ji}S+SE_{ij})_{aj}=S_{ai},
\]
因为 \(a\neq j\)。
故第 \(i\) 列的所有元素都为零。
列指标 \(i\) 任意，因而 \(S=0\)，即 \(B=0\)。证毕。
\leanverified{paper\_window6\_edge\_flux\_no\_invariant\_bilinear\_form}
\end{proof}

\leanverified{paper\_window6\_edge\_flux\_no\_invariant\_binary\_product}
\begin{theorem}[满 edge--flux 动力学下不存在非零不变乘法]\label{thm:window6-edge-flux-no-invariant-binary-product}
不存在非零双线性映射
\[
m:V\times V\to V
\]
使其对 \(\mathfrak g_6=\mathfrak{sl}(V)\) 的标准作用不变。
等价地，
\[
\operatorname{Hom}_{\mathfrak{sl}(V)}(V\otimes V,V)=0.
\]
因此，在 window-\(6\) 的 full edge--flux dynamics 下，不存在非零不变的交换乘法、Jordan 乘法、Leibniz 乘法，亦不存在非零不变 Lie 括号。
\end{theorem}
\begin{proof}
将一切复化到 \(\mathbb C\) 上。
则 \(V_{\mathbb C}\) 是 \(\mathfrak{sl}(V_{\mathbb C})\cong \mathfrak{sl}_{21}(\mathbb C)\) 的标准不可约模，其最高权为 \(\omega_1\)。
另一方面，
\[
V_{\mathbb C}\otimes V_{\mathbb C}\cong S^2V_{\mathbb C}\oplus \Lambda^2V_{\mathbb C},
\]
其中 \(S^2V_{\mathbb C}\) 与 \(\Lambda^2V_{\mathbb C}\) 分别是最高权 \(2\omega_1\) 与 \(\omega_2\) 的不可约模。
二者都不含最高权 \(\omega_1\) 的分量，因此
\[
\operatorname{Hom}_{\mathfrak{sl}_{21}(\mathbb C)}(V_{\mathbb C}\otimes V_{\mathbb C},V_{\mathbb C})=0.
\]
再由复化忠实性，便得
\[
\operatorname{Hom}_{\mathfrak{sl}(V)}(V\otimes V,V)=0.
\]
证毕。
\end{proof}

\begin{corollary}[intrinsic \texorpdfstring{$\mathrm{Spin}(7)$}{Spin(7)} 括号的动力学异常]\label{cor:window6-spin7-bracket-dynamic-anomaly}
在定理 \ref{thm:window6-root-dictionary-pullback-bracket} 的 \(B_3\) 分支上，记
\[
[\cdot,\cdot]_{\mathrm{int}}:\Lambda^2V\to V
\]
为拉回的内禀括号，则 \((V,[\cdot,\cdot]_{\mathrm{int}})\cong \mathfrak{so}(7)\cong \operatorname{Lie}(\mathrm{Spin}(7))\)。
此时 \(\mathfrak g_6=\mathfrak{sl}(V)\) 的标准作用
\[
X\cdot v:=Xv\qquad (X\in \mathfrak g_6,\ v\in V)
\]
不可能提升为 \((V,[\cdot,\cdot]_{\mathrm{int}})\) 的导子作用。
等价地，不存在 Lie 代数同态
\[
\rho:\mathfrak g_6\to \operatorname{Der}(V,[\cdot,\cdot]_{\mathrm{int}})
\]
满足 \(\rho(X)v=Xv\) 对所有 \(X\in \mathfrak g_6\)、\(v\in V\) 成立。
因此，window-\(6\) 中的低秩内禀 Lie 几何与六坐标 pushforward 所生成的满控制动力学之间存在不可消去的异常裂缝：前者给出运动学括号，后者给出动力学包络，二者不能在同一标准作用下兼容。
\leanverified{paper\_window6\_spin7\_bracket\_dynamic\_anomaly}
\end{corollary}
\begin{proof}
反设存在这样的导子作用。
则对一切 \(X\in \mathfrak g_6\) 与 \(u,v\in V\) 都有
\[
X[u,v]_{\mathrm{int}}=[Xu,v]_{\mathrm{int}}+[u,Xv]_{\mathrm{int}},
\]
这恰说明括号映射
\[
[\cdot,\cdot]_{\mathrm{int}}:\Lambda^2V\to V
\]
是 \(\mathfrak g_6\)-等变的非零双线性映射。
但由定理 \ref{thm:window6-edge-flux-no-invariant-binary-product}，
\[
\operatorname{Hom}_{\mathfrak{sl}(V)}(V\otimes V,V)=0,
\]
故其在反对称部分上的限制
\(
\operatorname{Hom}_{\mathfrak{sl}(V)}(\Lambda^2V,V)
\)
也必为零。
于是 \([\cdot,\cdot]_{\mathrm{int}}\) 只能是零映射，这与 \((V,[\cdot,\cdot]_{\mathrm{int}})\cong \mathfrak{so}(7)\) 的非阿贝尔性矛盾。证毕。
\end{proof}

\begin{conjecture}[动态--运动学二层刚性]\label{conj:window-even-dynamic-kinematic-two-layer-rigidity}
设 \(m\) 为任意偶数，\(V_m:=\RR^{X_m}\)，并令 \(L_1^{(m)},\dots,L_m^{(m)}\in \mathrm{End}(V_m)\) 为对应的 pushforward edge operators。
假设同时满足以下两条：
\begin{enumerate}
  \item \(V_m\) 携带某个刚性的低秩内禀 Lie 完成 \(K_m\)；
  \item 六位情形的现象在 \(m\) 维上稳定为
  \[
  \Lie\langle L_1^{(m)},\dots,L_m^{(m)}\rangle=\mathfrak{sl}(V_m).
  \]
\end{enumerate}
则应同时推出：
\begin{enumerate}
  \item \(\RR\langle I,L_1^{(m)},\dots,L_m^{(m)}\rangle=\mathrm{End}(V_m)\)；
  \item \(\operatorname{Comm}(L_1^{(m)},\dots,L_m^{(m)})=\RR I\)，并且标准作用无非平凡公共不变子空间；
  \item \(\operatorname{Hom}_{\mathfrak{sl}(V_m)}(V_m\otimes V_m,V_m)=0\)；
  \item \(K_m\) 的内禀括号只能作为可被精确综合的投影阴影（projector-shadow）出现，而不可能成为完整 edge--flux 标准作用下的不变乘法。
\end{enumerate}
window-\(6\) 给出该猜想的首个完全审计实例：其一方面具有四类仿射纤维几何与 \(B_3\)/\(\mathrm{Spin}(7)\) 型的低秩运动学结构，另一方面又由六个坐标 pushforward 触发满的 \(\mathfrak{sl}_{21}\) 乃至 \(M_{21}(\RR)\) 动力学包络。
\end{conjecture}

\begin{remark}[window-\texorpdfstring{$6$}{6} 的二层结构]\label{rem:window6-affine-dynamics-two-layer-structure}
定理 \ref{thm:window6-affine-fiber-orbit-classification}--\ref{thm:window6-affine-completion-volume-4x23} 表明，window-\(6\) 的内禀几何层首先被压缩为一个有限的仿射模型簇：全部纤维仅分成 \(L\)、\(P^\circ\)、\(P\) 与 \(\Sigma\) 四类，最大退化壳层进一步按 \(1:2\) 裂解为平面型与单纯形型。
在此之上，定理 \ref{thm:window6-edge-flux-full-matrix-saturation}--推论 \ref{cor:window6-spin7-bracket-dynamic-anomaly} 又表明，完整 edge--flux 动力学在同一 \(21\) 维载体上直接饱和为全部矩阵代数，并由此抹平全部非平凡对称扇区、非零不变双线性型与非零不变乘法。
因此，window-\(6\) 的最深层结构不是单层 Lie 对象，而是一个严格二层系统：低秩几何层真实存在，但在满控制动力学下只能以投影阴影（projector-shadow）的形式被精确综合，而不能作为守恒的动力学对称保留下来。
\end{remark}

\endinput


% This file is intended to be \input{} from subsubsec__window6_lie_bracket_closure.tex.

\paragraph{Window-\(6\) 续证：6D 超立方体推送交换子包络与兼容闭合的唯一性链}
本段续接“window-\(6\) 的李括号构造与唯一闭合”的证明链条，并将论证拆分为两段互补但不混同的对象链。
第一段在 \(6\) 维超立方体 \(Q_6\) 与 dyadic bin-fold 标号 \(F_6\) 的推送口径下，
给出由有限边计数完全决定的交换子包络，从而把李括号固定为矩阵交换子 \([A,B]=AB-BA\)，并给出可复核的闭合识别证书。
第二段在 orbit--Levi 分块与阈值轻/重分裂要求同一包络结构兼容闭合时，推出标准模型型 Levi 子结构的强制出现与同构意义唯一性。
\par
为便于核验，沿用定义 \ref{def:window6-coordinate-pushforward-lie-envelope} 的符号，并回忆：
\[
\Omega_6:=\{0,1\}^6,\qquad
\mathrm{int}_6(\omega_1\cdots\omega_6):=\sum_{k=1}^{6}\omega_k\,2^{6-k}\in\{0,\dots,63\},
\qquad
F_6(\omega):=\mathrm{Fold}^{\mathrm{bin}}_6(\mathrm{int}_6(\omega))\in X_6.
\]
对每个坐标方向 \(i\in\{1,\dots,6\}\) 的翻转边 \(\omega\mapsto \omega\oplus e_i\)，推送计数算子 \(L_i\) 由
\[
(L_i)_{uv}=\#\{\omega\in\Omega_6:\ F_6(\omega)=u,\ F_6(\omega\oplus e_i)=v\}\qquad (u,v\in X_6)
\]
确定，并以交换子作为李括号定义
\(
\mathfrak g_6:=\Lie\langle L_1,\dots,L_6\rangle\subset\mathfrak{gl}(V)
\)。

\begin{theorem}[推送交换子包络的可核验闭合]\label{thm:window6-push-envelope-certificate-upgrade}
令 \(V=\RR\langle X_6\rangle\) 且 \(|X_6|=21\)。
对定义 \ref{def:window6-coordinate-pushforward-lie-envelope} 的生成元 \(\{L_i\}_{i=1}^{6}\)，有
\[
\mathfrak{k}_6:=\Lie\langle [L_i,L_j]:1\le i<j\le 6\rangle=\mathfrak{so}(V),\qquad
\dim \mathfrak{k}_6=210,
\]
\[
\mathfrak{g}_6:=\Lie\langle L_1,\dots,L_6\rangle=\mathfrak{sl}(V),\qquad
\dim \mathfrak{g}_6=440.
\]
并且相应维数与结构不变量由
\eqref{eq:window6_lie_envelope_orthogonal_part_dim}
和
\eqref{eq:window6_lie_envelope_invariants}
给出可复核证书。
\end{theorem}
\begin{proof}
由引理 \ref{lem:window6-li-symmetric}，生成元 \(L_i\) 为对称矩阵，
因而任意交换子 \([L_i,L_j]\) 为斜对称矩阵并落在 \(\mathfrak{so}(V)\)。
由引理 \ref{lem:window6-li-traceless}，每个 \(L_i\) 均迹零，故 \(\mathfrak g_6\subseteq\mathfrak{sl}(V)\)。
\par
交换子部分 \(\mathfrak{k}_6=\mathfrak{so}(V)\) 由定理 \ref{thm:window6-lie-envelope-so21-from-commutators} 给出；
全包络部分 \(\mathfrak{g}_6=\mathfrak{sl}(V)\) 由定理 \ref{thm:window6-lie-envelope-sl21} 给出。
相应维数证书可由模素数闭包核验输出（脚本 \texttt{\detokenize{scripts/exp_window6_lie_envelope_orthogonal_part.py}} 与 \texttt{\detokenize{scripts/exp_window6_lie_envelope.py}}），
并以 \eqref{eq:window6_lie_envelope_orthogonal_part_dim} 与 \eqref{eq:window6_lie_envelope_invariants} 的形式嵌入正文供离线复核。
证毕。
\end{proof}

\begin{remark}[与低秩预期对象的差异定位]\label{rem:window6-conjectural-gap-object-layer}
定理 \ref{thm:window6-push-envelope-certificate-upgrade} 说明：
若对象固定为定义 \ref{def:window6-coordinate-pushforward-lie-envelope} 的无约束交换子闭包，
则其结果是 \(\mathfrak{sl}_{21}\) 的极大包络，而非秩 \(3\)、维数 \(21\) 的低秩简单对象。
因此，低秩根型结构必须落在以下两类对象中之一：
\begin{enumerate}
  \item 内禀括号链（猜想 \ref{conj:window6-intrinsic-bracket-uniqueness}）；
  \item 在 \(\mathfrak{sl}_{21}\) 内加入可观测压缩/兼容约束后抽取的紧致子对象（猜想 \ref{conj:window6-sm-gauge-as-minimal-compact-observable}）。
\end{enumerate}
这一定义层分离与注记 \ref{rem:window6-conj-10170-object-correction} 一致。
\end{remark}

\begin{remark}[orbit--Levi 与阈值轻/重分裂的兼容闭合强制 SM Levi]\label{rem:window6-orbit-levi-threshold-compatible-closure}
交换子包络链条给出的是最大闭合背景 \(\mathfrak g_6=\mathfrak{sl}(V)\)，其本身不产生低秩紧致规范对象（命题 \ref{prop:window6-lowrank-needs-observable-compression}）。
因此，“唯一闭合”的可读紧致子结构必须来自在该最大背景上施加的实现无关兼容约束。
\par
在 window-\(6\) 的离散审计口径下，orbit--Levi 分块 \(21=15\oplus 3\oplus 3\) 与阈值轻/重分裂 \(X_6=1\oplus 8\oplus 12\) 为两条彼此独立的分解接口。
当要求二者在同一连续包络/可观测语义下同时出现并兼容闭合，且轻扇区 \(1\oplus 8\) 在 \(15\) 维块内闭合时，
命题 \ref{prop:window6-common-refinement-sm-levi} 强制给出 SM 型 Levi 子结构
\(
\mathfrak{su}(3)\oplus\mathfrak{su}(2)_L\oplus\mathfrak u(1)
\)
的出现与同构意义唯一性。
进一步的 \(\mathfrak u(1)\) 方向刚性可由 \(\mathfrak{su}(4)\supset \mathfrak{su}(3)\oplus\mathfrak u(1)\) 的嵌入共轭唯一性（推论 \ref{cor:window6-su4-embedding-rigidity}）
以及几何 \(\ZZ_2\)/sheet parity 的审计输入共同锁定（推论 \ref{cor:window6-levi-rigidity-sm-residual}）。
\end{remark}

\begin{proposition}[内禀括号存在唯一性的有限整数系统归约]\label{prop:window6-intrinsic-bracket-finite-integer-reduction}
取基 \(B=\{e_a\}_{a\in X_6}\) 并设
\[
[e_a,e_b]_6=\sum_{c\in X_6}C_{ab}^{c}\,e_c,\qquad C_{ab}^{c}\in\ZZ.
\]
则猜想 \ref{conj:window6-intrinsic-bracket-uniqueness} 的存在与唯一性可等价归约为
\(\mathbf C=(C_{ab}^{c})\) 的有限整系数方程组
\[
\mathcal F(\mathbf C)=0
\]
的可解性与唯一性，其中 \(\mathcal F\) 由以下约束并成：
\begin{enumerate}
  \item \textbf{反对称约束：}
  \(
  C_{ab}^{c}+C_{ba}^{c}=0
  \)；
  \item \textbf{局域支撑约束：}
  若 \((a,b;c)\) 不在一跳推送邻接与边多重度确定的允许模板 \(\mathcal S\) 中，则
  \(
  C_{ab}^{c}=0
  \)；
  \item \textbf{parity 约束：}
  在 \(\delta=34\) 分级下，
  \(
  \deg(a)+\deg(b)\not\equiv \deg(c)\ (\mathrm{mod}\ 2)
  \Rightarrow C_{ab}^{c}=0
  \)
  （推论 \ref{cor:terminal-delta-parity-rule}）；
  \item \textbf{根--Cartan 兼容约束：}
  与定理 \ref{thm:window6-root-dictionary-pullback-bracket} 的根字典接口一致；
  \item \textbf{分块兼容约束：}
  与 orbit--Levi 与 edge--flux 双骨架的分块闭包一致
  （命题 \ref{prop:window6-common-refinement-sm-levi}）；
  \item \textbf{Jacobi 约束：}
  \[
  [e_a,[e_b,e_c]_6]_6+[e_b,[e_c,e_a]_6]_6+[e_c,[e_a,e_b]_6]_6=0
  \qquad (\forall a,b,c\in X_6).
  \]
\end{enumerate}
此外，类型邻接图刚性 \(\mathrm{Aut}(\Gamma_6)=\{1\}\)（命题 \ref{prop:terminal-gamma6-rigidity}）
消除类型重标号自由度，使“同构意义唯一”可压缩为固定标号下的代数唯一性判别。
\end{proposition}
\begin{proof}
命题 \ref{prop:window6-bracket-certificate-reduction} 已给出“反对称 + Jacobi + 局域/分级/分块”到有限多项式系统的归约框架。
本命题只是将该框架以结构常数 \(\mathbf C\) 显式化，并把 parity 与双分块接口写成线性稀疏约束。
根字典接口提供坐标系，邻接图刚性排除重标号吸收非唯一性，故唯一性可在固定标号上判别。证毕。
\end{proof}

\begin{theorem}[Lee--Yang 式模素数提升判据（内禀括号版）]\label{thm:window6-intrinsic-bracket-modp-lift-criterion}
设 \(\mathcal F(\mathbf C)=0\) 为命题
\ref{prop:window6-intrinsic-bracket-finite-integer-reduction}
中的整系数系统，
并加入线性归一化条件 \(\mathcal N(\mathbf C)=0\) 以固定有限规范自由度。
若存在奇素数 \(p\) 使得：
\begin{enumerate}
  \item 约化系统
  \(
  (\mathcal F,\mathcal N)\equiv 0\pmod p
  \)
  在 \(\FF_p\) 上有唯一解 \(\bar{\mathbf C}\)；
  \item Jacobi 矩阵
  \[
  J_p(\bar{\mathbf C})
  :=\frac{\partial(\mathcal F,\mathcal N)}{\partial\mathbf C}(\bar{\mathbf C})
  \]
  在 \(\FF_p\) 上满秩；
  \item 存在整数向量 \(\mathbf C^\star\) 满足
  \[
  (\mathcal F,\mathcal N)(\mathbf C^\star)=0,\qquad
  \mathbf C^\star\equiv \bar{\mathbf C}\pmod p,
  \]
  且经直接代入可逐式核验。
\end{enumerate}
则 \(\mathbf C^\star\) 在该归一化下为唯一整数解；因此
\((V,[-,-]_6)\) 的内禀括号在特征 \(0\) 上存在且唯一。
\end{theorem}
\begin{proof}
由 (1)(2)，\(\bar{\mathbf C}\) 是模 \(p\) 的简单根，按 Hensel 提升定理存在唯一 \(p\)-进提升。
任意整数解在模 \(p\) 约化后都必须等于 \(\bar{\mathbf C}\)，故其 \(p\)-进像与该唯一提升一致。
由 (3) 已给出一个满足全部整系数方程的整数代表 \(\mathbf C^\star\)。
若存在另一整数解 \(\mathbf C'\)，则其模 \(p\) 约化同为 \(\bar{\mathbf C}\)，因此 \(\mathbf C'\) 与 \(\mathbf C^\star\) 在 \(\ZZ_p\) 中对应同一唯一提升根。
由于嵌入 \(\ZZ^N\hookrightarrow \ZZ_p^N\) 为单射，必有 \(\mathbf C'=\mathbf C^\star\)。
故唯一性成立。证毕。
\end{proof}

\begin{corollary}[内禀括号定理化的可执行证书链]\label{cor:window6-intrinsic-bracket-theoremization-pipeline}
在定理 \ref{thm:window6-push-envelope-certificate-upgrade} 的推送包络证书基础上，
将猜想 \ref{conj:window6-intrinsic-bracket-uniqueness} 升格为定理可按如下链条执行：
\begin{enumerate}
  \item 由 \(L_i\)、\(\Gamma_6\)、\(\delta=34\) 分级与双分块骨架生成 \(\mathcal S\) 与 \((\mathcal F,\mathcal N)\)；
  \item 在大素数域上计算唯一模 \(p\) 解与 Jacobian 满秩证书；
  \item 回到整数域执行直接代入核验并输出结构常数表；
  \item 结合定理 \ref{thm:window6-bracket-uniqueness-mixed-certificate-reduction} 的 Toeplitz--PSD 侧证书，完成代数侧与谱侧的交叉一致性审计。
\end{enumerate}
由此，“存在唯一性”被压缩为有限对象上的可判定任务，并与 Lee--Yang 式整数/模素数证书方法保持同型。
\end{corollary}
\begin{proof}
第 (1) 步由命题
\ref{prop:window6-intrinsic-bracket-finite-integer-reduction}
给出；第 (2)(3) 步由定理
\ref{thm:window6-intrinsic-bracket-modp-lift-criterion}
给出；第 (4) 步由定理
\ref{thm:window6-bracket-uniqueness-mixed-certificate-reduction}
给出。组合即得。证毕。
\end{proof}


\begin{remark}[射影不变量、局部因子与 \(2\)-进地址化：对“存在/唯一”的三类增量约束]\label{rem:window6-leyang-integer-audit-triple-rigidity}
在 window-\(6\) 的审计范式下，Lee--Yang 模块与整数证书链的增益并不体现在引入新的连续结构，而在于把若干原本依赖结构选择的环节，重写为三类可交叉核验的刚性约束之交汇：射影不变量约束、算术判别式/局部因子约束、以及 \(2\)-进动力学地址化约束。
该重写的作用是将“存在性/唯一性”从纯形式的代数闭包问题，推进到可输出证书的复核任务。

\paragraph{证明对象的可审计化起点}
交换子包络链条以定义 \ref{def:window6-coordinate-pushforward-lie-envelope} 的整数计数算子 \(L_i\) 为生成元，并以交换子闭包定义 \(\mathfrak g_6\)。
该对象链在结构上已由定理 \ref{thm:window6-lie-envelope-sl21} 完成识别：\(\mathfrak g_6=\mathfrak{sl}_{21}\)。
因此，秩 \(3\) 简单根型 \(B_3/C_3\) 的 \(21\) 维结构应被放置在抽象括号链条中（定理 \ref{thm:window6-root-dictionary-pullback-bracket} 与猜想 \ref{conj:window6-intrinsic-bracket-uniqueness}），或作为 \(\mathfrak{sl}_{21}\) 内受限可观测紧致部分的抽取目标（猜想 \ref{conj:window6-sm-gauge-as-minimal-compact-observable}）。

\paragraph{整数证书链已给出的两类端点刚性}
一方面，若“可观测紧致部分”的维数审计落在 \(D=12\) 且满足 Zeckendorf 原子性与逐项对齐口径，则其同构型被命题 \ref{prop:sm-zeckendorf-lie-algebra-rigidity} 强制为
\(
\mathfrak u(1)\oplus\mathfrak{su}(2)\oplus\mathfrak{su}(3)
\)；
该结论为“唯一闭合”的终点提供了完全离散的刚性约束。
另一方面，\(\delta=34\) 的分级约束同时由 uplift 覆盖、几何稳定子掩码恒等式与 parity 规则给出（推论 \ref{cor:foldbin6-geo-mask-34} 及推论 \ref{cor:terminal-delta-parity-rule}），从而把猜想 \ref{conj:window6-intrinsic-bracket-uniqueness} 中的分级相容从选择性输入提升为硬结构约束。

\paragraph{射影不变量的消混合接口}
在 Lee--Yang 零点的射影化消混合框架中，变量混合修正的主导项具有严格的 \(1\)-余边界结构（命题 \ref{prop:xi-leyang-variable-mixing-coboundary}），并且有限交比集合在射影等价（\(\mathrm{PGL}_2(\CC)\) 作用）意义下为完备不变量（命题 \ref{prop:xi-leyang-crossratio-projective-fingerprint-complete}）。
进一步，交比在两体积两零点的极限中对线性混合参数不敏感（定理 \ref{thm:xi-leyang-two-scale-crossratio-invariant}）。
因此，一旦某一 window-\(6\) 可观测耦合（或括号结构常数的等价类）能够被函子化地编码为某个“零点/分歧点配置”的射影轨道，则射影不变量为“同构意义下唯一”提供一种可复核的离散化口径：同构自由度被压缩为有限交比集合的匹配问题，并与余边界闭路一致性方程形成互补约束。

\paragraph{算术判别式核与局部因子约束}
在椭圆化的 Lee--Yang 模型中，谱四次 \(\Pi(\lambda,y)\) 的结构被固定椭圆曲线 \(E_{\min}\) 与其判别式数据所锚定（定理 \ref{thm:xi-terminal-zm-leyang-elliptic-structure}）：
其 \(y\)-判别式 \(4\lambda^3-4\lambda+1\) 同时实现为短 Weierstrass 模型的 \(2\)-除法多项式，并给出仅在素数 \(37\) 处坏约化的判别式/导子锚点。
在证书语义上，这类“判别式核—局部因子”数据为零点几何提供一条实现无关的算术指纹。
若将 window-\(6\) 的完成化多项式/谱对象与该椭圆锚点通过一个明确的函子化桥接对齐，则可将若干原本依赖于结构选择的分支（例如分级通道与可见的分歧配置）进一步压缩为“局部因子一致性”的可审计条件，从而增强排他性。

\paragraph{\(2\)-进动力学地址化约束}
除零点几何与局部因子外，文稿在 \(2\)-进通道给出一类可复核的地址化动力学指纹：碰撞矩序列在模 \(2^k\) 上进入终端纯周期相并满足严格的里程表律与线性阈值律（见段落 \ref{subsubsec:pom-collision-moment-2adic-odometer} 及其汇总模板）。
与递归地址化“可见 \(\sigma\)-代数不扩张”的框架（\S\ref{sec:recursive-addressing}）合并后，这类 \(2\)-进指纹可被用作 transport/renormalization 的外部刚性输入：任何候选括号若宣称与协议读出兼容，则其诱导的可观测谱对象与分支选择应在该地址化动力学下保持协变或不变；该约束为 Levi 分块相容与分级相容提供一条可并行核验的动力学证书通道。

\paragraph{证书化推进的可执行形态}
在上述三类增量约束并入后，猜想 \ref{conj:window6-intrinsic-bracket-uniqueness} 的“存在/唯一”可被组织为一个混合证书可行性问题：
\begin{itemize}
  \item \textbf{多项式约束：}反对称性与 Jacobi 恒等式，加上局域支撑模板、parity 分级封闭与 Levi 分块相容（命题 \ref{prop:window6-bracket-certificate-reduction}）。
  \item \textbf{半正定约束：}由 unitary slice 的 Toeplitz--PSD 证书链给出的有限维可行性（定理 \ref{thm:app-horizon-toeplitz-lmi} 与定理 \ref{thm:xi-toeplitz-det-verblunsky}；并参见注记 \ref{rem:window6-leyang-toeplitz-mixed-certificates}）。
  \item \textbf{射影/算术/地址化约束：}由交比指纹、判别式核与 \(2\)-进动力学地址化提供的实现无关交叉核验（本注记前三段）。
\end{itemize}
该组织方式的关键在于：它把“同构意义下唯一”从纯代数消元，扩展为“多项式 + PSD + 不变量匹配”的混合排他性判别，从而更贴近本文的“整数—证书”审计口径。

\paragraph{仍需补齐的关键构造}
要将上述增量约束上升为可引用的定理链条，缺口可集中为以下几类明确的构造问题：
\begin{itemize}
  \item \textbf{函子化桥接：}给出从 \((X_6,\Gamma_6,\delta=34)\) 及其可观测谱对象到某一 Lee--Yang 零点/分歧配置的显式映射，使射影指纹与判别式核可被严格解释为 window-\(6\) 侧的不变量。
  \item \textbf{不变量同源性：}在算术口径下证明 parity 通道与局部因子/判别式核之间的强相容（从而把分级相容的证书来源由“几何可实现性”提升为“局部因子必然性”的交叉约束）。
  \item \textbf{证书输出接口：}将候选括号与可观测紧致部分的核验数据整理为可复核的有限证书（结构常数支撑表、分级闭包核验、以及与 Toeplitz--PSD/射影指纹的一致性接口）。
\end{itemize}
\end{remark}
